\documentclass[11pt]{article}
\usepackage{amsmath,amssymb,amsthm,mathtools}
\usepackage{enumitem}
\usepackage{booktabs,array}
\usepackage[margin=2.4cm]{geometry}
\usepackage[hidelinks]{hyperref}
\setlist{itemsep=1pt,topsep=3pt}

% The handbook's macros, so the family's notation is available.
\newcommand{\cE}{\mathcal E}
\newcommand{\cB}{\mathcal B}
\newcommand{\cH}{\mathcal H}
\newcommand{\one}{\mathbf 1}
\newcommand{\mvec}{\mathbf m}
\newcommand{\bmu}{\boldsymbol\mu}
\newcommand{\ip}[2]{\langle #1,#2\rangle}
\newcommand{\oE}{\mathbb E}
\newcommand{\sG}{\mathsf G}
\newcommand{\sZ}{\mathsf Z}
\newcommand{\sB}{\mathsf B}
\newcommand{\sT}{\mathsf T}
\newcommand{\Gst}[4]{\Gamma_{\substack{#1\\#2}}\!\begin{bmatrix}#3\\#4\end{bmatrix}}
\newcommand{\IZ}{\mathbb I_{\mathbb Z}}
\newcommand{\Ilat}[1]{\mathbb I_{#1}}
\newcommand{\Tx}[3]{(\sT_{#1;(1)}[x]\,#2)((#3,\infty))}
\DeclareMathOperator{\Reg}{Reg}
\DeclareMathOperator{\re}{Re}
\DeclareMathOperator{\im}{Im}
\DeclareMathOperator*{\CT}{CT}
% Typing aids only; each prints the full expression it stands for.
\newcommand{\Pc}[1]{P_{\cB}(#1\,\one;#1,#1)}
\newcommand{\Clo}[1]{[\alpha^0]\big(K(\alpha)\Psi(\alpha)/\alpha^{\cE}\big)^{#1}}
\newcommand{\Cor}{\sum_\pi\prod_{k=2}^nc(U_k)^{-1}}
\newcommand{\Pth}{P_{\binom{[3]}1}(\one;1,1)}
\newcommand{\Pfour}{P_{\binom{[4]}1}(\one;1,1)}

\theoremstyle{plain}
\newtheorem{proposition}{Proposition}[section]
\newtheorem{theorem}[proposition]{Theorem}
\newtheorem{corollary}[proposition]{Corollary}
\theoremstyle{definition}
\newtheorem{definition}[proposition]{Definition}
\newtheorem{procedure}[proposition]{Procedure}
\theoremstyle{remark}
\newtheorem{remark}[proposition]{Remark}
\renewcommand{\thesection}{15}
\renewcommand{\theproposition}{15.\arabic{proposition}}
\newcommand{\derived}{\textup{[derived]}}
\newcommand{\computed}{\textup{[computed]}}
\newcommand{\recalled}{\textup{[recalled]}}
\newcommand{\notyet}{\textup{[not yet computed]}}

\title{Foreman notebook, \S\S15--19\\[4pt]\large The centre ray read off the list; the five-mark lists along their hinge trees; the masters of $\binom{[5]}2$ and two marks doubled, closed; the reading of a ray through the origin; the three-mark centre, set as an examination}
\date{22--23 October 2026}
\author{}

\begin{document}
\maketitle

\noindent\emph{This entry follows \S14.17 of the running notebook, which was not available in this session; it is written to be appended there. Handbook citations are by printed number and page of the handbook PDF of 20 October 2026; ``the note'' is the tree-coaction note of 22 October 2026. Every object is a member of the serif family or a construction of the handbook, written with its slots. The receipts are in the script centre\_ray.py.}

\section{The centre ray read off the list}

The neighbours of the five-mark centre are the values of the centre ray of the list (Def.~35.9, p.~309), the period $P_{\cB}(\sigma\one;\sigma,\sigma)$. The interest of the list is one Gramma (Thm.~8.3, p.~54), the period is its residue on the balance wall, and the residue is the Gramma of the letter columns (Thm.~2.9, p.~15). Its share side (Prop.~4.2, p.~25) is the working form of this entry: a fibre integral of Gammas, one per basis for each letter, pinned by equal incidence. The pin is what the list contributes. Closing the fibre at the negative strengths (Thm.~4.5, p.~26) lands on the pairs of $k$-tuples of bases with equal incidence; the residue at zero shares is the corner of Thm.~38.5(i); the native shares (Thm.~4.3, p.~25) are uniform, the share side of the table's gauge; and the banana is the list whose pin identifies the two letters mark by mark, which turns its share side into the Mellin read of a power of one Gramma.

\subsection*{15.1\quad The centre period on its share side}

\begin{proposition}[The share side of the centre]\label{p:share}\derived; receipts \computed\
Let $(\cE,\cB)$ be a list, $n=|\cE|\ge2$, $e_0\in\cE$, and $\re\sigma>0$. Then
\[
\Gamma(\sigma)^2\,\Pc\sigma=\int_{\substack{|\mathbf a|=\sigma,\ |\mathbf b|=\sigma\\ \sum_Ta_T\one_T=\sum_Tb_T\one_T}}\ \prod_{T\in\cB}\Gamma(a_T)\,\Gamma(b_T),
\]
the fibre integral of Def.~4.1, p.~25, with one share $a_T$ for the letter $K$ and one share $b_T$ for the letter $\Psi$ at every basis, pinned by the two letter totals and by the incidence at the marks $e\neq e_0$; it has dimension $2|\cB|-n-1$, and the uniform shares $a_T=b_T=\sigma/|\cB|$ are an admissible real point.
\end{proposition}

\begin{proof}
By Thm.~8.3, p.~54, $\Gamma(s)\Gamma(t)\,J_{\cB}(\mvec;s,t)=\sG_{(\bmu;s;t);t_{\cB}}(\one)$ with the unit columns $(\mathbf e_e;0;0)$ and the letter columns $(\chi_T/2;1;0)$, $(-\chi_T/2;0;1)$. The linear form $\ell(x;v;w)=\sum_ex_e-\frac{2r-n}2(v-w)$ is $1$ on the unit columns and $0$ on the letter columns, and at the index it is $\kappa_{\cE}$ (Thm.~8.2(iii), p.~54): the balance wall is the first wall of this Gramma along $\ell$, with the letter columns as facet terms. By Thm.~8.10, p.~56, the period is the residue of $J_{\cB}$ there, and by Thm.~2.9, p.~15, at $p=0$ that residue is the Gramma of the facet terms. Its share side (Prop.~4.2) has one share per letter column; the pins of the two letter scales are $|\mathbf a|=s$, $|\mathbf b|=t$, and those of the marks at $\bmu=0$ are $\sum_T(a_T-b_T)\chi_T/2\in\mathbb C\one$, which with $|\mathbf a|=|\mathbf b|$ is $\sum_T(a_T-b_T)\one_T=0$. At the centre $s=t=\sigma$. The uniform shares satisfy the pins and have positive real parts.
\end{proof}

\noindent Receipts. For $\binom{[2]}1$ the fibre has dimension one and the integral is four Gammas on a line (Cor.~4.7, p.~26),
\[
\frac1{2\pi i}\int\Gamma(a)^2\,\Gamma(\sigma-a)^2\,da=\Gst41{\sigma,\ \sigma,\ \sigma,\ \sigma}{2\sigma},\qquad\text{so}\qquad \Pc\sigma=\Gst21{\sigma,\ \sigma}{2\sigma}\quad\derived.
\] For $\binom{[3]}1$ at $\sigma=1$ the two-dimensional fibre integral is $2.343907238689463\ldots$, the value of the receipts table below, to fourteen digits \computed.

\begin{proposition}[The banana's pin]\label{p:bananapin}\derived; receipt \computed\
For $\binom{\cE}1$ the pin is $a_e=b_e$ at every mark, so each mark carries $\Gamma(a_e)^2$, which is the share kernel of the Gramma $\sG_{0;(1,-1)}(1,y)=\int_0^\infty x^{-1}e^{-x-y/x}dx$ (Prop.~4.2), and
\[
\Gamma(\sigma)^2\,\Pc\sigma=\int_{|\mathbf a|=\sigma}\prod_{e\in\cE}\Gamma(a_e)^2=\int_0^\infty y^{\sigma-1}\,\sG_{0;(1,-1)}(1,y)^n\,dy .
\]
At $n=5$, $\sigma=1$ the right side is $39.945588011996245699\ldots$ \computed.
\end{proposition}

\begin{proof}
With $T=\{e\}$ the pins are $a_e-b_e=0$ for $e\neq e_0$ and the two totals, so $a=b$ and the fibre is the share integral $|\mathbf a|=\sigma$ (Def.~4.1). The share side of $\sG_{0;(1,-1)}(1,y)$ is $\frac1{2\pi i}\int\Gamma(c)^2y^{-c}dc$, and the read $\int_0^\infty y^{\sigma-1}(\cdot)\,dy$ of a product of $n$ such fibre integrals pins the sum of their shares at $\sigma$.
\end{proof}

So the ``Bessel moments'' are the banana's share side: the pin of a single mark turns the two letters into one Gramma per mark (alias $\sG_{0;(1,-1)}(1,y)=2K_0(2\sqrt y)$). A general list keeps a share per basis and per letter and couples them only through the incidence; that coupling is the broader arena.

\begin{proposition}[The corner is the residue at zero shares]\label{p:corner}\derived\
With $U_k=\{\pi(k),\dots,\pi(n)\}$ the tails of an order $\pi$ and $c$ the width (Def.~7.10, p.~49),
\[
\lim_{\sigma\downarrow0}\sigma^{n-1}\Pc\sigma=\Cor ,
\]
and $\sigma^{n-1}\Pc\sigma$ has no term of order $\sigma$. This is Thm.~38.5(i), p.~330, at the centre question, whose walls are $\sigma c(S)$ (Prop.~8.21(ii), p.~59); the term of order $\sigma$ there is a sum of $\log\theta$ at greedy bases, zero for a list. On the share side, $\sigma^{n+1}\Gamma(\sigma)^2\Pc\sigma$ tends to the same number: it is the residue of the fibre integral of Proposition~\ref{p:share} at the origin, where all $2|\cB|$ Gammas have their first pole. For $\binom{\cE}1$ the corner is $n!$, the $5!$ of \S14.17; for $\binom{[5]}2$ the tails of every order have widths $1,2,2,1$ and it is $30$.
\end{proposition}

\begin{proposition}[Closing the centre at the negative strengths]\label{p:negcorner}\derived\
For $k\in\mathbb N$, at $\sigma=-k$ the share side closes (Thm.~4.5, p.~26) onto the lattice of pairs $(\boldsymbol\alpha,\boldsymbol\beta)\in\mathbb N^{\cB}\times\mathbb N^{\cB}$ with $|\boldsymbol\alpha|=|\boldsymbol\beta|=k$ and $\sum_T\alpha_T\one_T=\sum_T\beta_T\one_T$, where every Gamma is at a pole, $a_T=-\alpha_T$, $b_T=-\beta_T$, with the residues $\prod_T(-1)^{\alpha_T+\beta_T}/(\alpha_T!\,\beta_T!)$ (Prop.~1.3(iv), p.~11). The continued period has a pole of order at most $n-1$ at $\sigma=-k$, and
\begin{align*}
\lim_{\sigma\to-k}(\sigma+k)^{n-1}\,\Pc\sigma&=\Cor\cdot\sum_{(\boldsymbol\alpha,\boldsymbol\beta)}\frac{(k!)^2}{\prod_T\alpha_T!\,\beta_T!}\\
&=\Cor\cdot\Clo k ,
\end{align*}
the sum over that lattice: it counts the pairs of $k$-tuples of bases with equal incidence.
\end{proposition}

\begin{proof}
On the balance wall the period is $\sum_\pi\int_{[0,1]^{n-1}}\prod_{k'=2}^nq_{k'}^{\kappa_{U_{k'}}-1}h_\pi(q)\,dq$ in the sector coordinates $\alpha_{\pi(k')}=\rho\,q_2\cdots q_{k'}$ (Thm.~8.7(ii), p.~55, and Thm.~8.10). At the centre $\kappa_{U_{k'}}=\sigma c(U_{k'})$ and $h_\pi=(k_\pi\psi_\pi)^{-\sigma}$, and $k_\pi\psi_\pi=\prod_{k'}q_{k'}^{c(U_{k'})}K(\alpha)\Psi(\alpha)/\alpha^{\cE}$ is a polynomial in $q$ with nonnegative coefficients and constant term one (Lemma~8.5, p.~55, weights one). By Lemma~8.6, p.~55, each cube integral continues with at most a simple pole in each $\kappa_{U_{k'}}$ at $-\mathbb N$, and at $\sigma=-k$ the $n-1$ poles meet only through the term $q^{kc}$, $c=(c(U_{k'}))_{k'}$, of $(k_\pi\psi_\pi)^k(k_\pi\psi_\pi)^{-(\sigma+k)}$, with coefficient $\prod_{k'}c(U_{k'})^{-1}[q^{kc}](k_\pi\psi_\pi)^k$. The map $\alpha^v\mapsto q^{(v(U_{k'}))_{k'}}$ is injective on exponents of degree zero, so $[q^{kc}](k_\pi\psi_\pi)^k=\Clo k$. Expanding each letter's $k$-th power (the share kernel of Cor.~4.4, p.~26, closed at the exponent $-k$) gives the lattice sum. On the share side the factor $(k!)^{-2}$ is the residue of $\Gamma(\sigma)^2$ at $-k$, and the common factor of the lattice points is the residue at zero shares, the corner (Proposition~\ref{p:corner}).
\end{proof}

\begin{corollary}[One law for the period and its closing]\label{c:law}\derived\
If polynomials $c_0,\dots,c_d$ satisfy
\[
\sum_{i=0}^dc_i(\sigma)\,\Pc{(\sigma+i)}=0
\]
identically, then for every $k\in\mathbb N$
\[
\sum_{i=0}^{\min(d,k)}c_i(-k)\,\Clo{k-i}=0 .
\]
So the least law of the period in its strength is at least as long as the least recurrence of its closings, and the closings obey the reflected law from $k=0$, with nothing below.
\end{corollary}

\begin{proof}
Multiply by $(\sigma+k)^{n-1}$ and let $\sigma\to-k$: the terms with $i\le k$ tend to $c_i(-k)$ times the corner times the closing at $k-i$ by Proposition~\ref{p:negcorner}, and the others are regular there.
\end{proof}

The laws come from the family's own relations. By the letter presentation (Prop.~3.6, p.~20) the centre period is a Breta with one bunch for the shape and one singleton bunch per letter; its letter splits and scale exchanges (Thm.~3.5, p.~19) are the letter and mark relations of the interest (Thm.~8.3(i)--(iii)), and along the centre ray they are the relations of Thm.~35.3, p.~308, whose eliminations are the laws of Cor.~35.6. The converse of Corollary~\ref{c:law}, that the least recurrence of the closings, read at $k=-\sigma$, is a law of the period, is not derived here; for the banana through seven marks it is not needed, the laws being derived by adding marks (Theorem~\ref{t:cert}). For every list below the minimal solution it gives agrees with an independent quadrature to fourteen digits, and the closings obey it from $k=0$ with nothing below \computed.

\begin{procedure}[The centre ray from the list]\label{p:proc}
(i) Close the share side at $\sigma=-k$: $x_k=\Clo k$, $k\le k_{\max}$, the letter $x_k$ abbreviating them in this procedure. (ii) Find their least recurrence $\sum_{j=0}^dp_j(k)x_{k+j}=0$ by exact linear algebra and check it on the surplus equations. (iii) Read it at $k=-\sigma$: $c_i(\sigma)=p_{d-i}(-\sigma-d)$, $\sum_ic_i(\sigma)\Pc{(\sigma+i)}=0$. (iv) The period is its minimal solution: $\Pc{(\sigma+1)}/\Pc\sigma\to K(\one)^{-2}$ (Prop.~35.7, p.~309), the native value of least modulus. (v) Scale it by the residue at zero shares: $c_0$ has a zero of order $n-1$ at $0$, and $[\sigma^{n-1}]c_0\cdot\Cor+\sum_{i\ge1}c_i(0)\Pc i=0$.
\end{procedure}

\noindent Receipts \computed, no quadrature in the procedure:
\begin{center}
\begin{tabular}{lll}
\toprule
list & $P_{\cB}(\one;1,1)$ & reading\\ \midrule
$\binom{[3]}1$ & $2.343907238689458890601562\ldots$ & $3L(\chi_{-3},2)$ (alias)\\
$\binom{[4]}1$ & $8.414398322117159997798167\ldots$ & $\Tx{3}{\Ilat{\frac12+\mathbb Z}}{0}=7\,\Tx{3}{\IZ}{0}$\\
$\binom{[4]}2$ & $2.343907238689458890601562\ldots$ & as $\binom{[3]}1$: Cor.~19.17 (p.~143) and duality\\
$\binom{[5]}1$ & $39.94558801199624569894402\ldots$ & $M_5(1)$ of \S14.17\\
\bottomrule
\end{tabular}
\end{center}
The four-mark values are those of the note (Ex.~6.3 and 6.2, p.~39). There they are closings of the share side at positive strength: the Gammas of a basis meet in double poles, whose Laurent data are Trailmoments of the counting input and of lattice inputs (Prop.~5.14, p.~33; Lemma~35.20, p.~314).

\subsection*{15.2\quad The masters are the signs}

A native gauge of the centre ray is a point where its direction $(\one;1,1)$ is native (Def.~35.4, p.~308): the defect $1-R_e(\alpha)-R^\ast_e(\alpha)$ vanishes at every mark, and since $R^\ast_e(\alpha)=1-R_e(\alpha^{-1})$ (Thm.~7.14, p.~50) this reads
\[
R_e(\alpha)=R_e(\alpha^{-1})\qquad(e\in\cE):
\]
inverting the gauge leaves the occupation unchanged. Its native value is $\big(K(\alpha)K(\alpha^{-1})\big)^{-1}$.

\begin{proposition}[The banana has only signs]\label{p:onlysigns}\derived\
The native gauges of the centre ray of $\binom{\cE}1$ are exactly the signs $\eta\in\{\pm1\}^{\cE}$, $\eta_{e_0}=1$, with $K(\eta)=\sum_e\eta_e\neq0$; their native values are $K(\eta)^{-2}$ (Thm.~35.12, p.~310).
\end{proposition}

\begin{proof}
For $K=\sum_e\alpha_e$ the occupation is $R_e(\alpha)=\alpha_e/K(\alpha)$, so $R_e(\alpha)=R_e(\alpha^{-1})$ for every $e$ says $\alpha_e^2=K(\alpha)/K(\alpha^{-1})$ for every $e$.
\end{proof}

So along the banana sequence the native values are $(n-2j)^{-2}$, $j$ the number of minus signs, $\lceil n/2\rceil$ of them, all simple. At $n=5$ they are $\frac1{25},\frac19,1$, the symbol of the ray law vanishes there (Cor.~35.6, p.~308), the order is three, and the corner leaves two masters. Every list has the signs with $K(\eta)\neq0$ among the native gauges of its centre ray; what is special about $\binom{\cE}1$ is that it has nothing else.

\begin{remark}[The lowered settings of \S14.17]\label{r:lowered}
With $g=\sG_{0;(1,-1)}(1,\cdot)$ and $h(y)=\sum_my^m/\Gamma(1+m)^2$, the settings $x^{(j)}(s)=\int_0^\infty y^{s-1}g^{5-j}h^j\,dy$, $j=0,1,2$, satisfy the ray law in the normalization $M(s)=\Gamma(s)^2\Pc s$, and $x^{(0)}=M$ is Proposition~\ref{p:bananapin}. Their corners are $120,0,0$ \computed, and $x^{(j)}(s+1)/\big(s^2x^{(j)}(s)\big)\to(5-2j)^{-2}$: they are the sign classes with $j$ minus signs. ``The $g$'s outnumber the $h$'s'' is $K(\eta)\neq0$; the tie at four marks is $K(\eta)=0$, which is not a native gauge.
\end{remark}

\begin{proposition}[The lowered-setting block is the frame]\label{p:frame}\derived; receipt \computed\
Let $\Phi(s)=\det\big[x^{(i)}(s+l)\big]_{i,l=0}^2$. By Prop.~35.8, p.~309, applied to the ray law of $M$ (end coefficients $s^6$ and $-225$), $\Phi(s+1)=\frac{s^6}{225}\Phi(s)$ and $\Phi(s)=\Phi(1)\,\Gamma(s)^6\,225^{1-s}$. The amplitudes at the three sign classes, $x^{(j)}(s)\sim2\pi^{(5-2j)/2}2^{-j}\rho_j^{5/2}\,\Gamma(2s-\frac52)\,\rho_j^{-2s}$ with $\rho_j=2(5-2j)$, give
\[
\Phi(1)=\frac{128\sqrt{15}}{16875}\,\pi^3=0.910880564798203032890397187975\ldots,
\]
which the settings computed directly (quadrature, as a check) reproduce to thirty digits. The corners reduce the third column, so
\[
\Phi(1)=\frac{120}{225}\Big(x^{(1)}(1)\,x^{(2)}(2)-x^{(1)}(2)\,x^{(2)}(1)\Big),\qquad x^{(1)}(1)x^{(2)}(2)-x^{(1)}(2)x^{(2)}(1)=\frac{16\sqrt{15}}{1125}\,\Gamma(\tfrac12)^6 :
\]
the block of \S14.17 is the frame of the ray over its corner. The masters do not enter $\Phi$; their coefficients are minors that the zero corners of the lowered settings kill. So the frame does not shorten the last minor of \S14.17.
\end{proposition}

\subsection*{15.3\quad A parallel class closes through the banana}

\begin{proposition}\label{p:parallel}\derived\
Let $\cB$ be obtained from a list $\cH$ on the marks $1,\dots,m$ by replacing mark $i$ with a parallel class $P_i$ of $s_i$ marks, so that $K_{\cB}(\alpha)=K_{\cH}\big((\alpha(P_i))_i\big)$ with $\alpha(P_i)=\sum_{e\in P_i}\alpha_e$. Then $[\alpha^v]K_{\cB}^k=\big([A^w]K_{\cH}^k\big)\prod_i\big(w_i!/\prod_{e\in P_i}v_e!\big)$ with $w_i=v(P_i)$, and the closings of $\cB$ are
\[
\sum_w\big([A^w]K_{\cH}(A)^k\big)^2\prod_i[\alpha^0]\Big(K_{P_i}(\alpha)\Psi_{P_i}(\alpha)/\alpha^{P_i}\Big)^{w_i},
\]
with $K_{P_i}=\sum_{e\in P_i}\alpha_e$ and $\Psi_{P_i}$ the letters of the banana $\binom{P_i}1$: its closings at $-w_i$.
\end{proposition}

The banana enters the closing only through its own closings, as the parallel node enters the network through its full leans (Cor.~25.9, p.~209).

\subsection*{15.4\quad The five-mark centres}

All connected lists on five marks, up to duality, which leaves the centre unchanged (Prop.~8.14, p.~57); ``doubled'' and ``tripled'' are parallel classes. \computed
\begin{center}
\small
\begin{tabular}{@{}l c l l c c c@{}}
\toprule
list & $|\cB|$ & $K(\eta)^2$ at the signs (mult.) & roots of the symbol & order & corner & masters\\ \midrule
$\binom{[5]}1$ & 5 & $25\,(1),\,9\,(5),\,1\,(10)$ & $25,9,1$ & 3 & 120 & 2\\
$\binom{[4]}2$, a mark doubled & 9 & $81\,(1),\,9\,(6),\,1\,(9)$ & $81,9,1$ & 3 & 36 & 2\\
$\binom{[5]}2$ & 10 & $100\,(1),\,4\,(15)$ & $100,4,1,-8$ & 4 & 30 & 3\\
$\binom{[3]}2$, two marks doubled & 8 & $64\,(1),\,16\,(2),\,4\,(8),\,0\,(5)$ & $64,16,4,4,1$ & 5 & 44 & $3^\dagger$\\
$\binom{[3]}2$, a mark tripled & 7 & $49\,(1),\,25\,(1),\,9\,(3),\,1\,(11)$ & $49,25,9,1,1,1$ & 6 & 60 & $4^\dagger$\\
\bottomrule
\end{tabular}

\medskip
\begin{tabular}{@{}l l@{}}
\toprule
list & $P_{\cB}(\one;1,1)$\\ \midrule
$\binom{[5]}1$ & $39.94558801199624569894402849842399589184978292477\ldots$\\
$\binom{[4]}2$, a mark doubled & $5.4939011435808440124442576078206911958536386406801\ldots$\\
$\binom{[5]}2$ & $4.0297141707110496073206957344284069331398675636396\ldots$\\
$\binom{[3]}2$, two marks doubled & $7.8658373683640215924219746974461682109823828094581\ldots$\\
$\binom{[3]}2$, a mark tripled & $12.848722141915555957312339931253534636619123934709\ldots$\\
\bottomrule
\end{tabular}
\end{center}
The roots of the symbol are those of the recurrence of the closings, the values $K(\alpha)K(\alpha^{-1})$ at the native gauges; the ray law of the period has their reciprocals. ``Masters'' is $m-1$ for the least $m$ with a rational relation among $1,P_{\cB}(\one;1,1),\dots,\Pc m$ (PSLQ at $120$ digits, coefficients below $10^{15}$, values good to $150$ digits): over $\mathbb Q$, for the symmetric insertion, on the centre ray. In the first three rows the one relation is the corner.

$^\dagger$\,In the last two rows $c_0$ and $c_1$ both carry $(\sigma+1)^2$, so the ray law at $\sigma=-1$ reaches the period at neither $-1$ nor $0$, and it is a second relation, a \emph{second corner}: the boundary term, the constant term at $\sigma=-1$ (Thm.~1.6, p.~12),
\[
\CT_{\sigma=-1}\Big(c_0(\sigma)\Pc\sigma+c_1(\sigma)\Pc{(\sigma+1)}\Big)=-\sum_{i\ge2}c_i(-1)\,\Pc{(i-1)},
\]
is rational \computed. In lowest terms the two relations read
\begin{align*}
&2503-566\,P_{\cB}(\one;1,1)+88248\,\Pc2-908872\,\Pc3+1425600\,\Pc4=0,\\
&13800002-1848045\,P_{\cB}(\one;1,1)+227950815\,\Pc2-3601110627\,\Pc3\\
&\qquad\qquad+11177654265\,\Pc4-8565102000\,\Pc5=0,
\end{align*}
so the masters are $5-2=3$ and $6-2=4$. Write $c_0=(\sigma+1)^2A$ and $c_1=(\sigma+1)^2B$. Proposition~\ref{p:negcorner} at $\sigma=-1$ requires $A(-1)\,|\cB|+B(-1)=0$, the closing at $1$ being $|\cB|$; it holds, $B(-1)/A(-1)=-8$ and $-7$ \computed. There the recursion of Corollary~\ref{c:law} leaves the corner at $-1$ free, so this is a check of Proposition~\ref{p:negcorner} that the recursion does not supply.

$\binom{[5]}2$ has native gauges off the signs: on the pattern $\alpha\propto(a,a,b,b,1)$ the native equations have solutions with $K(\alpha)K(\alpha^{-1})=1$ and $-8$. There is no rational relation among $1$, the two masters of $\binom{[5]}1$ and the three of $\binom{[5]}2$ (coefficients below $10^8$, 150 digits), nor with $\pi^2$ adjoined.

\subsection*{15.5\quad The native shares, the roots, and an exact pairing}

Every root of the symbol of the closings' recurrence is the reciprocal native value $K(\alpha_c)K(\alpha_c^{-1})$ of a native gauge $\alpha_c$ of the centre ray (Cor.~35.6, p.~308).

\begin{proposition}[Both reads sit at the native shares]\label{p:edge}\derived; receipts \computed\
Let the list be connected, and let $H=2\big(\kappa_{\one}(X_e,X_{e'})\big)_{e,e'\neq e_0}$, twice the covariances of the indicators under the weight law (Def.~9.2, p.~61), the Hessian of Thm.~35.11, p.~310.
\begin{enumerate}[label=(\roman*)]
\item On the real fibre of Proposition~\ref{p:share} the native shares of Thm.~4.3(ii), p.~25, are the uniform shares $a_T=b_T=\sigma/|\cB|$: every basis takes the same share from each letter. This is the table's gauge $\one$ read on the share side.
\item The period reads the fibre along its imaginary directions. By Thm.~4.3(iii), which at the period is Thm.~38.5(ii), p.~330,
\[
\Pc k\sim K(\one)^{-2k}\Big(\frac{2\pi}k\Big)^{\frac{n-1}2}(\det H)^{-\frac12}.
\]
\item The closing reads the same fibre on its integer points. Dropping the pin, the multinomial theorem gives $\sum_{|\boldsymbol\alpha|=|\boldsymbol\beta|=k}(k!)^2/\prod_T\alpha_T!\beta_T!=K(\one)^k\Psi(\one)^k=|\cB|^{2k}$; the pin keeps the Gaussian of the fibre at the native shares scaled by $k$, and
\[
\Clo k\sim K(\one)^{2k}(2\pi k)^{-\frac{n-1}2}(\det H)^{-\frac12}.
\]
\item So $k^{n-1}\,\Pc k\,\Clo k\to(\det H)^{-1}$: one Gaussian, read once along a contour and once on a lattice, so that the powers of $2\pi$ cancel and the determinant appears twice.
\end{enumerate}
\end{proposition}

\begin{proof}
(i) The uniform point satisfies the pins, and $\nabla S=-\log\mathbf w$ is constant there, a combination of the two rows of the letter totals; $S$ is strictly concave on the real fibre (Thm.~4.3(ii)). (ii) Thm.~38.5(ii) at the centre: $X(0)=-2\log K(\one)$ and $-\nabla_z^2X(0)=8\kappa_{\one}=4H$ on the chart, $y=2z$; Thm.~35.11 is the same edge read on the generating function of the ray. (iii) The terms $(k!)^2/\prod\alpha_T!\beta_T!$ on the pinned lattice peak at the uniform shares scaled by $k$; their Gaussian there is the quadratic form of $\sum\log\Gamma$ at the native shares restricted to the fibre, the same form as in (ii), now summed over the integer points of the fibre; $\nabla^2_zE(0)=8\kappa_{\one}$ (Thm.~7.14, p.~50) identifies it with $H$.
\end{proof}

\noindent Receipts \computed: $(\det H)^{-1}=\frac{3125}{16},\frac{729}{16},\frac{3125}{81},\frac{512}9,\frac{16807}{192}$ in the order of the table of \S15.4. At $k=70$ the product agrees with these to about $10^{-4}$, the error falling like $k^{-2}$. For $\binom{[2]}1$ the product is exact at every $k$: $k\,\Gst21{k,\ k}{2k}\binom{2k}k=2=(\det H)^{-1}$ \derived.

\begin{remark}[The signs are the pin read modulo two]
Averaging over the signs detects the incidence modulo two: for every $k$,
\[
2^{-n}\sum_{\eta\in\{\pm1\}^{\cE}}K(\eta)^{2k}=\sum_{\substack{|\boldsymbol\alpha|=|\boldsymbol\beta|=k\\ \sum_T\alpha_T\one_T\equiv\sum_T\beta_T\one_T\ (2)}}\frac{(k!)^2}{\prod_T\alpha_T!\,\beta_T!}\ \ \ge\ \Clo k\quad\derived,
\]
since $K(\eta)\Psi(\eta)/\eta^{\cE}=K(\eta)^2$. So the sign values $K(\eta)^2$ are the growth rates of the pin's reduction modulo two, with $K(\one)^2$ the largest. The full pin keeps them as roots (Thm.~35.12) and adds the native gauges off the signs, below.
\end{remark}

\begin{proposition}[An exact pairing of the period with its closing for the five-mark banana]\label{p:pairing}\derived\ from the ray law; receipt \computed.\ For $\binom{[5]}1$ the ray law $\sum_{i=0}^3c_i(\sigma)\Pc{(\sigma+i)}=0$ has
\begin{gather*}
c_0=-\sigma^4,\qquad c_1=35\sigma^4+70\sigma^3+63\sigma^2+28\sigma+5,\\
c_2=-(\sigma+1)^2(259\sigma^2+518\sigma+285),\qquad c_3=225(\sigma+1)^2(\sigma+2)^2,
\end{gather*}
and it is self-adjoint under the reflection of the strength: $c_i(-\sigma-i)=c_i(\sigma)$ for every $i$. So the closings, read as a function of $\sigma$, solve the adjoint law, and for every $\sigma\in\{1,2,3,\dots\}$
\[
\sum_{i=1}^3\sum_{j=1}^ic_i(\sigma-j)\ \Clo{\sigma-j}\ \Pc{(\sigma-j+i)}=120,
\]
the residue at zero shares, with the closings at negative indices read as $0$.
\end{proposition}

\begin{proof}
The left side, $B(\sigma)$ in this proof, is the bilinear concomitant of the law and its adjoint: $B(\sigma+1)-B(\sigma)$ is the law applied to the period times the closing at $\sigma$, less the period at $\sigma$ times the adjoint law applied to the closings. The first vanishes by the law. The adjoint law at $\sigma$ is $\sum_ic_i(\sigma-i)x_{\sigma-i}$, which by $c_i(\sigma-i)=c_i(-\sigma)$ is the recurrence of Corollary~\ref{c:law} at $k=\sigma$, zero from $k=0$ on. At $\sigma=1$ only $j=1$ survives, and $B(1)=\sum_ic_i(0)\Pc i=120$ is the corner relation of Procedure~\ref{p:proc}(v). The self-adjointness is an identity of the four polynomials; $B(\sigma)=120$ holds to seventy digits for $\sigma\le55$ \computed.
\end{proof}

At large $\sigma$ the closings and the period move by the factors $25$ and $\frac1{25}$ per step, so $B(\sigma)\to\lambda\chi'(\lambda)\cdot\lim_kk^4\,\Pc k\,\Clo k$, with $\chi(\lambda)=\sum_i[\sigma^4]c_i\,\lambda^i=-1+35\lambda-259\lambda^2+225\lambda^3$ the symbol and $\lambda=\frac1{25}$. With Proposition~\ref{p:edge}:
\[
120\cdot\det H=\lambda\chi'(\lambda)\big|_{\lambda=1/25}=\frac{384}{625},\qquad \det H=\frac{16}{3125}\quad\derived.
\]
The residue at zero shares (the low end of strength, Thm.~38.5(i)) and the Gaussian at the native shares (the high end, Thm.~38.5(ii)) are the two readings of one exact identity between the fibre and its lattice. For the other four five-mark lists the closings solve no adjoint of the ray law in the forms tried: no combination $\sum_{l<d}r_l(\sigma)\Clo{\sigma+l}$ with polynomial $r_l$ of degree at most $20$, nor with the apparent factors of $c_0$ as denominators up to degree $14$, solves it \computed. In that framework the pairing of the period with its closing is not a single product there; it must carry other solution classes.

\begin{proposition}[Signs through parallel classes]\label{p:signsparallel}\derived\
In the setting of Proposition~\ref{p:parallel}, $K_{\cB}(\eta)=K_{\cH}\big(s_1(\eta),\dots,s_m(\eta)\big)$ with $s_i(\eta)=\sum_{e\in P_i}\eta_e\in\{s_i,s_i-2,\dots,-s_i\}$: the sign values of the bananas $\binom{P_i}1$ feed the letter of $\cH$.
\end{proposition}

For the triangle $\cH=\binom{[3]}2$, $K_{\cH}=e_2$, and with one mark tripled, $s=s_1(\eta)\in\{3,1,-1,-3\}$:
\[
e_2(s,1,1)=2s+1\in\{7,3,-1,-5\},\qquad e_2(s,1,-1)=-1\quad\text{for every }s,
\]
which is the census $49\,(1),\,25\,(1),\,9\,(3),\,1\,(11)$ of \S15.4. The second form does not see $s$.

\begin{proposition}[Letters that forget a class]\label{p:families}\derived\
Write $K_{\cB}=A_i\,K_{\cH/i}(A)+K_{\cH\setminus i}(A)$ with $A_i=\alpha(P_i)$ (Def.~7.3, p.~48), and $A_i'=\sum_{e\in P_i}\alpha_e^{-1}$.
\begin{enumerate}[label=(\roman*)]
\item Tripled triangle. On $\alpha_3=-\alpha_2$ the contraction letter $K_{\cH/1}=A_2+A_3$ vanishes, so every copy has occupation zero at $\alpha$ and at $\alpha^{-1}$, and $K(\alpha)=-\alpha_2^2$, $K(\alpha^{-1})=-\alpha_2^{-2}$ whatever the copies are. The native gauges there are
\[
\big\{\alpha_3=-\alpha_2,\ A_1=\alpha_2^2A_1'\big\},
\]
a family of dimension two modulo scale, all of native value $1$, containing the eight sign classes with $\eta_3=-\eta_2$.
\item Two marks doubled. On $P_2=(t,-t)$ the class letter vanishes, $A_2=A_2'=0$, and $K=A_1\alpha_3$ is the letter of $\cH\setminus2$, in which both remaining classes are coloops. The native gauges there are
\[
\Big\{P_1=(c,c),\ P_2=(t,-t),\ t^2=\frac{c\,\alpha_3\,(2\alpha_3+c)}{2c+\alpha_3}\Big\},
\]
a curve of native value $\frac14$ which on the torus $|\alpha_e|=1$ is a circle through the signs with $(s_1,s_2)=(\pm2,0)$.
\end{enumerate}
\end{proposition}

\begin{proof}
Nativity is $R_e(\alpha)=R_e(\alpha^{-1})$ at every mark. (i) At a copy $x\in P_1$, $R_x(\alpha)=\alpha_x(A_2+A_3)/K(\alpha)$ and $R_x(\alpha^{-1})=\alpha_x^{-1}(A_2'+A_3')/K(\alpha^{-1})$, both zero on $\alpha_3=-\alpha_2$; at $\alpha_2$ and $\alpha_3$ the two sides differ by $\mp(A_1/\alpha_2-\alpha_2A_1')$. (ii) At $\alpha_3$ both occupations are $1$ on $A_2=A_2'=0$; at $P_1$ the equation holds when $\alpha_{1a}^2=\alpha_{1b}^2=A_1/A_1'$, and $\alpha_{1b}=-\alpha_{1a}$ gives $K=0$; at $P_2$ it holds exactly at the stated $t^2$. When $|c|=|\alpha_3|=1$, $|2\alpha_3+c|=|2c+\alpha_3|$, so $|t|=1$.
\end{proof}

\begin{remark}[Local exponents]
Put $x_k\approx\lambda^kk^\beta$ into the recurrence of the closings \computed. At every simple root of every row, $\beta=-2=-\frac{n-1}2$, the exponent of a nondegenerate native gauge (Thm.~35.5(i), p.~308). At the repeated roots the indicial polynomials are
\[
(\beta+1)(\beta+2)(\beta+3)\quad(\text{tripled triangle},\ \lambda=1),\qquad(2\beta+3)(2\beta+5)\quad(\text{two marks doubled},\ \lambda=4).
\]
A family of real dimension $d$ on the torus $|\alpha_e|=1$, transversally nondegenerate, gives $\beta=-\frac{n-1-d}2$. The Hessian of $\log\big(K(\alpha)K(\alpha^{-1})\big)$ there has rank two along the family of (i) and rank three along the circle of (ii) \computed, so these are $-1$ and $-\frac32$. The isolated sign orbit at $\lambda=1$ (one copy plus, two minus, $\eta_2=\eta_3$) has full rank and gives $-2$. The exponents $-3$ and $-\frac52$: \notyet. These families are integral directions in the exceptional set of Thm.~35.19(i), p.~313: the order of the ray law counts native values with these multiplicities, and the families raise the order, and with it the masters, past the number of native values.
\end{remark}

\begin{remark}[Native gauges at complex points off the signs \computed]
$\binom{[5]}2$: $K(\alpha)K(\alpha^{-1})=1$ at $(\omega^2,\omega^2,\omega,\omega,1)$, $\omega=e^{2\pi\mathrm i/3}$, where $K=K(\alpha^{-1})=1$; and $-8$ at $(1,1,b,b,1)$, $b=-(2+\sqrt3)$, and at its inverse, where $K=-2-2\sqrt3$ and $K(\alpha^{-1})=-2+2\sqrt3$. Two marks doubled: $1$ at $(\omega,\omega^2,\omega,\omega^2,1)$. The signs are the points of order two of the torus; points of order three are the next layer that occurs here.
\end{remark}

\subsection*{15.6\quad Closing the banana at strength one}

\begin{proposition}[Four marks close on the half-lattice]\label{p:fourclose}\derived\
For $\binom{[4]}1$, merging the marks in pairs by four Gammas on a line (Cor.~4.7, p.~26), $\int_{a_1+a_2=c}\Gamma(a_1)^2\Gamma(a_2)^2=\Gst41{c,\ c,\ c,\ c}{2c}$, leaves a fibre of dimension one, and
\begin{align*}
P_{\cB}(\one;1,1)&=\frac1{2\pi i}\int_{(1/2)}\Gst41{c,\ c,\ c,\ c}{2c}\Gst41{1-c,\ 1-c,\ 1-c,\ 1-c}{2-2c}dc\\
&=\Tx{3}{\Ilat{\frac12+\mathbb Z}}{0}=7\,\Tx{3}{\IZ}{0}.
\end{align*}
\end{proposition}

\begin{proof}
By Lemma~5.16, p.~34, and the shift and reflection of the stacked Gamma (Prop.~1.3(iii),(v), p.~11) the integrand is
\[
\Gst40{c,\ c,\ 1-c,\ 1-c}{}^2\,\frac{\sin2\pi c}{\pi(1-2c)}=\frac{2\pi^3\cos\pi c}{(1-2c)\sin^3\pi c}.
\]
On the lines $\re c=N+\frac12$ it is $O\big(N^{-1}e^{-2\pi|\im c|}\big)$, so it closes to the right (Thm.~4.5, p.~26). At $c=m\ge1$, $2\pi^3\cos\pi c/\sin^3\pi c=2(c-m)^{-3}+O(c-m)$, with no parts of orders one and two, as in Lemma~5.16; the half-integers are zeros. So the residue is $2\,[\varepsilon^2](1-2m-2\varepsilon)^{-1}=-8/(2m-1)^3$, and the closing is $\sum_{m\ge1}8(2m-1)^{-3}=\sum_{x\in\frac12+\mathbb Z,\,x>0}x^{-3}$. The last equality is the coset relation at the step $2$ (Thm.~5.12(i), p.~32). The half-lattice is the doubled share $2c$ of a merged pair, the lower entry of its stacked Gamma.
\end{proof}

\noindent This is the note's value (Ex.~6.2, p.~39), now read as a closing of the share side.

\begin{proposition}[A mark added is a critical-line kernel]\label{p:addmark}\derived; receipts \computed\
For every $n\ge1$,
\[
P_{\binom{[n+1]}1}(\one;1,1)=-\sum_{m\ge1}\partial_\sigma P_{\binom{[n]}1}(\sigma\one;\sigma,\sigma)\big|_{\sigma=m}.
\]
\end{proposition}

\begin{proof}
In the share side of $\binom{[n+1]}1$ (Proposition~\ref{p:bananapin}) merge the first $n$ marks: their fibre at the total $d$ is $\Gamma(d)^2P_{\binom{[n]}1}(d\one;d,d)$, by the same proposition at $n$ marks. The last mark carries $\Gamma(1-d)^2$, and $\Gamma(d)^2\Gamma(1-d)^2=\Gst40{d,\ d,\ 1-d,\ 1-d}{}$, whose right closing reads the derivative of the rest at the positive integers (Lemma~5.16(i)). The rest is holomorphic on $\re d>0$, bounded on each vertical line by its value at the real part, and decays as $\re d\to\infty$, so the closing applies (Thm.~4.5).
\end{proof}

\noindent Receipts \computed. At $n=2$, with $\partial_\sigma\Gst21{\sigma,\ \sigma}{2\sigma}=2\Gst21{\sigma,\ \sigma}{2\sigma}\big((\log\Gamma)'(\sigma)-(\log\Gamma)'(2\sigma)\big)$ and $(\log\Gamma)'(j+1)=-\Tx{1}{\IZ}{j}$ (Prop.~5.14, p.~33),
\[
P_{\binom{[3]}1}(\one;1,1)=2\sum_{m\ge1}\Gst21{m,\ m}{2m}\,(\sT_{1;(1)}[x]\,\IZ)((m-1,2m-1])=2.343907238689458890601562\ldots,
\]
the period of two marks times a head of the counting input, to $25$ digits. At $n=3$ the sum is $7\Tx{3}{\IZ}{0}$ and at $n=4$ it is $39.945588011996245699\ldots$, both to $20$ digits.

\begin{remark}[The five-mark centre is a Trailmoment of a live input]\label{r:live}
By Proposition~\ref{p:addmark},
\[
P_{\binom{[5]}1}(\one;1,1)=\big(\sT_{0;(1)}[x]\,f\,\IZ\big)((0,\infty)),\qquad f(m)=-\partial_\sigma P_{\binom{[4]}1}(\sigma\one;\sigma,\sigma)\big|_{\sigma=m},
\]
a live input (Chapter~6): its density is read from the four-mark centre period, and the four-mark ray law differentiated in $\sigma$ gives it a recurrence of order two whose inhomogeneity is the values. At $\sigma=0$ the differentiated law is a relation among the first two slopes, the corner and the values, the term of order $\sigma$ of Proposition~\ref{p:corner} being zero. So every $f(m)$ lies in $\mathbb Q+\mathbb Q\Tx{3}{\IZ}{0}+\mathbb Q f(1)$ \derived. For example $P_{\binom{[4]}1}(2\,\one;2,2)=(7\zeta(3)-6)/16$ and $f(2)=-\big(36+21\zeta(3)-2f(1)\big)/32$ (alias $\zeta(3)=\Tx{3}{\IZ}{0}$) \computed. The first slope is
\[
-f(1)=48\,\mathrm{Li}_4(\tfrac12)+2\log^42-2\pi^2\log^22-\tfrac{17}{30}\pi^4=-39.38158212312216562493980934508\ldots
\]
by PSLQ at $60$ digits \computed. In the family these aliases are Trailmoments of index at most four of the counting input at the frequencies $0$ and $\frac12$ and of the lattice input of density $2^{-x}$ (Def.~5.2, p.~28). So the four-mark slope sits on the half-lattice of Proposition~\ref{p:fourclose}, at index four. The masters of the five-mark centre are values of this live Trailmoment, and this is where the closing stops landing on plain lattice inputs.
\end{remark}

\subsection*{15.7\quad Closing the live input of five marks}

Throughout this subsection $p(\sigma)=P_{\binom{[4]}1}(\sigma\one;\sigma,\sigma)$ for real $\sigma>0$ and $f=-p'$; $p\,\IZ$ and $f\,\IZ$ are the \emph{values input} and the \emph{slope input} (Def.~5.4, p.~29). The four-mark ray law is
\[
\sum_{i=0}^2c_i(\sigma)\,p(\sigma+i)=0,\qquad c_0=\sigma^3,\quad c_1=-2(2\sigma+1)(5\sigma^2+5\sigma+2),\quad c_2=64(\sigma+1)^3,
\]
read off the closings by Procedure~\ref{p:proc} and derived in Theorem~\ref{t:cert} below. For a function $R$ on a coset of $\mathbb Z$ put $(L^\ast R)(x)=\sum_ic_i(x-i)\,R(x-i)$, and let $L'^\ast$ be the same with the derivatives $c_i'$ in place of the $c_i$.

\begin{proposition}[The live recurrence]\label{p:liverec}\derived\
The slope input has a recurrence of order two whose inhomogeneity is the values input:
\[
\sum_{i=0}^2c_i(\sigma)\,f(\sigma+i)=\sum_{i=0}^2c_i'(\sigma)\,p(\sigma+i),\qquad c_0'=3\sigma^2,\quad c_1'=-6(10\sigma^2+10\sigma+3),\quad c_2'=192(\sigma+1)^2 .
\]
\end{proposition}

\begin{proof}
The law holds identically in $\sigma$ (Theorem~\ref{t:cert}); differentiate it.
\end{proof}

\begin{proposition}[A mark added, at every strength]\label{p:cosetmark}\derived; receipts \computed\
For real $s>0$ let $W_s(d)=\Gst22{d,\ d}{1+d-s,\ 1+d-s}$. Then
\begin{align*}
\Gamma(s)^2\,P_{\binom{[5]}1}(s\one;s,s)&=-\big(\sT_{0;(1)}[x]\,(W_sp)'\,\Ilat{s+\mathbb Z}\big)((s-1,\infty)),\\
x^{(1)}(s)&=\phantom{-}\big(\sT_{0;(1)}[x]\,W_sp\,\Ilat{s+\mathbb Z}\big)((s-1,\infty)),
\end{align*}
with $x^{(1)}$ the lowered setting with one $h$ of Remark~\ref{r:lowered}. The period reads the slopes of the word $W_sp$ on the coset of its strength, and the lowered setting reads its values. At $s\in\{1,2,\dots\}$ the word $W_s(x)=\prod_{j=1}^{s-1}(x-j)^2$ is a polynomial with double zeros at $1,\dots,s-1$, the cut moves to $0$, and
\[
\Gamma(s)^2P_{\binom{[5]}1}(s\one;s,s)=\Tx{0}{(W_sf-W_s'p)\,\IZ}{0},\qquad x^{(1)}(s)=\Tx{0}{W_sp\,\IZ}{0}.
\]
\end{proposition}

\begin{proof}
Merge the first four marks in the share side of five (Proposition~\ref{p:bananapin}): $\Gamma(s)^2P_{\binom{[5]}1}(s\one;s,s)=\frac1{2\pi i}\int_{(c)}\Gamma(d)^2p(d)\,\Gamma(s-d)^2\,dd$ with $0<c<s$, as in the proof of Proposition~\ref{p:addmark}. By the reflection (Prop.~1.3(v), p.~11), $\Gamma(s-d)\Gamma(1-s+d)=\pi/\sin\pi(s-d)$, so
\[
\Gamma(d)^2\,\Gamma(s-d)^2=W_s(d)\,\Gst40{s-d,\ s-d,\ 1-s+d,\ 1-s+d}{},
\]
the stacked Gamma of Lemma~5.16, p.~34, at $a=s-d$: double poles at $d\in s+\mathbb Z$, with no simple parts. The double zeros of $W_s$ at $d=s-1,s-2,\dots$ remove the poles between the contour and $s$. The rest is as in Proposition~\ref{p:addmark}: $W_sp$ is holomorphic on $\re d>0$, of polynomial growth along vertical lines, where the kernel decays like $e^{-2\pi|\im d|}$, and it decays like $16^{-\re d}$ to the right; closing to the right (Thm.~4.5, p.~26) reads $-(W_sp)'$ at $d=s,s+1,\dots$. For $x^{(1)}$, expand $h=\sum_ky^k/\Gamma(1+k)^2$ against $g^4$ (Proposition~\ref{p:bananapin} at four marks): $x^{(1)}(s)=\sum_{k\ge0}\Gamma(s+k)^2p(s+k)/\Gamma(1+k)^2$, and $\Gamma(s+k)^2/\Gamma(1+k)^2=W_s(s+k)$.
\end{proof}

\noindent Receipts \computed: at $s=\frac32$ both sides of the first identity are $2.27912905076377$, and at $s=\frac94$ both are $0.129204079141924$ (the left side by one-dimensional quadrature of the share side; fifteen digits).

\begin{proposition}[Summation by parts at the cut: what lands on lattice inputs]\label{p:sbp}\derived\
For every polynomial $R$,
\begin{gather*}
\Tx{0}{(L^\ast R)\,p\,\IZ}{0}=-24\,R(0),\\
\Tx{0}{(L^\ast R)\,f\,\IZ}{0}=\Tx{0}{(L'^\ast R)\,p\,\IZ}{0}.
\end{gather*}
\end{proposition}

\begin{proof}
By the shifted inputs (Thm.~6.2, p.~42) a shift of the input is a shift of the word, and the cut relation (Prop.~5.3, p.~29) returns the points that cross the cut: for a density $\varphi$,
\[
\sum_{x\ge1}R(x)\sum_ic_i(x)\,\varphi(x+i)=\Tx{0}{(L^\ast R)\,\varphi\,\IZ}{0}-R(0)\sum_{i\ge1}c_i(0)\,\varphi(i)-R(-1)\,c_2(-1)\,\varphi(1),
\]
and $c_2(-1)=c_2'(-1)=0$. For $\varphi=p$ the left side vanishes by the law. For $\varphi=f$ it is $\sum_xR(x)\sum_ic_i'(x)\,p(x+i)$ by Proposition~\ref{p:liverec}, to which the same move applies. So the boundaries are
\[
\sum_{i\ge1}c_i(0)\,p(i)=-4p(1)+64p(2),\qquad \sum_{i\ge1}\big(c_i(0)f(i)-c_i'(0)p(i)\big)=-4f(1)+64f(2)+18p(1)-192p(2).
\]
The first is $-[\sigma^3]c_0$ times the corner $24$ (Procedure~\ref{p:proc}(v)); the second is the term of order $\sigma$ of $\sigma^3p(\sigma)$, which vanishes (Proposition~\ref{p:corner}). With the values of Remark~\ref{r:live} they are $-24$ and $0$: the lattice Trailmoments in them, $p(1)=7\Tx{3}{\IZ}{0}$, $f(1)$ and $f(2)$, cancel identically.
\end{proof}

So what the closing of the live Trailmoment puts on lattice inputs is the residue at zero shares, and nothing else: neither $\Tx{3}{\IZ}{0}$ nor the four-mark slope $f(1)$ survives.

\begin{proposition}[The fixed block]\label{p:block}\derived; receipts \computed\
$L^\ast$ raises the degree of a polynomial by three, with leading factor $1-20+64=45$, the symbol of the law at $\lambda=1$, which is not a native value of four marks; so the polynomials are $L^\ast(\mathbb Q[x])\oplus\mathrm{span}\{1,\,x-1,\,(x-1)^2\}$. By Proposition~\ref{p:sbp}, every Trailmoment at index zero of a polynomial word against the values input or the slope input, cut at $0$, is a rational number plus a rational combination of the six
\[
\Tx{0}{(x-1)^j\,p\,\IZ}{0},\qquad \Tx{0}{(x-1)^j\,f\,\IZ}{0}\qquad(j=0,1,2).
\]
The share side adds one relation:
\[
5\,\Tx{0}{p\,\IZ}{0}-5\,\Tx{0}{(x-1)\,f\,\IZ}{0}=\Tx{0}{f\,\IZ}{0}.
\]
\end{proposition}

\begin{proof}
The splitting: $\deg L^\ast R=\deg R+3$, so $L^\ast$ is injective on polynomials and its image misses exactly three dimensions in each degree. The relation says that each mark takes on average a fifth of the strength. On the fibre of Proposition~\ref{p:bananapin} the shares sum to $\sigma$, and the integrand and the pins are symmetric in the five marks; so the fibre integral with one share $a_e$ inserted is the same for every $e$, and the five sum to $\sigma\,\Gamma(\sigma)^2P_{\binom{[5]}1}(\sigma\one;\sigma,\sigma)$. Merging the first four marks, the fifth share is $\sigma-d$, and at $\sigma=1$ the closing of Proposition~\ref{p:cosetmark} reads its insertion as $-\sum_{x\ge1}\partial_d\big((1-d)p(d)\big)\big|_{d=x}=\Tx{0}{p\,\IZ}{0}-\Tx{0}{(x-1)f\,\IZ}{0}$; the right side is $\frac15P_{\binom{[5]}1}(\one;1,1)=\frac15\Tx{0}{f\,\IZ}{0}$.
\end{proof}

\noindent Receipts \computed. At the strengths $2,\dots,5$ the same insertion reduces to rational multiples of this relation. There is no rational relation among $1$, $\Tx{3}{\IZ}{0}$, $f(1)$ and the five Trailmoments other than $\Tx{0}{(x-1)f\,\IZ}{0}$ with coefficients below $10^{30}$ (PSLQ, $360$ digits): in that framework, over $\mathbb Q$, the block is five-dimensional. Its six values:
\begin{center}
\small
\begin{tabular}{@{}lll@{}}
\toprule
$j$ & $\Tx{0}{(x-1)^j\,p\,\IZ}{0}$ & $\Tx{0}{(x-1)^j\,f\,\IZ}{0}$\\ \midrule
$0$ & $8.570280443374461268149695824246537286\ldots$ & $39.94558801199624569894402849842399589\ldots$\\
$1$ & $0.161069702276674858851310666248514198\ldots$ & $0.581162840975212128360890124561738107\ldots$\\
$2$ & $0.171874581383536904847683491575293916\ldots$ & $0.616835511971624849092924541153028126\ldots$\\
\bottomrule
\end{tabular}
\end{center}

\begin{theorem}[A mark added carries the ray law]\label{t:cert}\derived; receipts \computed\
Normalize the ray law of $\binom{[n]}1$, $\sum_{i=0}^{d_n}c^{(n)}_i(\sigma)\,P_{\binom{[n]}1}((\sigma+i)\one;\sigma+i,\sigma+i)=0$, so that $c_0^{(n)}=(-1)^n\sigma^{n-1}$, as for four marks above and for five marks in Proposition~\ref{p:pairing}, and let $L_n^\ast$ be its adjoint. For $n=1,\dots,6$,
\[
\sum_{i=0}^{d_{n+1}}c^{(n+1)}_i(s)\binom{x-1}{s-1+i}^{\!2}=\big(L_n^\ast A_s\big)(x),\qquad A_s(x)=\big((n+1)x-(n+2)s\big)\binom xs^{\!2},
\]
an identity of rational functions of $x$ and $s$ once both sides are divided by $\binom xs^2$. Read one way it produces the law: the words $\binom{x-1}{s-1+i}^2$ have distinct degrees, so the coefficients of $L_n^\ast A_s$ in them are the $c^{(n+1)}_i(s)$. Read the other way it proves it: for $2\le n\le6$, if the ray law of $\binom{[n]}1$ holds identically in $\sigma$, so does that of $\binom{[n+1]}1$, and the lowered setting $x^{(1)}$ of $n+1$ marks, divided by $\Gamma(s)^2$, obeys it too. The two-mark period $P_{\binom{[2]}1}(\sigma\one;\sigma,\sigma)=\Gst21{\sigma,\ \sigma}{2\sigma}$ obeys $\sigma P(\sigma)-2(2\sigma+1)P(\sigma+1)=0$ by the shift of the stacked Gamma (Prop.~1.3(iii), p.~11). So the ray laws of $\binom{[n]}1$, $n\le7$, are derived from the stacked Gamma, one mark at a time.
\end{theorem}

\begin{proof}
Write $p_n$ for the centre period of $\binom{[n]}1$, and let $s>0$ be non-integral. Proposition~\ref{p:cosetmark}, with $n$ marks merged in place of four, gives $P_{\binom{[n+1]}1}((s+i)\one;s+i,s+i)=-\sum_{d\in s+\mathbb N}(w_{s,i}\,p_n)'(d)$ with $w_{s,i}=W_{s+i}/\Gamma(s+i)^2$, the terms $d=s,\dots,s+i-1$ being zero by the double zeros of $W_{s+i}$. On the coset $w_{s,i}(d)$ is $\binom{d-1}{s-1+i}^2$ written with Gammas, so the displayed identity holds there. On $s+\mathbb N$ the move of Proposition~\ref{p:sbp} has no boundary for the words $A_s$ and $A_s'$, which vanish at $d=s-1,s-2,\dots$: $\sum_d(L_n^\ast A)\,\varphi=\sum_dA\,(L_n\varphi)$, and likewise with $L_n'$. With $(L_n^\ast A)'=L_n'^\ast A+L_n^\ast(A')$ and $L_np_n'=-L_n'p_n$ (the law differentiated),
\begin{multline*}
\sum_ic^{(n+1)}_i(s)\,P_{\binom{[n+1]}1}((s+i)\one;s+i,s+i)=-\sum_{d\in s+\mathbb N}\Big((L_n'^\ast A_s)\,p_n+A_s'\,L_np_n+A_s\,L_np_n'\Big)(d)\\
=-\sum_{d}(L_n'^\ast A_s)\,p_n+\sum_dA_s\,L_n'p_n=0,
\end{multline*}
and $\sum_ic^{(n+1)}_i(s)\,\Gamma(s+i)^{-2}x^{(1)}(s+i)=\sum_d(L_n^\ast A_s)\,p_n=\sum_dA_s\,L_np_n=0$. Both sides are holomorphic in $s$, so the laws hold identically. The displayed identity is a finite computation with rational functions of $(x,s)$.
\end{proof}

\noindent Receipts \computed. The certificate $A_s$ was found by exact linear algebra at each $n\le6$ and $s\le5$, the same at every $n$ up to the factor $(n+1)x-(n+2)s$, and the identity was then checked with $s$ indeterminate. At $n=1$ ($p_1=1$, $c^{(1)}=(-1,1)$) it holds as an identity, but the sums of the proof diverge. The five-mark law at $s=\frac12$ and $s=\frac3{10}$, from the share side by quadrature: residuals $4.6\cdot10^{-14}$ and $2.3\cdot10^{-14}$ against terms of size $10^2$. At $s=0$, with the cut at $0$, the same reduction returns the corner relation of $n+1$ marks as $A_0'(0)=n+1$ times the boundary of Proposition~\ref{p:sbp}: the corners $(n+1)!=(n+1)\cdot n!$ of Proposition~\ref{p:corner}, here $120=5\cdot24$.

\begin{corollary}[Two masters over the corner, two lowered numbers]\label{c:masters}\derived\
For every $s\in\{1,2,\dots\}$ there are rationals $a_s,b_s,q_s$ with
\begin{align*}
\Gamma(s)^2P_{\binom{[5]}1}(s\one;s,s)&=q_s+a_s\,P_{\binom{[5]}1}(\one;1,1)+b_s\,P_{\binom{[5]}1}(2\one;2,2),\\
x^{(1)}(s)&=a_s\,x^{(1)}(1)+b_s\,x^{(1)}(2),
\end{align*}
and in the block
\begin{gather*}
P_{\binom{[5]}1}(\one;1,1)=\Tx{0}{f\,\IZ}{0},\\
P_{\binom{[5]}1}(2\one;2,2)=\Tx{0}{(x-1)^2f\,\IZ}{0}-2\,\Tx{0}{(x-1)\,p\,\IZ}{0},\\
x^{(1)}(1)=\Tx{0}{p\,\IZ}{0},\qquad x^{(1)}(2)=\Tx{0}{(x-1)^2p\,\IZ}{0}.
\end{gather*}
Modulo $L^\ast(x\,\mathbb Q[x])$ every word $W_s$ lies in $\mathrm{span}\{1,(x-1)^2\}$: the ray reaches the direction $x-1$ of the block only through the value term $-2\Tx{0}{(x-1)p\,\IZ}{0}$ of its second master. For example $(a_3,b_3,q_3)=\big(-\frac1{45},\frac{19}{15},\frac8{15}\big)$ and $(a_4,b_4,q_4)=\big(-\frac{113}{1125},\frac{1907}{375},\frac{944}{375}\big)$.
\end{corollary}

\begin{proof}
The four block identities are Proposition~\ref{p:cosetmark} at $s=1,2$, with $W_1=1$, $W_2=(x-1)^2$, $W_2'=2(x-1)$. Write $W_s/\Gamma(s)^2=L^\ast B_s+r_s$ with $\deg r_s\le2$. Since the display of Theorem~\ref{t:cert} is an identity of polynomials and the splitting of Proposition~\ref{p:block} is direct, the $r_s$ obey the five-mark recurrence for $s\ge1$, and at $s=0$, where $A_0=5x$, they obey $\sum_{i\ge1}c^{(5)}_i(0)\,r_i=0$. As $r_1=1$ and $r_2=(x-1)^2$, every $r_s$ lies in their span, and the same recursion, from $B_1=B_2=0$, gives $B_s(0)=0$. Then $x^{(1)}(s)=\Gamma(s)^2\Tx{0}{r_s\,p\,\IZ}{0}$, and by Proposition~\ref{p:sbp} $P_{\binom{[5]}1}(s\one;s,s)=\Tx{0}{(r_sf-r_s'p)\,\IZ}{0}+24B_s'(0)$, where $\sum_{i\ge1}c^{(5)}_i(0)B_i'(0)=A_0'(0)=5$ carries the corner.
\end{proof}

\begin{proposition}[The lowered half is a stacked Gamma]\label{p:stackedhalf}\computed\
\[
x^{(1)}(1)=\Tx{0}{p\,\IZ}{0}=\frac1{120}\Gst84{\tfrac1{15},\ \tfrac2{15},\ \tfrac4{15},\ \tfrac8{15},\ \tfrac15,\ \tfrac25,\ \tfrac35,\ \tfrac45}{\tfrac12,\ \tfrac12,\ \tfrac12,\ \tfrac12}
\]
to $364$ digits, and
\begin{gather*}
x^{(1)}(2)=\frac{65\,x^{(1)}(1)-32\,\Gamma(\frac12)^8/x^{(1)}(1)}{1125}\quad(360\text{ digits}),\\
x^{(2)}(1)=\frac{\sqrt{15}}{4\,\Gamma(\frac12)^2}\,x^{(1)}(1),\qquad x^{(2)}(2)=\frac{\sqrt{15}}{4\,\Gamma(\frac12)^2}\cdot\frac{65\,x^{(1)}(1)+32\,\Gamma(\frac12)^8/x^{(1)}(1)}{1125}\quad(100\text{ digits}).
\end{gather*}
So the block of Proposition~\ref{p:frame} is $x^{(1)}(1)x^{(2)}(2)-x^{(1)}(2)x^{(2)}(1)=\frac{\sqrt{15}}{4\Gamma(1/2)^2}\cdot\frac{64\,\Gamma(1/2)^8}{1125}=\frac{16\sqrt{15}}{1125}\Gamma(\frac12)^6$: the stacked Gamma cancels from the frame, which is why the frame saw only $\Gamma(\frac12)$.
\end{proposition}

The values half that the ray reaches is two stacked-Gamma numbers, $x^{(1)}(1)$ and $\Gamma(\frac12)^8/x^{(1)}(1)$. The direction $x-1$ is not among them: $\Tx{0}{(x-1)\,p\,\IZ}{0}$ has no rational relation with $1$, $\Tx{3}{\IZ}{0}$, $\Gamma(\frac12)^4$, $x^{(1)}(1)$ and $\Gamma(\frac12)^8/x^{(1)}(1)$, nor with ten products of $x^{(1)}(1)^{\pm1}$, powers of $\Gamma(\frac12)^2$ and $\sqrt{15}$, with coefficients below $10^{12}$ \computed.

\begin{proposition}[The block on the four-mark chart and its pencil]\label{p:chart}\derived; receipts \computed\
On the chart $y_{e_0}=0$ of four marks let $X=K(\alpha)K(\alpha^{-1})$, the reciprocal of the stride of the centre ray (Def.~35.9, p.~309), so that $p(\sigma)=\int X^{-\sigma}dy$, the value along the ray from the real origin $(\sigma\one;\sigma,\sigma)$ (Def.~35.1, p.~307), with $X\ge16$ on the chart (Thm.~35.10(i), p.~309). For a polynomial $R$,
\begin{gather*}
\Tx{0}{R\,p\,\IZ}{0}=\int\rho_R(X)\,dy,\qquad\Tx{0}{R\,f\,\IZ}{0}=\int\log X\,\rho_R(X)\,dy,\\
\rho_R(X)=\sum_{m\ge1}R(m)X^{-m},
\end{gather*}
a rational function of $X$ whose only pole is $X=1$. In particular
\begin{gather*}
P_{\binom{[5]}1}(\one;1,1)=\int\frac{\log X}{X-1}\,dy,\qquad P_{\binom{[5]}1}(2\one;2,2)=\int\frac{(X+1)\log X-2(X-1)}{(X-1)^3}\,dy,\\
x^{(1)}(1)=\int\frac{dy}{X-1},\qquad x^{(1)}(2)=\int\frac{X+1}{(X-1)^3}\,dy,\qquad \Tx{0}{(x-1)\,p\,\IZ}{0}=\int\frac{dy}{(X-1)^2}.
\end{gather*}
The first is the left closing of Lemma~5.16(ii), p.~34, read on the chart. With the generating function of the four-mark ray from the origin $(\varepsilon\one;\varepsilon,\varepsilon)$, $G_\varepsilon(q)=\sum_{k\ge1}p(k+\varepsilon)\,q^{k-1}=\int X^{-\varepsilon}\,dy/(X-q)$ for $|q|<16$, carried by the pencil $K\Psi-q\alpha^{\one}$ (Thm.~35.10(ii)),
\begin{gather*}
\Tx{0}{(x-1)^j\,p\,\IZ}{0}=(q\partial_q)^jG_0(q)\big|_{q=1},\\
\Tx{0}{(x-1)^j\,f\,\IZ}{0}=-\partial_\varepsilon(q\partial_q)^jG_\varepsilon(q)\big|_{q=1,\ \varepsilon=0}:
\end{gather*}
the block is the jet at $q=1$ of the four-mark generating function and of its slope in the origin.
\end{proposition}

\begin{proof}
Def.~35.1 along the ray from $(\sigma\one;\sigma,\sigma)$, the positivity of every term, and $\sum_{m\ge1}R(m)X^{-m}$ summed for $X>1$: $\rho_1=1/(X-1)$, $\rho_{x-1}=1/(X-1)^2$, $\rho_{(x-1)^2}=(X+1)/(X-1)^3$. For the masters use Corollary~\ref{c:masters}; for the pencil, Thm.~35.10(ii) with $k_0=1$ and the origin form $X^{-\varepsilon}dy$.
\end{proof}

\noindent Receipts \computed, by three-dimensional quadrature on the chart, to twelve digits:
\[
\int\frac{dy}X=8.414398322115,\qquad\int\frac{\log X}{X-1}\,dy=39.945588011991,\qquad\int\frac{dy}{X-1}=8.570280443372,
\]
against $7\Tx{3}{\IZ}{0}$, the five-mark centre at strength one, and $x^{(1)}(1)$.

\begin{remark}[What the closing found]\label{r:found}
(i) The live Trailmoment lands on lattice inputs only through the residue at zero shares (Proposition~\ref{p:sbp}).

(ii) The words of the five-mark centre ray reach a plane of the five-dimensional block on each side: on the values side the lowered setting, two stacked-Gamma numbers (Proposition~\ref{p:stackedhalf}); on the slopes side the period, whose two masters are
\begin{gather*}
P_{\binom{[5]}1}(\one;1,1)=\Tx{0}{f\,\IZ}{0},\\
P_{\binom{[5]}1}(2\one;2,2)=\Tx{0}{(x-1)^2f\,\IZ}{0}-2\,\Tx{0}{(x-1)\,p\,\IZ}{0}:
\end{gather*}
the slope in the origin of the four-mark generating function at $q=1$, and its second jet less twice the first jet of the generating function itself. They have no rational relation with $1$, $\Tx{3}{\IZ}{0}$, $f(1)$, $\Gamma(\frac12)^4$, $\Gamma(\frac12)^8$, $\Tx{5}{\IZ}{0}$, $\Gamma(\frac12)^4\Tx{3}{\IZ}{0}$, the two stacked-Gamma numbers and the direction $x-1$ of the block, with coefficients below $10^{20}$ (PSLQ, $360$ digits) \computed: in that framework, over $\mathbb Q$, they are new numbers.

(iii) The last minor of \S14.17 is now written in the masters over the stacked Gamma (Proposition~\ref{p:stackedhalf}):
\begin{multline*}
P_{\binom{[5]}1}(\one;1,1)\,x^{(1)}(2)-P_{\binom{[5]}1}(2\one;2,2)\,x^{(1)}(1)\\
=\frac{65\,x^{(1)}(1)-32\,\Gamma(\frac12)^8/x^{(1)}(1)}{1125}\,P_{\binom{[5]}1}(\one;1,1)-x^{(1)}(1)\,P_{\binom{[5]}1}(2\one;2,2)\\
=4.34000293153566203407\ldots,
\end{multline*}
and it is not a rational combination of $1$, $\Gamma(\frac12)^4$, $\Gamma(\frac12)^8$, $\Tx{3}{\IZ}{0}$ and seven products of $x^{(1)}(1)^{\pm1}$ or $x^{(1)}(1)^{\pm2}$ with powers of $\Gamma(\frac12)^4$, with coefficients below $10^{15}$ \computed. Its scale is the masters' scale; nothing smaller is missing.

(iv) The rational words and the pole columns of Chapter~6 did not enter: every word of the centre ray is a polynomial, read at index zero. What closed the live Trailmoment is the shifted inputs (Thm.~6.2), the cut relation (Prop.~5.3), summation by parts in its adjoint form (Thm.~6.5, p.~43), and a fixed block at nonpositive index, the counterpart of Thm.~6.9, p.~45.

(v) Aliases, last: $x^{(1)}(1)=8\int_0^\infty tI_0(t)K_0(t)^4\,dt$, whose evaluation by the Gammas at fifteenths is the classical one through the singular modulus of $\mathbb Q(\sqrt{-15})$ \recalled; Theorem~\ref{t:cert} gives the recurrences of the moments $\int_0^\infty t^{2s-1}K_0(t)^n\,dt$, $n\le7$, by adding marks.
\end{remark}

\subsection*{15.8\quad Open}

\begin{enumerate}[label=(\roman*)]
\item The certificate for every $n$. Theorem~\ref{t:cert} for all $n$ would derive every banana ray law from the stacked Gamma, the law of each list produced by the one before. At each $n\le6$ the same $A_s$ appeared; an operator proof for general $n$: \notyet. From five to six marks the symbol of the five-mark law vanishes at $\lambda=1$ (a native value of five marks), $L_5^\ast$ raises degrees by less than four, and the block of six marks is built differently: \notyet. A mark added in parallel to a mark of another list: for $\binom{[4]}2$ the class decouples and the second row of \S15.4 is the square of the three-mark ray (Proposition~\ref{p:u24}); in general the class enters through the law of its parallel cell (\S16).
\item The direction $x-1$ of the block, $\Tx{0}{(x-1)\,p\,\IZ}{0}=(q\partial_q)G_0(q)|_{q=1}=\int dy/(X-1)^2=0.16106970227667485885\ldots$, which the ray reaches only through the value term of its second master. Expanding $h$ as in Proposition~\ref{p:cosetmark}, it is the lowered setting with $h$ replaced by $y\,\partial_yh$: $\int_0^\infty g^4\,y\,h'(y)\,dy$ \derived\ (alias $4\int_0^\infty t^2I_1(t)K_0(t)^4dt$, to $25$ digits \computed). Which member of the family it is, and whether a lean of the five-mark list reads it as a period: \notyet
\item Three marks: reduce the live Trailmoment of the receipts of Proposition~\ref{p:addmark} at $n=2$ to the Trailmoment of the odd periodic density modulo $3$ at index $2$ (alias $3L(\chi_{-3},2)$) inside the family; Thm.~5.17 does not reach it, the parity being mismatched. \notyet
\item The three masters of $\binom{[5]}2$: no identification yet. \notyet
\item The second corners, from the list. The boundary term at $\sigma=-1$ is a constant term of the continued period (Thm.~1.6), assembled from its Laurent coefficients of orders $-(n-1)$, $-(n-2)$, $-(n-3)$ at $\sigma=0$ and $\sigma=-1$. The leading ones are the corner and the corner times $|\cB|$ (Proposition~\ref{p:negcorner}); the next vanishes at $0$ (Proposition~\ref{p:corner}) and the ray law fixes it at $-1$. On the share side the third ones are Laurent data of the Gammas at the lattice points of the closing, which are heads of the counting input and Trailmoments of the counting input cut at the data (Prop.~5.13, p.~32; Prop.~5.14, p.~33); the rationality of the two boundary terms says that the Trailmoments cancel and the heads remain. Carrying this out: \notyet. Also: which lists put $(\sigma+1)^2$ into both $c_0$ and $c_1$. Here they are the two triangles $\binom{[3]}2$ with parallel classes, and not $\binom{[4]}2$ with a mark doubled. \notyet
\item The last minor of \S14.17, $M_5(1)x^{(1)}(2)-M_5(2)x^{(1)}(1)=4.34000293153566\ldots$, is the minor of the solutions of native values $\frac1{25}$ and $\frac19$, so it is the minimal solution of the exterior square of the ray law, of native value $\frac1{225}$ \derived. It is now the masters over the stacked Gamma (Remark~\ref{r:found}(iii)), so a relation of degree two fixing it would be a relation between the masters and $x^{(1)}(1)$. The exact pairing of Proposition~\ref{p:pairing} is a relation of degree one between solutions, fixed by the residue at zero shares. Whether pairings of the sign classes with closings of the share side twisted by the signs fix relations of degree two such as this minor: \notyet
\item The pairing for the other lists: the partner of the period in the adjoint law, when the closings are not. \notyet
\item Lean rays. At a lean with uniform masses only $\eta=\one$ stays native for $\binom{\cE}1$: at a sign the native equation reads $(m-t)K(\eta)=(s-t)\eta_e$ for every $e$ \derived. The gauges that move off the signs and their values. \notyet
\item Families of native gauges. In both cases of Proposition~\ref{p:families} the letter collapses at the gauge to the letter of a minor that leaves directions free: a class whose contraction letter vanishes is free, and a class whose own letter vanishes leaves a deletion with coloops, whose relative scales are free. A rule for every list, and the exponents $-3$ and $-\frac52$ of the remark on local exponents. \notyet
\item A third flavour: $\mathrm M(K_4)$ on six marks, an atom, whose full cobasis lean $(0,2)$ at unit masses is $6\,\Tx{3}{\IZ}{0}$ (alias $6\zeta(3)$) \recalled; its centre ray. \notyet
\end{enumerate}


\renewcommand{\thesection}{16}
\setcounter{proposition}{0}
\renewcommand{\theproposition}{16.\arabic{proposition}}

\section{The five-mark lists along their hinge trees}

A connected list is glued from prime tables and series and parallel nodes at its bunches of width one (Thm.~20.9, p.~166), and its canonical hinge tree (Def.~22.14, p.~189) is read off the Skeleton: the width $c(S)=r(S)+r(\cE\setminus S)-r(\cE)$ of Def.~7.10, p.~49, is a count of the circuit and cocircuit layers, and a bunch of width one is where the list composes itself out of two pieces. The period is read from the leaves up (Thm.~22.15, p.~190). A cell enters only through its hinge wall and its law (Thm.~22.7, p.~185), a series node multiplies the transforms of its cells and a parallel node their dual transforms (Thm.~22.13, p.~188), and a mark left in place is a one-mark cell whose two transforms at the centre are the Gramma $\sG_{0;(1,-1)}(1,y)$ (Def.~22.17(1), p.~191). This entry reads the five lists of \S15.4 that way. Three are series-parallel, one carries a prime node that acts as a mark, and one is prime.

\begin{center}
\small
\begin{tabular}{@{}l l@{}}
\toprule
list & canonical hinge tree \derived\\ \midrule
$\binom{[5]}1$ & one parallel node\\
$\binom{[4]}2$, a mark doubled & the prime node $\binom{[4]}2$, a parallel node of three marks at one marker\\
$\binom{[5]}2$ & one prime node\\
$\binom{[3]}2$, two marks doubled & a series node of three marks, parallel nodes of three marks at two markers\\
$\binom{[3]}2$, a mark tripled & a series node of three marks, a parallel node of four marks at one marker\\
\bottomrule
\end{tabular}
\end{center}
The bunches of width one with two marks on each side are the parallel classes, of width $1+r(\cE\setminus P)-r(\cE)=1$; in $\binom{[4]}2$ and $\binom{[5]}2$ every set of two marks has width $2$, so both are prime.

\subsection*{16.1\quad A prime node that is a mark}

\begin{proposition}[The doubled mark of $\binom{[4]}2$]\label{p:u24}\derived; receipts \computed\
Replace a mark $i$ of $\binom{[4]}2$ by a parallel class of $k\ge1$ marks. Then
\[
P_{\binom{[4]}2[i\to k]}(\sigma\one;\sigma,\sigma)=P_{\binom{[3]}1}(\sigma\one;\sigma,\sigma)\;P_{\binom{[k+1]}1}(\sigma\one;\sigma,\sigma).
\]
At $k=1$ this is $P_{\binom{[4]}2}=P_{\binom{[3]}1}\,\Gst21{\sigma,\ \sigma}{2\sigma}$, the agreement of the table of \S15.1 at every strength; at $k=2$ it is the second row of \S15.4,
\[
P_{\binom{[4]}2,\ \text{a mark doubled}}(\sigma\one;\sigma,\sigma)=P_{\binom{[3]}1}(\sigma\one;\sigma,\sigma)^2 .
\]
\end{proposition}

\begin{proof}
$\binom{[4]}2$ is identically self-dual. The cell of $\binom{[4]}2$ hinged at $i$ has the hinge minors $\binom{[4]}2\setminus i=\binom{[3]}2$ and $\binom{[4]}2/i=\binom{[3]}1$, dual with the identity and the same weights, so its ratio is identically one, and in every gluing it acts as $\phi(\Delta)$ times one mark, with $\phi(\Delta)=P_{\binom{[3]}2}(\sigma\one;\kappa,2\sigma-\kappa)$ (Cor.~22.8(i),(ii), p.~186). Its hinge wall is $\kappa=\sigma$, balanced against the parallel cell of the class, whose hinge wall is $k\sigma-(k-1)\sigma=\sigma$ (Thm.~22.13(ii)). By Prop.~8.14, p.~57, $P_{\binom{[3]}2}(\sigma\one;\sigma,\sigma)=P_{\binom{[3]}1}(\sigma\one;\sigma,\sigma)$. What is left is the class glued to one mark, the banana $\binom{[k+1]}1$.
\end{proof}

\noindent Receipts \computed: Procedure~\ref{p:proc} for $\binom{[4]}2$ with a mark doubled agrees with $P_{\binom{[3]}1}(\sigma\one;\sigma,\sigma)^2$ to $60$ digits at $\sigma=1,\dots,4$, and for $\binom{[4]}2$ with $P_{\binom{[3]}1}\Gst21{\sigma,\ \sigma}{2\sigma}$ at $\sigma=1,2,3$.

\begin{corollary}\label{c:u24}\derived\
The ray law of $\binom{[4]}2$ with a mark doubled is the symmetric square of the three-mark law: its native values are the products of $\frac19$ and $1$ in pairs, the roots $81,9,1$ of \S15.4, and its corner is $6^2=36$. Every value lies in $\mathbb Q+\mathbb Q\,P_{\binom{[3]}1}(\one;1,1)+\mathbb Q\,P_{\binom{[3]}1}(\one;1,1)^2$: its two masters are the three-mark centre and its square.
\end{corollary}

The same argument at the centre of any identically self-dual list splits off a mark (Cor.~22.8(iii) and Prop.~8.14) \derived:
\[
P_{\binom{[2r]}r}(\sigma\one;\sigma,\sigma)=\Gst21{\sigma,\ \sigma}{2\sigma}\,P_{\binom{[2r-1]}{r-1}}(\sigma\one;\sigma,\sigma).
\]
So $\binom{[6]}3$ carries exactly the masters of $\binom{[5]}2$.

\subsection*{16.2\quad A series root over banana leaves}

\begin{proposition}[The triangle rows]\label{p:triangles}\derived; receipts \computed\
Let $f_{P_k}$ be the transform at the centre of the parallel cell of a class of $k$ marks; $\int_0^\infty y^{w-1}f_{P_k}(y)\,dy=\Gamma(w)^2P_{\binom{[k]}1}((\sigma-w)\one;\sigma-w,\sigma-w)$ for $0<\re w<\sigma$. Then
\begin{gather*}
\Gamma(\sigma)^2P_{\binom{[3]}2,\ \text{a mark tripled}}(\sigma\one;\sigma,\sigma)=\int_0^\infty y^{\sigma-1}\,\sG_{0;(1,-1)}(1,y)^2\,f_{P_3}(y)\,dy\\
\qquad=\frac1{2\pi i}\int_{(c)}\Gamma(d)^2\Gamma(\sigma-d)^2\,P_{\binom{[4]}2}(d\one;d,d)\,dd,\\
\Gamma(\sigma)^2P_{\binom{[3]}2,\ \text{two marks doubled}}(\sigma\one;\sigma,\sigma)=\int_0^\infty y^{\sigma-1}\,\sG_{0;(1,-1)}(1,y)\,f_{P_2}(y)^2\,dy .
\end{gather*}
So the tripled triangle is the closing of Proposition~\ref{p:cosetmark} with the ray of the four-mark banana replaced by the ray of $\binom{[4]}2$.
\end{proposition}

\begin{proof}
The triangle is a series root (Def.~22.17(5), p.~191): its cells are the parallel cells of the classes and the marks left in place, the latter with the transform $\sG_{0;(1,-1)}(1,y)$. The dual transform of a parallel cell is the product of the $k$ Grammas of its marks (Thm.~22.13(ii)), whose Mellin read is $\Gamma(z)^2P_{\binom{[k]}1}(z\one;z,z)$ (Proposition~\ref{p:bananapin}); its hinge wall is $\sigma$, and at $s=t=\kappa=\sigma$ the change of transform reflects $z\mapsto\sigma-z$ (Prop.~22.5(iii),(v), p.~184). The Mellin read of a product pins the sum of the shares, as in the proof of Proposition~\ref{p:bananapin}, and $P_{\binom{[2]}1}P_{\binom{[3]}1}=P_{\binom{[4]}2}$ is Proposition~\ref{p:u24} at $k=1$.
\end{proof}

\noindent Receipts \computed: the closing $-\sum_{d\ge1}(W_s\,P_{\binom{[4]}2})'(d)/\Gamma(s)^2$ of Proposition~\ref{p:cosetmark} agrees with Procedure~\ref{p:proc} for the tripled triangle to $60$ digits at $s=1,2,3$; for two marks doubled, opening the list at its unmarked mark (Prop.~19.19, read as in Proposition~\ref{p:u25split} below: the contraction is the four-mark banana, the deletion two parallel pairs) gives $7.8658373683630$ by three-dimensional quadrature, against $7.86583736836402159\ldots$.

\begin{proposition}[Closing the tripled triangle over its input]\label{p:tripled}\derived; receipts \computed\
Let $q(\sigma)=P_{\binom{[4]}2}(\sigma\one;\sigma,\sigma)$ and $f_q=-q'$.
\begin{enumerate}[label=(\roman*)]
\item $q$ obeys the three-mark law times the shift of the stacked Gamma,
\[
\sigma^3q(\sigma)-2(2\sigma+1)(10\sigma^2+10\sigma+3)\,q(\sigma+1)+36(\sigma+1)(2\sigma+1)(2\sigma+3)\,q(\sigma+2)=0,
\]
with corner $12=2\cdot6$ and no term of order $\sigma$, and its symbol at $\lambda=1$ is $105\neq0$.
\item Summation by parts at the cut $0$, as in Proposition~\ref{p:sbp}: for every polynomial $R$,
\begin{gather*}
\Tx{0}{(L_q^\ast R)\,q\,\IZ}{0}=-12\,R(0),\\
\Tx{0}{(L_q^\ast R)\,f_q\,\IZ}{0}=\Tx{0}{(L_q'^\ast R)\,q\,\IZ}{0}+36\,R(-1)\,P_{\binom{[3]}1}(\one;1,1).
\end{gather*}
The last coefficient of the law has only a simple zero at $\sigma=-1$: the shift of the stacked Gamma, $P_{\binom{[2]}1}(\sigma+2)/P_{\binom{[2]}1}(\sigma+1)=(\sigma+1)/(2(2\sigma+3))$, cancels one of the two factors $\sigma+1$ of the three-mark coefficient $9(\sigma+1)^2$. So the slope boundary there reads $q(1)=P_{\binom{[3]}1}(\one;1,1)$.
\item Every value of the tripled triangle is a rational number, plus a rational multiple of $P_{\binom{[3]}1}(\one;1,1)$, plus a rational combination of the three slope Trailmoments $-\Tx{0}{(Rq)'\,\IZ}{0}$, $R\in\{1,\,x,\,x^2/2\}$. The words $W_s$, $s\le8$, reach all three; the five-mark banana reached two (Corollary~\ref{c:masters}).
\item In these coordinates the rational relation of \S15.4 for this row, the second corner, is the relation among $1,P(1),\dots,P(5)$, coordinate by coordinate.
\end{enumerate}
\end{proposition}

\begin{proof}
(i) $\sigma\,P_{\binom{[2]}1}(\sigma)=2(2\sigma+1)\,P_{\binom{[2]}1}(\sigma+1)$ (Prop.~1.3(iii)) multiplies the three-mark law coefficient by coefficient into $\sigma$ times the display. The corner is the product of the corners $\lim\sigma P_{\binom{[2]}1}=2$ and $6$; $\sigma\Gst21{\sigma,\ \sigma}{2\sigma}=2\Gamma(1+\sigma)^2/\Gamma(1+2\sigma)$ and $\sigma^2P_{\binom{[3]}1}$ have no terms of order $\sigma$ (Proposition~\ref{p:corner}). (ii) The move of Proposition~\ref{p:sbp}: the value boundary is $-6q(1)+108q(2)=-12$; for the slopes the boundary at the cut $0$ is the term of order $\sigma$, zero, and the one at $-1$ is $-R(-1)\,c_2'(-1)\,q(1)$ with $c_2'(-1)=-36$. (iii) As in Proposition~\ref{p:block}, with $\mathrm{span}\{1,x,x^2\}$ as the complement of $L_q^\ast(\mathbb Q[x])$, and Proposition~\ref{p:cosetmark} for the values. (iv) Exact arithmetic on the eight coordinates.
\end{proof}

\noindent Receipts \computed. The coordinates reproduce Procedure~\ref{p:proc} to $60$ digits at $s\le5$. There is no rational relation among $1$, the three-mark centre $P_{\binom{[3]}1}(\one;1,1)$, its slope $\partial_\sigma P_{\binom{[3]}1}(\sigma\one;\sigma,\sigma)|_{\sigma=1}$ and the six Trailmoments
\[
\Tx{0}{x^jq\,\IZ}{0},\qquad \Tx{0}{x^jf_q\,\IZ}{0}\qquad(j\le2)
\]
with coefficients below $10^{25}$, and none among $1$, the three-mark centre and the three slope Trailmoments of (iii) below $10^{40}$ (PSLQ, $300$ digits). So the four masters of the tripled triangle are the three-mark centre, entering at the boundary $\sigma=-1$, and the three slope Trailmoments.

\subsection*{16.3\quad The prime node at a mark}

\begin{proposition}[$\binom{[5]}2$ opened at a mark]\label{p:u25split}\derived; receipt \computed\
For a mark $e$ of $\binom{[5]}2$ the deletion is $\binom{[4]}2$ and the contraction $\binom{[4]}1$. Write $X_{\setminus}=X_{\binom{[4]}2}$ and $X_{/}=X_{\binom{[4]}1}$ for their strides on the chart of the other four marks. Then
\[
P_{\binom{[5]}2}(\sigma\one;\sigma,\sigma)=\int dy\int_0^1w^{\sigma-1}(1-w)^{\sigma-1}\big(w\,X_{\setminus}+(1-w)\,X_{/}\big)^{-\sigma}dw,
\]
and on the share side of the shape,
\[
\Gamma(\sigma)^2P_{\binom{[5]}2}(\sigma\one;\sigma,\sigma)=\frac1{2\pi i}\int_{(c)}\Gamma(d)^2\Gamma(\sigma-d)^2\,\Phi_\sigma(d)\,dd,\qquad \Phi_\sigma(d)=\int X_{/}^{-d}X_{\setminus}^{-(\sigma-d)}\,dy,
\]
$0<c<\sigma$, with $\Phi_\sigma(d)=P_F\big(\sigma\one;(d,\sigma-d),(d,\sigma-d)\big)$ the centre of the family $F$ of the two children (Def.~8.1, p.~53). At $\sigma=1$,
\[
P_{\binom{[5]}2}(\one;1,1)=\int\frac{\log X_{\setminus}-\log X_{/}}{X_{\setminus}-X_{/}}\,dy .
\]
\end{proposition}

\begin{proof}
The mark is the Breta of its children's even letters (Prop.~19.19, p.~144): at the centre the prefactor and the stacked Gamma are one, and $E_{\setminus e}=\alpha^{\cE\setminus e}X_{\setminus}$, $E_{/e}=\alpha^{\cE\setminus e}X_{/}$. Splitting the coefficient $wX_{\setminus}+(1-w)X_{/}$ (Cor.~4.4, p.~26) and integrating the shape gives the Gammas, one pair for each child.
\end{proof}

\noindent Receipt \computed: the last integral is $4.0297141707108$ by three-dimensional quadrature, against $4.02971417071104960732\ldots$ of \S15.4.

\begin{proposition}[The contraction under the deletion]\label{p:under}\derived\
On the chart $X_{\setminus}-X_{/}=(e_2^2-e_1e_3)/e_4$ with
\[
e_2^2-e_1e_3=\sum_{i<j}\alpha_i^2\alpha_j^2+\sum_i\sum_{\{j,k\}\not\ni i}\alpha_i^2\alpha_j\alpha_k+2\,\alpha_1\alpha_2\alpha_3\alpha_4,
\]
a polynomial with positive coefficients, so $X_{/}<X_{\setminus}$ everywhere. The walls of $F$ at the strengths $(d,1-d)$ are $1$ at $|S|$ odd and $2-d$ at $|S|=2$, so $\Phi=\Phi_1$ is holomorphic on $\re d<2$, and closing to the left (Thm.~4.5, Lemma~5.16),
\[
P_{\binom{[5]}2}(\one;1,1)=\sum_{m\ge0}\Phi'(-m),\qquad \Phi(-m)=\int X_{/}^{\,m}\,X_{\setminus}^{-1-m}\,dy :
\]
the slopes of the deletion's centre ray with the contraction's even letter inserted $m$ times.
\end{proposition}

\begin{proof}
$X_{/}=e_1e_3/e_4$, $X_{\setminus}=e_2^2/e_4$, and the polynomial is read off the monomials of $e_2^2$ and $e_1e_3$. The widths are $1$ for every $S$ in $\binom{[4]}1$ and $1,2,1$ at $|S|=1,2,3$ in $\binom{[4]}2$. On the line $\re d=-N-\frac12$,
\[
|\Phi(d)|\le\int\big(X_{/}/X_{\setminus}\big)^{N+1/2}X_{\setminus}^{-1}\,dy\to0,
\]
so the left closing applies; at $d=-m$ the kernel $\Gamma(d)^2\Gamma(1-d)^2$ of Lemma~5.16 has the principal part $(d+m)^{-2}$ and no simple part. Equivalently, expand $1/(X_{\setminus}-X_{/})=\sum_mX_{/}^mX_{\setminus}^{-m-1}$ in the last display of Proposition~\ref{p:u25split}.
\end{proof}

The ratio $X_{/}/X_{\setminus}$ is largest at the table's gauge, where it is $\frac{16}{36}=\frac49$, so the closing converges like $(\frac49)^m$ (alias: Newton's inequality $e_1e_3\le\frac49e_2^2$) \recalled.

\begin{proposition}[The values at the lattice points are leans of the three-mark banana]\label{p:leans}\derived; receipts \computed\
For $m\in\mathbb N$, with $(e_1e_3)^m=\sum_vc_v\alpha^v$,
\[
\Phi(-m)=\sum_vc_v\,P_{\binom{[4]}2}\big(\one+v;1+m,1+m\big),
\]
and for masses $\bmu$ on the balance wall and every mark $g$,
\[
P_{\binom{[4]}2}(\bmu;s,s)=\frac1{\Gamma(2s)}\int_0^\infty y^{\mu_g-1}\prod_{e\neq g}\sG_{2s-\mu_g-\mu_e;(1,-1)}(1,y)\,dy :
\]
$\Phi(-m)$ is a finite sum of leans of the three-mark banana, Grammas of integer index read three at a time. The first values are
\begin{gather*}
\Phi(0)=P_{\binom{[3]}1}(\one;1,1),\qquad \Phi(-1)=\frac{2+2P_{\binom{[3]}1}(\one;1,1)}9,\\
\Phi(-2)=\frac{71+56P_{\binom{[3]}1}(\one;1,1)}{810}\quad\computed .
\end{gather*}
\end{proposition}

\begin{proof}
On the chart $X_{/}^mX_{\setminus}^{-1-m}=\alpha^{\one}(e_1e_3)^m\,e_2^{-2(1+m)}$, and $e_2$ is both letters of $\binom{[4]}2$; each monomial $\alpha^v$ is an insertion moving the masses (Lemma~35.2(iii), p.~308). Then Cor.~22.8(iii) at $g$, $P_{\binom{[4]}2}(\bmu;s,s)=\Gst21{\mu_g,\ 2s-\mu_g}{2s}P_{\binom{[3]}2}(\bmu|;2s-\mu_g,\mu_g)$; Prop.~8.14 turns $\binom{[3]}2$ into $\binom{[3]}1$ at the exponents $(\mu_g,2s-\mu_g)$; and the share side of the three-mark banana at a lean, $\Gamma(a)\Gamma(b)P_{\binom{[3]}1}(\bmu';a,b)=\int_0^\infty y^{a-1}\prod_e\sG_{b-\mu'_e;(1,-1)}(1,y)\,dy$ (Prop.~4.2 with the pins $a_e-b_e=\mu'_e-b$, as in Proposition~\ref{p:bananapin}), whose Gammas cancel the Beta corner.
\end{proof}

\noindent Receipts \computed: $\Phi(-1)$ by three-dimensional quadrature is $0.7430904974869$, against $\frac29(1+P_{\binom{[3]}1}(\one;1,1))=0.74309049748654642013\ldots$; the leans themselves by one-dimensional quadrature to $40$ digits, with PSLQ for the rational coefficients; on the way, $\int_0^\infty\sG_{0;(1,-1)}(1,y)\,\sG_{1;(1,-1)}(1,y)^2\,dy=\frac13$.

\begin{remark}[What the prime node reads]
The prime node reads the three-mark banana at leans. Its values at the lattice points land on $\mathbb Q+\mathbb Q\,P_{\binom{[3]}1}(\one;1,1)$, and its period at strength one is the sum of the slopes there: the division of Propositions~\ref{p:sbp} and~\ref{p:block}, values on lattice inputs and slopes carrying the masters, now along the deletion's ray with the contraction's even letter as its stride, of native value $\frac49$ at the table's gauge. The three masters of $\binom{[5]}2$ are slope Trailmoments of that ray, as the two of $\binom{[5]}1$ were of the four-mark ray; the law in $m$ and its closing: \notyet.
\end{remark}

\begin{remark}[What the tree decided]
A parallel node multiplies Grammas: the banana. An identically self-dual prime node is a mark of ratio one, and the period factors: the doubled mark of $\binom{[4]}2$ is the square of the three-mark ray, with no new number. A series root over parallel cells multiplies transforms; its closing has the banana's shape over a product input, and the stacked Gamma in that input's law opens the boundary at $\sigma=-1$, which is where the second corner of the tripled triangle comes from (for two marks doubled: \notyet). A prime node that is not self-dual is opened at a mark into its children, weighed by the Breta of their even letters, and reads leans. The three-mark centre enters the rows whose trees carry a three-mark piece, the class of three or the deletion $\binom{[3]}2$ inside $\binom{[4]}2$: rows two, three and five. The row with two marks doubled has only pairs and one mark at its leaves; what its masters are: \notyet. Aliases, last: $P_{\binom{[3]}1}(\one;1,1)=3L(\chi_{-3},2)$, and $\sG_{\nu;(1,-1)}(1,y)=2y^{\nu/2}K_\nu(2\sqrt y)$, so the leans are moments of products of three Bessel functions $K_\nu$ of integer orders.
\end{remark}

\subsection*{16.4\quad Open}

\begin{enumerate}[label=(\roman*)]
\item The ray of $\binom{[4]}2$ with the stride $X_{/}/X_{\setminus}$: the law of $\Phi(-m)$ in $m$, the summation by parts of its slope input, and with it the three masters of $\binom{[5]}2$ (item (iii) of \S15.8). \notyet
\item Two marks doubled: the closing over its input of two parallel cells, and its second corner as a boundary at $\sigma=-1$. \notyet
\item The certificate of Theorem~\ref{t:cert} for a series root: there the input was one ray, here it is a product, and the laws have orders $6$ and $5$ (\S15.4). \notyet
\item $\mathrm M(K_4)$ is self-dual only up to a relabelling of its marks, so Cor.~22.8(iii) does not split it (as the handbook notes after Cor.~22.8); opened at a mark its children are $K_4\setminus e$ and $K_4/e$, the latter the triangle with two marks doubled of \S15.4. Its centre ray read this way. \notyet
\end{enumerate}


\renewcommand{\thesection}{17}
\setcounter{proposition}{0}
\renewcommand{\theproposition}{17.\arabic{proposition}}

\section{The masters of $\binom{[5]}2$; two marks doubled, closed}

Both lists are read at a mark (Prop.~19.19, p.~144): the period is the Breta of the even letters of the two children, its share side is a Mellin--Barnes integral over the split $d$ with the kernel of Lemma~5.16, p.~34, and the closing (Thm.~4.5, p.~26) reads a values input and a slope input on the lattice of the split. For $\binom{[5]}2$ the children are $\binom{[4]}2$ and $\binom{[4]}1$ (\S16.3). For the triangle with two marks doubled the mark is taken in a doubled class; the children are then those of the triangle $H$ with one class doubled, opened at its single mark: the three-mark banana $H/1$ and a pair with a coloop $H\setminus1$, and the class enters through its Mellin read. In both lists the lattice values of the input are leans of the three-mark banana. This entry finds the law of each input along its lines, the boundaries of its slope input, and the masters. Throughout, $\Pth$ is the three-mark centre and $\Pfour$ the four-mark centre.

\subsection*{17.1\quad The input of $\binom{[5]}2$ along its lines}

For $s\in\{1,2,\dots\}$ put $\Phi_s(d)=\Phi(d,s-d)=\int X_{/}^{-d}X_{\setminus}^{-(s-d)}\,dy$, the family of Proposition~\ref{p:u25split} on the line $s$; $\Phi_1$ is the $\Phi$ of Proposition~\ref{p:under}. By Proposition~\ref{p:leans} every lattice value is a finite sum of leans of the three-mark banana,
\[
\Phi_s(-m)=r_s(m)+c_s(m)\,\Pth,\qquad r_s(m),\,c_s(m)\in\mathbb Q\quad(m\in\mathbb N),
\]
and $c_s$ is the \emph{three-mark coordinate} of the values input.

\begin{proposition}[The closing at every integer strength]\label{p:u25clos}\derived\
For $s\in\{1,2,\dots\}$ let $V_s(m)=\big((m+1)(m+2)\cdots(m+s-1)\big)^2$. Then
\[
\Gamma(s)^2P_{\binom{[5]}2}(s\one;s,s)=\sum_{m\ge0}\Big(V_s(m)\,\Phi_s'(-m)-V_s'(m)\,\Phi_s(-m)\Big),
\]
and $x_s=\Gamma(s)^{-2}\sum_{m\ge0}V_s(m)\,\Phi_s(-m)$ is the lowered setting, the values half of the closing.
\end{proposition}

\begin{proof}
Proposition~\ref{p:u25split} and the left closing of Proposition~\ref{p:under} at the strength $s$. By the reflection (Prop.~1.3(v), p.~11), $\Gamma(d)^2\Gamma(s-d)^2=V_s(-d)\,\Gamma(d)^2\Gamma(1-d)^2$, with double poles at $d=-m$ and no simple parts; the residue at $d=-m$ is $\big(V_s(-d)\Phi_s(d)\big)'$ there.
\end{proof}

At $s=1$ this is Proposition~\ref{p:under}.

\begin{proposition}[The law on a line]\label{p:u25law}\derived\ in Proposition~\ref{p:certs}\
On the line $s=1$,
\[
\sum_{i=0}^4c_i(m)\,\Phi_1(-m-i)=0,
\]
\begin{align*}
c_0&=-8(m+1)^3(m+2)(m+3)(7m+16),\\
c_1&=(m+2)(m+3)(630m^4+4583m^3+12351m^2+14584m+6336),\\
c_2&=-(m+3)(2m+5)(1295m^4+11241m^3+35869m^2+49386m+24408),\\
c_3&=3(2m+5)(2m+7)(385m^4+3855m^3+13913m^2+21153m+11286),\\
c_4&=-54(m+2)(m+4)(2m+5)(2m+7)(2m+9)(7m+9),
\end{align*}
and $r_1$, $c_1$ obey it separately. Its symbol is
\[
216\rho^4-330\rho^3+185\rho^2-45\rho+4=(2\rho-1)(9\rho-4)(3\rho-1)(4\rho-1).
\]
The lines $s=2,3$ have laws of order four and degree eight with the same symbol.
\end{proposition}

\noindent Receipts \computed. Found by exact linear algebra on the values $\Phi_s(-m)$, $m\le80$, of Proposition~\ref{p:leans} (at $s=1$: $138$ equations, $35$ unknowns, one solution). At $m=-1$, outside the data, $c_0$ has a triple zero and the law holds exactly, $\sum_{i\ge1}c_i(-1)\Phi_1(1-i)=0$ with $(c_1,\dots,c_4)(-1)=(300,-5670,25920,-34020)$. It was read off first; Proposition~\ref{p:certs} derives it from the relations along a ray.

\begin{proposition}[The roots are the native values of the ratio of the strides]\label{p:u25roots}\derived\
On the chart $X_{/}/X_{\setminus}=e_1e_3/e_2^2$. At the table's gauge its value is $\frac49$, the maximum. On the pattern $(1,1,x,x)$ its native equation is $(1-x)(1+x^2)=0$, and at $(1,1,i,i)$ the value is $\frac12$. Where one mark goes to zero it becomes $e_1e_3/e_2^2$ of three marks, of maximum $\frac13$; where one mark goes to zero and another to infinity it becomes $\alpha_1\alpha_2/(\alpha_1+\alpha_2)^2$, the reciprocal stride of the two-mark banana, of maximum $\frac14$. These four values are the roots of the symbol of Proposition~\ref{p:u25law}.
\end{proposition}

The values input decays like $(\frac49)^m$; its coordinates $r_1$ and $c_1$ decay like $(\frac12)^m$, and
\[
\Pth=-\lim_{m\to\infty}\frac{r_1(m)}{c_1(m)}\quad\computed:
\]
the three-mark centre is the multiplier that removes the component along the native value $\frac12$ of the complex gauge.

\begin{proposition}[Lines tied by a first-order relation]\label{p:u25mix}\derived\ for $s\le4$ in Proposition~\ref{p:certs}\
With $F(m,n)=\Phi(-m,n)$,
\begin{multline*}
4(m+1)(3m+2n+4)\,F(m,n+1)-12(n+1)(2n+3)\,F(m+1,n+2)-(m+1)(2n-m)\,F(m,n)\\
+(6n^2+4mn-3m^2+16n-4m+5)\,F(m+1,n+1)-6(n+1)(2n+3)\,F(m+2,n+2)=0.
\end{multline*}
\end{proposition}

On the lines $n=m+s$ it ties the line $s+1$ to the line $s$, so with the law every line is read off the line $1$. Receipts \computed: found at $s\le4$, the coefficients polynomial in $s$; the lines $2$ to $5$ read off the line $1$ agree with Proposition~\ref{p:leans} at every $m\le80$.

\subsection*{17.2\quad The slope input and its boundaries}

The law holds as a function of the split (Proposition~\ref{p:certs}), so its derivative is the recurrence of the slope input,
\[
\sum_{i=0}^4c_i(m)\,\Phi_1'(-m-i)=\sum_{i=0}^4c_i'(m)\,\Phi_1(-m-i),
\]
as in Proposition~\ref{p:liverec}. Its four initial slopes $\Phi_1'(0),\dots,\Phi_1'(-3)$ are tied by the boundaries of the law.

\begin{proposition}[The cut]\label{p:u25cut}\derived\ from Proposition~\ref{p:u25law}\
At $m=-1$, where the four-mark centre $\Phi_1(1)=\Pfour$ sits, $c_0$ has a triple zero; neither $\Phi_1(1)$ nor its slope enters, and
\[
300\,\Phi_1'(0)-5670\,\Phi_1'(-1)+25920\,\Phi_1'(-2)-34020\,\Phi_1'(-3)=-\frac{21380}{63}+\frac{8866}{63}\,\Pth .
\]
The lines $2$ and $3$ give $\frac12$ and $-\frac59$ times the same relation.
\end{proposition}

\begin{proposition}[The wall at $d=2$]\label{p:u25wall}residue \derived; relation \computed\
$\Phi_1$ has a simple pole at $d=2$, the wall $2-d$ of the two-element $S$ (Proposition~\ref{p:under}), with residue $-6$: each of the six two-element $S$ collapses (Cor.~8.8, p.~56) to a product of two two-mark centres, $P_{\binom{[2]}1}(\one;1,1)^2=1$; the law at $m=-2$ returns the same $-6$. There the law meets the slopes only through
\[
c_2(-2)\,\Phi_1'(0)+c_3(-2)\,\Phi_1'(-1)=96\,\Phi_1'(0)-432\,\Phi_1'(-1),
\]
and
\[
6\,\Phi_1'(0)-27\,\Phi_1'(-1)=7\,\Pth-22
\]
(PSLQ, $140$ digits), which is the same as the finite part
\[
\Phi_1(d)=-\frac{6}{d-2}+\frac{21}2+\frac34\,\Pfour+O(d-2).
\]
\end{proposition}

\begin{proposition}[One free slope, along the three-mark coordinate]\label{p:u25free}\computed\
A third relation, $1680\,\Phi_1'(0)-24300\,\Phi_1'(-2)=2308\,\Pth-7207$ (PSLQ), with Propositions~\ref{p:u25cut} and~\ref{p:u25wall}, says that $\Phi_1'(-m)-\Phi_1'(0)\,c_1(m)\in\mathbb Q+\mathbb Q\,\Pth$ for $m\le3$. By the slope recurrence and Proposition~\ref{p:u25mix} it holds for every $m$ and every line:
\[
\Phi_s'(-m)=\rho_s(m)+\lambda_s(m)\,\Pth+c_s(m)\,\Phi_1'(0),\qquad \rho_s(m),\,\lambda_s(m)\in\mathbb Q .
\]
The one free direction of the slope input is the three-mark coordinate of the values input. Minimality fixes it: the slope input has no component along the native value $\frac12$.
\end{proposition}

\noindent Receipts \computed. With $\Phi_1'(0)$ from the masters below, $\Phi_1'(-m)/\Phi_1(-m)=0.81342$ at $m=600$, tending to $\log\frac94=0.81093$, where a component along $\frac12$ would grow like $(\frac98)^m$. The free slope is
\begin{gather*}
\Phi_1'(0)=\int\log\frac{X_{\setminus}}{X_{/}}\,\frac{dy}{X_{\setminus}}=2.794313984758270846154613553026697141639\ldots,\\
\Phi_1'(0)=2\,\Pth-\partial_\sigma P_{\binom{[3]}1}(\sigma\one;\sigma,\sigma)\big|_{\sigma=1}-\int\log X_{/}\,\frac{dy}{X_{\setminus}},\\
\int\log X_{/}\,\frac{dy}{X_{\setminus}}=9.9686527561522909453181622662\ldots,
\end{gather*}
by three-dimensional quadrature $2.7943139847583$; and by the same quadrature
\[
\Phi_1'(-1)=0.8280938236176\quad\text{against}\quad0.82809382361938566084\ldots .
\]
There is no rational relation among $1$, $\Pth$, its slope, $\Pfour$, the four-mark slope and $\Phi_1'(0)$ with coefficients below $10^{15}$.

\subsection*{17.3\quad The masters of $\binom{[5]}2$}

\begin{theorem}[The masters of $\binom{[5]}2$]\label{t:u25masters}(i) \derived\ from Propositions~\ref{p:u25clos} and~\ref{p:u25free}; (ii)--(iv) \computed\
\begin{enumerate}[label=(\roman*)]
\item Put
\begin{gather*}
A_s=\frac{\sum_m\big(V_s\rho_s-V_s'r_s\big)(m)}{\Gamma(s)^2},\quad B_s=\frac{\sum_m\big(V_s\lambda_s-V_s'c_s\big)(m)}{\Gamma(s)^2},\\
C_s=\frac{\sum_mV_s(m)c_s(m)}{\Gamma(s)^2},\quad R_s=\frac{\sum_mV_s(m)r_s(m)}{\Gamma(s)^2},
\end{gather*}
each convergent, the coordinates decaying like $(\frac12)^m$. Then
\[
P_{\binom{[5]}2}(s\one;s,s)=A_s+B_s\,\Pth+C_s\,\Phi_1'(0),\qquad x_s=R_s+C_s\,\Pth :
\]
the same $C_s$ weighs the three-mark centre in the lowered setting and the free slope in the period.
\item $C_1=\dfrac{\sqrt3}{80}\Gst44{\tfrac1{20},\ \tfrac3{20},\ \tfrac7{20},\ \tfrac9{20}}{\tfrac12,\ \tfrac12,\ \tfrac12,\ \tfrac12}$ \ ($200$ digits).
\item The ray law of \S15.4 and its corner are the boundary at the cut. As functions of the four initial slopes,
\begin{gather*}
\sum_{i=1}^4c^{\binom{[5]}2}_i(0)\,P(i)-1800=-\tfrac{55}{27}\,\beta,\qquad\sum_{i=0}^4c^{\binom{[5]}2}_i(1)\,P(1+i)=-\tfrac{898}{189}\,\beta,
\end{gather*}
with $\beta$ the left side minus the right side of Proposition~\ref{p:u25cut}, $P(i)=P_{\binom{[5]}2}(i\one;i,i)$, and $1800=60\cdot30$ the corner times $-[\sigma^4]c_0$. The lowered setting obeys the ray law with no corner.
\item A second corner lands on the three-mark centre:
\[
20\,P(1)-1672\,P(2)+5120\,P(3)=384-135\,\Pth .
\]
It is the ray law at $\sigma=-1$: $(c_2,c_3,c_4)(-1)=40\cdot(20,-1672,5120)$ and $\CT_{\sigma=-1}\big(c_0(\sigma)P(\sigma)+c_1(\sigma)P(\sigma+1)\big)=5400\,\Pth-15360$. In the coordinates of (i) it reads $5C_1-418C_2+1280C_3=0$, so the free slope drops out.
\end{enumerate}
Over $\mathbb Q+\mathbb Q\,\Pth$ the masters of $\binom{[5]}2$ are $P(1)$ and $P(2)$; the third master of \S15.4 is the three-mark centre.
\end{theorem}

\begin{proof}
(i) Insert Proposition~\ref{p:u25free} in Proposition~\ref{p:u25clos}; every coordinate is a combination of solutions of the law with the native values $\frac12,\frac49,\frac13,\frac14$, so each sum converges. (ii)--(iv) are identities among the sums of (i), checked to $140$ to $200$ digits.
\end{proof}

\begin{center}
\small
\begin{tabular}{@{}l l l@{}}
\toprule
& $s=1$ & $s=2$\\ \midrule
$A_s$ & $1.2617550081051847724190292271463942\ldots$ & $-0.3171370910356445997740190032938055\ldots$\\
$B_s$ & $-0.4062466940152853736690239421082401\ldots$ & $-0.0530255881622819812253694257861745\ldots$\\
$C_s$ & $1.3313334684984500811824592146018574\ldots$ & $0.1607947380347384161148704355054792\ldots$\\
$R_s$ & $0.3527480810189285064734282363418756\ldots$ & $-0.2668464215995376310247031065607639\ldots$\\
$P(s)$ & $4.0297141707110496073206957344284069\ldots$ & $0.0078868342010287786455991517685324\ldots$\\
$x_s$ & $3.4732702349419903397503533563416278\ldots$ & $0.1100415288232609996435971271184530\ldots$\\
\bottomrule
\end{tabular}
\end{center}

\noindent Receipts \computed. The four initial slopes solved from $P(1),P(2),P(3)$ and Proposition~\ref{p:u25cut} give $P(4)$ and $P(5)$ to $154$ digits against Procedure~\ref{p:proc}, and the relations of Propositions~\ref{p:u25wall} and~\ref{p:u25free} then come out of PSLQ with coefficients below $10^5$. Apart from (iv) there is no rational relation among $1$, $\Pth$, $\Phi_1'(0)$, $P(1)$, $P(2)$, $x_1$, $x_2$ with coefficients below $10^8$; $R_1$, $C_2$, $R_2$, $C_3$ are not in $\mathbb Q\,aC_1+\mathbb Q\,b\,\Gamma(\frac12)^2+\mathbb Q\,c\,\Gamma(\frac12)^4/C_1+\mathbb Q+\mathbb Q\,\Pth$ for any $a,b,c\in\{1,\sqrt3,\sqrt5,\sqrt{15}\}$, coefficients below $10^{12}$.

\subsection*{17.4\quad Two marks doubled, closed}

Let $H$ be the triangle with the class $2$ doubled. Opened at its mark $1$ (Prop.~19.19), its children on the chart of $\{2a,2b,3\}$ are the three-mark banana $H/1=\binom{[3]}1$, of stride $X_{H/1}=(B+\alpha_3)(B'+\alpha_3^{-1})$, and $H\setminus1$, the pair $\{2a,2b\}$ with $3$ a coloop, of stride $X_{H\setminus1}=BB'$, where $B=\alpha_{2a}+\alpha_{2b}$ and $B'=\alpha_{2a}^{-1}+\alpha_{2b}^{-1}$. As $X_{H/1}=BB'+B\alpha_3^{-1}+B'\alpha_3+1$, the deletion lies under the contraction. Put $\Psi(a,b)=\int X_{H/1}^{-a}X_{H\setminus1}^{-b}\,dy$ and $\Psi_s(d)=\Psi(d,s-d)$.

\begin{proposition}[Opened at a mark of a doubled class]\label{p:r4clos}\derived\
For the triangle $R$ with two marks doubled, $0<c<\sigma$,
\[
\Gamma(\sigma)^2P_R(\sigma\one;\sigma,\sigma)=\frac1{2\pi i}\int_{(c)}\Gamma(d)^2\Gamma(\sigma-d)^2\,P_{\binom{[2]}1}(d\one;d,d)\,\Psi_\sigma(d)\,dd,
\]
and for $s\in\{1,2,\dots\}$, with $W_s(x)=\prod_{j=1}^{s-1}(x-j)^2$ as in Proposition~\ref{p:cosetmark},
\[
\Gamma(s)^2P_R(s\one;s,s)=\sum_{x\ge s}\big(W_s\,f_s-W_s'\,q_s\big)(x),\qquad q_s(x)=P_{\binom{[2]}1}(x\one;x,x)\,\Psi_s(x),\quad f_s=-q_s' .
\]
At the lattice points, with $k=x-s$,
\[
q_s(x)=\frac1{\Gamma(2x)}\sum_{u=-k}^{k}\binom{2k}{k+u}\int_0^\infty y^{x-1-|u|}\,\sG_{|u|;(1,-1)}(1,y)^2\,\sG_{0;(1,-1)}(1,y)\,dy\ \in\ \mathbb Q+\mathbb Q\,\Pth .
\]
\end{proposition}

\begin{proof}
$R$ is $H$ with its mark $1$ replaced by a class of two. In the opening of $H$ at $1$ the mark contributes $\Gamma(d)^2$; the class contributes the Mellin read of the product of its two Grammas, $\Gamma(d)^2P_{\binom{[2]}1}(d\one;d,d)$ (Thm.~22.13(ii), p.~188, as in the proof of Proposition~\ref{p:triangles}). $\Psi_\sigma$ is holomorphic on $\re d>0$, its one wall being the coloop, and $X_{H\setminus1}/X_{H/1}<1$ while $P_{\binom{[2]}1}$ decays like $4^{-d}$: close to the right as in Proposition~\ref{p:cosetmark}. At $d=x$ the insertion $(BB')^k=\sum_u\binom{2k}{k+u}(\alpha_{2a}/\alpha_{2b})^u$ moves the masses (Lemma~35.2(iii), p.~308), and the share side of the three-mark banana at a lean (proof of Proposition~\ref{p:leans}) has the indices $(-u,u,0)$: $\sG_{-u}\sG_u\sG_0=y^{-|u|}\sG_{|u|}^2\sG_0$.
\end{proof}

\begin{proposition}[The input's law]\label{p:r4law}\derived\ in Proposition~\ref{p:certs}\
On the line $1$,
\begin{multline*}
x^2(2x-1)^2(60x^2+137x+65)\,q_1(x)-(2x+1)(1680x^5+4316x^4+2782x^3+231x^2-223x-56)\,q_1(x+1)\\
+(2x+1)(2x+3)(2940x^4+9113x^3+7350x^2-89x-1080)\,q_1(x+2)\\
-18(x+2)(2x+1)(2x+3)(2x+5)(60x^2+17x-12)\,q_1(x+3)=0,
\end{multline*}
with the symbol $36\rho^3-49\rho^2+14\rho-1=(\rho-1)(4\rho-1)(9\rho-1)$, and the lines are tied by
\begin{multline*}
-4(2x-2s-1)^2\,q_{s+1}(x)+2(2x+1)(16x-1)\,q_{s+1}(x+1)\\
+\big(4x^2-(8s+1)x+2s(2s+1)\big)\,q_s(x)-2(2x+1)(4x-1)\,q_s(x+1)=0 .
\end{multline*}
The roots are the native value $\frac14$ of the pair's period times the native values $4,1,\frac49$ of $X_{H\setminus1}/X_{H/1}$: $4$ at the sign $(1,1,-1)$ of the three-mark banana, where $BB'=4$ and $X_{H/1}=1$; $1$ where $\alpha_3\to0$ and the pair splits; $\frac49$ at the table's gauge \derived.
\end{proposition}

\noindent Receipts \computed: found on $x\le91$ ($176$ equations, $28$ unknowns) and at $s\le4$; the lines $2$ to $5$ read off the line $1$ agree with Proposition~\ref{p:r4clos}. The native value $1$ sits on the symbol, the case of \S15.8(i): the coordinates of $q_1$ decay only like powers, $\Pth=-\lim r/c$ again, and the closing cannot be summed coordinate by coordinate. It is summed whole, with the minimality below.

\begin{proposition}[Zero shares]\label{p:r4zero}Laurent data and boundary \derived; second relation \computed\
At $d=0$,
\[
P_{\binom{[2]}1}(d\one;d,d)=\frac2d-\frac{\Gamma(\frac12)^4}{3}\,d+O(d^3),\qquad \Psi_1(d)=\frac2d+2+O(d) .
\]
The pole of $\Psi_1$ is the coloop's wall: each end of $\alpha_3$ gives $\int X_{H\setminus1}^{-1}dy=P_{\binom{[2]}1}(\one;1,1)=1$, and the finite part is $\int\log X_{H\setminus1}\,X_{H\setminus1}^{-1}dy=2$ on the chart of the pair. So $q_1(d)=4/d^2+4/d+O(1)$ where $c_0$ has a double zero: the law at $x=0$ returns the $4$, and its derivative is the boundary
\[
56\,f_1(1)-3240\,f_1(2)+6480\,f_1(3)=-\frac{1726}3-\frac{16}{15}\,\Pth .
\]
The next relation carries $\Gamma(\frac12)^4$, the class's term of order $d$:
\[
3\,f_1(1)-27\,f_1(2)=11+9\,\Pth-2\,\Gamma(\tfrac12)^4\qquad(\text{PSLQ},\ 130\text{ digits}).
\]
Hence $f_1(x)-f_1(1)\,c_1(x)\in\mathbb Q+\mathbb Q\,\Pth+\mathbb Q\,\Gamma(\frac12)^4$ for every $x$ and, by the lines, on every line: as for $\binom{[5]}2$ the one free slope lies along the three-mark coordinate. Minimality fixes it: no component along the native value $1$.
\end{proposition}

\noindent The free slope is
\begin{gather*}
f_1(1)=2\,\Pth-\partial_\sigma P_{\binom{[3]}1}(\sigma\one;\sigma,\sigma)\big|_{\sigma=1}-\int\log X_{H\setminus1}\,\frac{dy}{X_{H/1}},\\
f_1(1)=7.45882387041008265101442831798585979305\ldots,\\
\int\log X_{H\setminus1}\,\frac{dy}{X_{H/1}}=5.3041428705004791404583475012260708\ldots,
\end{gather*}
by two-dimensional quadrature $5.3041428705004785$ \computed.

\begin{theorem}[The masters of two marks doubled]\label{t:r4masters}\computed\
\begin{enumerate}[label=(\roman*)]
\item $P_R(\one;1,1)-8\,P_R(2\one;2,2)=6\,\Pth-2-\frac49\,\Gamma(\frac12)^4$.
\item $521\,P_R(\one;1,1)-12288\,P_R(2\one;2,2)+28672\,P_R(3\one;3,3)=648+1350\,\Pth$.
\item Through the closing, the corner relation of \S15.4 is the minimality alone, with the value $4432269600=-[\sigma^4]c_0\cdot44$; the second corner is $-\frac{1741824}{13}$ times the boundary of Proposition~\ref{p:r4zero} plus the minimality, with the value $-80640\cdot2503$. The second corner is the boundary of the input's law at zero shares, where the class and the coloop put the double pole $4/d^2+4/d$.
\item The lowered settings $x_s=\Gamma(s)^{-2}\sum_xW_s(x)q_s(x)$ obey the ray law with no corner; at $\sigma=-1$ they have the rational second corner $80640\cdot432$, and $239-148\,x_1+5244\,x_2-14336\,x_3=0$.
\end{enumerate}
With the corner and the second corner of \S15.4, every $P_R(s\one;s,s)$ lies in $\mathbb Q+\mathbb Q\,\Pth+\mathbb Q\,\Gamma(\frac12)^4+\mathbb Q\,P_R(\one;1,1)$: one master over $\mathbb Q+\mathbb Q\,\Pth+\mathbb Q\,\Gamma(\frac12)^4$, and the $3^\dagger$ of \S15.4 are $\Pth$, $\Gamma(\frac12)^4$ and $P_R(\one;1,1)$.
\end{theorem}

\noindent Receipts \computed. The three initial slopes solved from the minimality at $x=640$ and $P_R(1),P_R(2)$ give $P_R(3),P_R(4),P_R(5)$ to $155$ digits against Procedure~\ref{p:proc}; the slope at $x=548$ is then $1.7\cdot10^{-336}$, of the size $4^{-x}$. The values:
\begin{gather*}
P_R(\one;1,1)=7.8658373683640215924219746974461682\ldots,\\
x_1=2.5094223834929320566277011203297088\ldots .
\end{gather*}
There is no rational relation among $1$, $\Pth$, $\Gamma(\frac12)^4$, $f_1(1)$ and $P_R(\one;1,1)$, nor with $\partial_\sigma P_{\binom{[3]}1}$ and $\Pfour$ adjoined, below $10^{10}$.

\begin{remark}[What the closing found]\label{r:closingfound}
(i) In both lists the lattice values are three-mark leans, and the slope input has one free direction, the three-mark coordinate of the values; the other initial slopes are fixed by the boundaries of the law, at the cut and the wall for $\binom{[5]}2$, at zero shares for the triangle, and the free one by minimality along the leading native value: $\frac12$ at the complex gauge $(1,1,i,i)$ for $\binom{[5]}2$, $1$ at the sign of the three-mark banana for the triangle. The words of the masters, $V_s$ with $s\le3$ and $W_s$ with $s\le2$, have degree below what the adjoint of the law adds (six for $\binom{[5]}2$; five for the triangle, whose symbol vanishes at $\rho=1$), so summation by parts in the sense of Proposition~\ref{p:sbp} has nothing to reduce; the content is in the boundaries.

(ii) The corners of the ray laws are those boundaries. For $\binom{[5]}2$ the corner and the law itself are the boundary at the cut; for the triangle the corner is the minimality and the second corner is the boundary at zero shares.

(iii) Over $\mathbb Q$ the masters are $\{\Pth,P(1),P(2)\}$ for $\binom{[5]}2$ and $\{\Pth,\Gamma(\frac12)^4,P_R(1)\}$ for the triangle with two marks doubled. The three-mark centre is a master of both; in the triangle it enters through the contraction $H/1$.

(iv) The weight $C_1$ of the three-mark centre in the lowered setting of $\binom{[5]}2$ is a stacked Gamma at twentieths, and the same number weighs the free slope in the period.

(v) Aliases, last: $\Pth=3L(\chi_{-3},2)$, $\Pfour=7\zeta(3)$, $\Gamma(\frac12)^4=\pi^2$; $C_1$ is the Chowla--Selberg period of $\mathbb Q(\sqrt{-5})$; $X_{\setminus}-X_{/}=(e_2^2-e_1e_3)/e_4$, so $x_1=\int dy/(X_{\setminus}-X_{/})$ is a period of the quartic surface $e_2^2=e_1e_3$; $\sG_{\nu;(1,-1)}(1,y)=2y^{\nu/2}K_\nu(2\sqrt y)$ \recalled.
\end{remark}

\begin{remark}[The leaf and the closing]\label{r:leaf}
In each of the four rays met so far, the three-mark ray, the four-mark ray, the input of $\binom{[5]}2$ along its lines and the input of the triangle with two marks doubled, the lattice values are $r+c\,\ell$ and the lattice slopes are $\rho+\lambda\,\ell+c\,\ell'$: $r,c,\rho,\lambda$ rational, the same $c$ in both, and $\ell$ the centre of the smallest piece, $\ell=\Pth$ for the three-mark ray and the two inputs and $\ell=\Pfour$ for the four-mark ray. For the two banana rays this is \derived: the cut relation and its derivative (the corner and its vanishing term of order $\sigma$, Proposition~\ref{p:corner}) have the same homogeneous part, so the one new slope $\ell'=\partial_\sigma P(\sigma\one;\sigma,\sigma)|_{\sigma=1}$ follows $c$ through the recursion (exactly, for every lattice point checked). For the two inputs it is Propositions~\ref{p:u25free} and~\ref{p:r4zero}.

In every case the coordinates $r$ and $c$ carry the leading native value of the law, and that native value sits at a gauge off the positive chart: the sign $(1,1,-1)$ of three marks (value $1$), the sign $(1,1,1,-1)$ of four marks (value $\frac14$), the complex gauge $(1,1,i,i)$ of $X_{/}/X_{\setminus}$ (value $\frac12$), the sign of the three-mark banana inside the triangle (value $1$). The leaf's centre is the multiplier that removes it,
\[
\ell=-\lim\frac rc\qquad\computed,
\]
and the true values decay like the extreme the positive chart reaches: $\frac19$ and $\frac1{16}$ at the gauges of the bananas, $\frac49$ at the table's gauge for $\binom{[5]}2$, and $\frac14$ for the triangle, from the stratum where $\alpha_3\to0$ and the pair splits.

So the closing pairs two kinds of number: Trailmoments of rational coordinate sequences against the closing's words, and the leaf's centre with its one slope. For $\binom{[5]}2$ this is Theorem~\ref{t:u25masters}(i). The stacked Gammas sit differently in the two prime lists. For the banana the whole lowered setting $x^{(1)}(1)=R+C\,\Pfour$ is the stacked Gamma at fifteenths (Proposition~\ref{p:stackedhalf}), while the weight $C=1.0705107798578621429652109212947808\ldots$ of the four-mark centre is not: there is no relation with that stacked Gamma, with $\Gamma(\frac12)^8$ over it, with $1$ and $\Pfour$, times $1,\sqrt3,\sqrt5,\sqrt{15}$, below $10^{12}$. For $\binom{[5]}2$ the weight $C_1$ of the three-mark centre is the stacked Gamma at twentieths and $x_1$ is not \computed.
\end{remark}

\subsection*{17.5\quad The laws from the relations along a ray}

The laws of \S\S17.1--17.4 were read off exact lattice values. They follow from the two relations along a ray (Thm.~35.3, p.~308), integration by parts in every mark and the shift by the stride. The relation space has finite codimension (Thm.~35.19, p.~313), so every insertion has a law that follows from it (Cor.~35.6, p.~308). Finding it is linear algebra, and what the linear algebra returns is an insertion that certifies the law.

\begin{proposition}[Certificates]\label{p:certs}\derived; the certificates \computed\
On the line $s$ of $\binom{[5]}2$ put $\langle g\rangle_m=\int g\,t^m\,X_{\setminus}^{-s}\,dy$ with the stride $t=X_{/}/X_{\setminus}=e_1e_3/e_2^2$, so that $\langle t^i\rangle_m=\Phi_s(-m-i)$ and $\langle t^iX_{\setminus}^{-1}\rangle_m=\Phi_{s+1}(-m-i)$. For each law and each line relation $\sum_jp_j(m)\langle g_j\rangle_m=0$ used in \S\S17.1--17.3 (the laws at $s\le3$; the first-order relations and the explicit relations of the lines at $s\le4$) there is an insertion
\[
H=\frac{P}{e_2^{\,d}},
\]
$P$ symmetric in three marks, of degree $2d$, with every exponent at most $d$ and coefficients in $\mathbb Q[m]$,
with $H_e$ its image under the exchange of the first mark with $e$, such that
\[
\sum_jp_j(m)\,g_j=\sum_{e\in\cE}\Big(\alpha_e\partial_{\alpha_e}H_e+H_e\big(m\,\alpha_e\partial_{\alpha_e}\log t-s\,\alpha_e\partial_{\alpha_e}\log X_{\setminus}\big)\Big).
\]
For the triangle with two marks doubled the same holds on the chart $\alpha_{2a}=1$ of $\{2a,2b,3\}$ for the laws of $\Psi_s$ (the law of Proposition~\ref{p:r4law} divided by the shift of $P_{\binom{[2]}1}$), with the stride $X_{H\setminus1}/X_{H/1}$, the origin form $X_{H\setminus1}^{-s}dy$, and $H_e=P_e/(L_1^{d_1}L_2^{d_2}L_4^{d_4})$, where $L_1,L_2$ are the letters of the three-mark banana, $L_4=\alpha_{2a}+\alpha_{2b}$, and $P_e$ lies inside the Newton polygon of the denominator.
\begin{center}
\small
\begin{tabular}{@{}l l l l@{}}
\toprule
& law & first-order relation & explicit relation\\ \midrule
$\binom{[5]}2$ & $d=7$, degree $5$ ($s=1$); $d=9$, degree $7$ ($s=2,3$) & $d=3$, degree $5$ & $d=5$, degree $5$\\
triangle & $(d_1,d_2,d_4)=(3,3,2)$, degree $6$ & $(2,2,1)$, degree $5$ & $(2,2,1)$, degree $5$\\
\bottomrule
\end{tabular}
\end{center}
The degrees are in the line coordinate.
\end{proposition}

\begin{proof}
The identity integrates to the law by Thm.~35.3, once $H_e\,t^m X_{\setminus}^{-s}$ tends to zero at the boundary of the chart. The hypothesis of Thm.~35.3 fails here. Along the line the contraction's exponent goes down as the deletion's goes up, and the direction has $\kappa_{1,S}=(d+d^\ast)_{\binom{[4]}1}(S)-(d+d^\ast)_{\binom{[4]}2}(S)$, which is $0$ for $|S|$ odd and $1$ for $|S|=2$; on the odd walls the stride does not decay but tends to the ratio of the face. The origin form carries the decay instead. Along a direction $w$ with $w_1\ge w_2\ge w_3\ge w_4$,
\begin{gather*}
\operatorname{ord}_wt=w_3-w_2\le0,\qquad \operatorname{ord}_wX_{\setminus}^{-s}=s(w_3+w_4-w_1-w_2)<0,\\
\operatorname{ord}_wH_e\le d(w_1+w_2)-d(w_1+w_2)=0,
\end{gather*}
the last because no exponent of $P$ exceeds $d$. So the boundary terms vanish for $m\ge0$. Both sides of the law are holomorphic in $m$ where the values converge ($\re m>-2$ on the line $1$), so the law holds identically in the split. For the triangle, $\operatorname{ord}_wH_e\le0$ by the Newton polygon, and $(X_{H\setminus1}/X_{H/1})^xX_{H\setminus1}^{-s}$ decays in every direction for $x>0$.
\end{proof}

\noindent Receipts \computed. Each certificate comes from exact linear algebra in under a second; the symmetric ansatz brings the law of the line $1$ of $\binom{[5]}2$ down to $444$ unknowns, and every equation is checked exactly. The certificates of the two line-$1$ laws were checked again by direct evaluation at random rational points, with the line coordinate negative and non-integral. The two banana rays, whose strides decay in every direction so that Thm.~35.3 applies as stated, have certificates of the same kind, with the denominators $e_1^2e_2$ (three marks) and $e_1^2e_3^2$ (four marks).

\begin{corollary}[What is now derived]\label{c:certs}\derived\
The laws of Propositions~\ref{p:u25law} and~\ref{p:r4law}, the line relations of Propositions~\ref{p:u25mix} and~\ref{p:r4law} for $s\le4$, the slope recurrences (their derivatives in the split), and the boundaries at the cut (Proposition~\ref{p:u25cut}) and at zero shares (Proposition~\ref{p:r4zero}, with its Laurent data) are derived. Still computed: the relations among the initial slopes found by PSLQ (Propositions~\ref{p:u25wall} and~\ref{p:u25free}, and the relation carrying $\Gamma(\frac12)^4$ in Proposition~\ref{p:r4zero}), the identities of the closings in Theorems~\ref{t:u25masters}(iii) and~\ref{t:r4masters}(iii), and the relations among the masters.
\end{corollary}

\begin{remark}[The centre read off the law]\label{r:paper}
With the laws derived, the exchange rate of Remark~\ref{r:leaf} is an exact identity at every lattice point, not only in the limit. On the three-mark centre ray $p(\sigma)=P_{\binom{[3]}1}(\sigma\one;\sigma,\sigma)$ the three-mark law (its certificate is among the receipts of Proposition~\ref{p:certs}) and the first two values,
\begin{gather*}
\sigma^2p(\sigma)-(10\sigma^2+10\sigma+3)\,p(\sigma+1)+9(\sigma+1)^2\,p(\sigma+2)=0,\\
p(1)=\Pth,\qquad p(2)=-\tfrac23+\tfrac13\,\Pth,
\end{gather*}
give $p(q)=r(q)+c(q)\,\Pth$ with $r,c$ rational, and $\Pth=-r(q)/c(q)+p(q)/c(q)$ for every $q$. The ratios $-r/c$ are
\[
0,\ \ 2,\ \ \tfrac{23}{10},\ \ \tfrac{145}{62},\ \ \tfrac{1331}{568},\ \ \tfrac{242353}{103400},\ \ \tfrac{5495507}{2344600},\ \ \tfrac{46422941}{19805800},\ \ \dots,
\]
all below $\Pth$, the error shrinking by $0.1112$ per step at $q=30$, towards $\frac19$, the native value of the gauge over that of the sign $(1,1,-1)$ \computed: about a digit of the three-mark centre per step of exact arithmetic. On the line $1$ of $\binom{[5]}2$, $\Phi_1(-m)=\int t^m X_{\setminus}^{-1}\,dy>0$ and $\Phi_1(0)=\Pth$, so $\Pth=-r_1(m)/c_1(m)+\Phi_1(-m)/c_1(m)$ \derived; the exact $c_1(m)$ is positive for $m\le80$, so these ratios are lower bounds as well, with error $6.7\cdot10^{-4}$ at $m=80$, shrinking per step towards $\frac89=\frac49\big/\frac12$ \computed. With $V_1=1$, Theorem~\ref{t:u25masters}(ii) is a series of the rationals that the law of Proposition~\ref{p:u25law} produces:
\[
\sum_{m\ge0}c_1(m)=1+\frac29+\frac{28}{405}+\frac{208}{8505}+\frac{304}{32805}+\frac{16672}{4546773}+\cdots=\frac{\sqrt3}{80}\Gst44{\tfrac1{20},\ \tfrac3{20},\ \tfrac7{20},\ \tfrac9{20}}{\tfrac12,\ \tfrac12,\ \tfrac12,\ \tfrac12}.
\]
\end{remark}

\subsection*{17.6\quad Corrections to \S\S15.4 and 16}

\begin{enumerate}[label=(\roman*)]
\item The masters column of \S15.4 counts over $\mathbb Q$. Over $\mathbb Q+\mathbb Q\,\Pth$ the list $\binom{[5]}2$ has two masters, and over $\mathbb Q+\mathbb Q\,\Pth+\mathbb Q\,\Gamma(\frac12)^4$ the triangle with two marks doubled has one (Theorems~\ref{t:u25masters} and~\ref{t:r4masters}).
\item In the remark ``What the tree decided'' of \S16.3 the three-mark centre enters every row of \S15.4 but the banana, row four included, through the contraction of $H$ at its single mark.
\item In the remark ``What the prime node reads'' of \S16.3, two of the masters of $\binom{[5]}2$ are slope Trailmoments of the deletion's ray; the third is $\Pth$.
\end{enumerate}

\subsection*{17.7\quad Open}

\begin{enumerate}[label=(\roman*)]
\item Done for the lines used here (Proposition~\ref{p:certs}). One certificate with $s$ a variable, as Theorem~\ref{t:cert} has for the banana. \notyet
\item The finite part $\frac{21}2+\frac34\Pfour$ at the wall $d=2$ from the collapse of Cor.~8.8 as a regularized lean of the four-mark banana; the third relation of Proposition~\ref{p:u25free} from the next order. \notyet
\item The relation of Proposition~\ref{p:r4zero} that carries $\Gamma(\frac12)^4$, from the next Laurent order at zero shares. \notyet
\item Theorem~\ref{t:r4masters}(i) and (ii) as boundaries: (i) ties the lines $1$ and $2$ with the coefficient $8=|\cB|$. Which insertion reads it (a mean share as in Proposition~\ref{p:block}?). \notyet
\item The other weights $A_s$, $B_s$, $R_s$, and $C_s$ for $s\ge2$, of $\binom{[5]}2$, as periods of the surface $e_2^2=e_1e_3$ at its $\mathbb Q(\sqrt{-5})$ point. \notyet
\item $\mathrm M(K_4)$, opened at a mark into $K_4\setminus e$ and $K_4/e$, the latter the triangle with two marks doubled closed here. \notyet
\end{enumerate}


\renewcommand{\thesection}{18}
\setcounter{proposition}{0}
\renewcommand{\theproposition}{18.\arabic{proposition}}

\section{The reading of a ray through the origin}

Remark~\ref{r:paper} read the three-mark centre off its law, and Remark~\ref{r:leaf} found the same exchange rate on four rays. This entry states that reading for every ray through the origin of every family, as derived statements, and runs it where \S\S15--17 had not looked: off the centre, at a lean of three marks, and on the one connected list on four marks that had not been read. The receipts are in the script ray\_reading\_s18.py (python-flint, exact arithmetic and ball arithmetic, no quadrature in the procedure).

Throughout, a \emph{ray through the origin} is a ray of Def.~35.1, p.~307, with origin $0$ and an integral direction $(\mvec_1;\mathbf s_1,\mathbf t_1)$ on the balance wall, $s_{1j},t_{1j}\in\mathbb N$, all of whose proper walls are positive, $\kappa_{1,S}>0$ for every nonempty proper $S$. Its values and its stride are
\[
P(k)=P_F(k\mvec_1;k\mathbf s_1,k\mathbf t_1)=\int_{y_{e_0}=0}f^k\,dy\quad(k\ge1),\qquad f=\alpha^{\mvec_1}\prod_jK_j^{-s_{1j}}\Psi_j^{-t_{1j}},
\]
with $\alpha^\star$ the native gauge of the direction on the orthant, where $f$ takes its largest value $f(\alpha^\star)$ there (Thm.~35.10(i), p.~309). The centre ray of Def.~35.9, p.~309, is the case $\mvec_1=\sigma\one$, $\mathbf s_1=\mathbf t_1=\boldsymbol\sigma$.

\subsection*{18.1\quad One pencil, two chains}

\begin{proposition}[The closings of a direction]\label{p:dirclose}\derived\
Put $x_j=[\alpha^0]\,(1/f)^j$ for $j\in\mathbb N$, the constant terms of the powers of the Laurent polynomial $1/f=\alpha^{-\mvec_1}\prod_jK_j^{s_{1j}}\Psi_j^{t_{1j}}$, and
\[
\mathrm{Cor}=\sum_\pi\prod_{i=2}^n\kappa_{1,U_i}^{-1},
\]
the sum over the orders $\pi$ of $\cE$, with tails $U_i=\{\pi(i),\dots,\pi(n)\}$. Then $P$ continues meromorphically in $k$, with its poles at negative rational $k$, and
\[
\lim_{k\to-j}(k+j)^{n-1}P(k)=\mathrm{Cor}\cdot x_j\qquad(j\in\mathbb N).
\]
At the centre this is Proposition~\ref{p:negcorner}: there $1/f=\prod_j(K_j\Psi_j/\alpha^{\cE})^{\sigma_j}$ and $\kappa_{1,S}=\sum_j\sigma_jc_j(S)$.
\end{proposition}

\begin{proof}
In the sector coordinates of Thm.~8.7(ii), p.~55, $\alpha_{\pi(1)}=\rho$, $\alpha_{\pi(i)}=\rho\,q_2\cdots q_i$, one has $K_j=\rho^{r_j}\prod_iq_i^{d_j(U_i)}k_{j,\pi}$, $\Psi_j=\rho^{n-r_j}\prod_iq_i^{d^\ast_j(U_i)}\psi_{j,\pi}$ and $\alpha^{\mvec_1}=\rho^{m_{1,\cE}}\prod_iq_i^{m_1(U_i)}$, so by the balance of the direction
\[
f=\prod_{i=2}^nq_i^{\kappa_{1,U_i}}\,\varphi_\pi(q),\qquad \varphi_\pi=\prod_jk_{j,\pi}^{-s_{1j}}\psi_{j,\pi}^{-t_{1j}},
\]
and by Thm.~8.10, p.~56, $P(k)=\sum_\pi\int_{[0,1]^{n-1}}\prod_iq_i^{k\kappa_{1,U_i}-1}\varphi_\pi(q)^k\,dq$. The sector letters have nonnegative coefficients and positive constant terms (Lemma~8.5, p.~55), so $(1/\varphi_\pi)^j=\sum_vc_{\pi,v}q^v$ is a polynomial. At $k=-j+\varepsilon$ write $\varphi_\pi^{k}=(1/\varphi_\pi)^j\varphi_\pi^{\varepsilon}$. By Lemma~8.6, p.~55, the term $q^v$ has a pole in $\varepsilon$ of order at most the number of $i$ with $v_i=j\kappa_{1,U_i}$, so the pole of order $n-1$ comes from $v=j\kappa_1$ alone, with the leading coefficient $c_{\pi,j\kappa_1}\prod_i\kappa_{1,U_i}^{-1}$. Finally $c_{\pi,j\kappa_1}=[q^0]\big(\prod_iq_i^{-\kappa_{1,U_i}}/\varphi_\pi\big)^j$ is the constant term of $(1/f)^j$ in the sector coordinates, and $\alpha^v\mapsto\prod_iq_i^{v(U_i)}$ is injective on exponents of degree zero, so it is $x_j$ for every $\pi$.
\end{proof}

\begin{corollary}[The law and its closings]\label{c:dirlaw}\derived\
If $\sum_{i=0}^hc_i(k)P(k+i)=0$ identically, then $\sum_{i=0}^{\min(h,j)}c_i(-j)\,x_{j-i}=0$ for every $j\in\mathbb N$, and $c_0$ vanishes to order at least $n-1$ at $0$, with
\[
[k^{n-1}]c_0\cdot\mathrm{Cor}+\sum_{i\ge1}c_i(0)\,P(i)=0 .
\]
\end{corollary}

\begin{proof}
As Corollary~\ref{c:law}, with Proposition~\ref{p:dirclose} in place of Proposition~\ref{p:negcorner}; the second statement is the case $j=0$, $x_0=1$.
\end{proof}

So a law of the values is the recurrence of the integers $x_j$ read backwards, and reading it is exact linear algebra on integers. The converse, that the least recurrence of the $x_j$, reflected, is a law of $P$, is what a certificate of Thm.~35.3, p.~308, settles, as in Proposition~\ref{p:certs}.

\begin{remark}[The torus and the chart]\label{r:twochains}\derived\
The two sequences are one integrand on two chains. For $|q|<f(\alpha^\star)^{-1}$, $\sum_{k\ge1}P(k)q^{k-1}=\int_{y_{e_0}=0}f/(1-qf)\,dy$ (Thm.~35.10(ii) with $k_0=1$). On the torus $T=\{|\alpha_e|=1\}$, where $|1/f|\le(1/f)(\one)$ because $1/f$ has nonnegative coefficients, expanding in $1/(qf)$ gives for $|q|>f(\one)^{-1}$
\[
\sum_{j\ge0}x_j\,q^{-j-1}=-\frac1{(2\pi i)^{n-1}}\oint_{T}\frac{f}{1-qf}\prod_{e\ne e_0}\frac{d\alpha_e}{\alpha_e}.
\]
The pencil $\prod_jK_j^{s_{1j}}\Psi_j^{t_{1j}}-q\,\alpha^{\mvec_1}$ of Thm.~35.10 carries both: the positive chart gives its solution at $q=0$, whose coefficients are the periods, and the torus its solution at $q=\infty$, whose coefficients are integers for a list. Corollary~\ref{c:dirlaw} is the one equation, inhomogeneous by its boundary terms, read at its two ends, and $\mathrm{Cor}$ is the ratio of the two normalizations.
\end{remark}

\subsection*{18.2\quad The slit}

\begin{proposition}[The ray lives on the slit plane]\label{p:slit}\derived\
$G(q)=\int_{y_{e_0}=0}f/(1-qf)\,dy$ is holomorphic on $\mathbb C\setminus[f(\alpha^\star)^{-1},\infty)$, and it continues $\sum_{k\ge1}P(k)q^{k-1}$.
\end{proposition}

\begin{proof}
On the chart $0<f\le f(\alpha^\star)$ (Thm.~35.10(i)). For $q$ off the slit the segment $q\cdot(0,f(\alpha^\star)]$ misses $1$, so $|1-qf|$ is bounded below by the distance from $1$ to that segment, locally uniformly in $q$, and $\int f\,dy=P(1)<\infty$; dominated convergence gives holomorphy. On $|q|<f(\alpha^\star)^{-1}$ the integral is the series.
\end{proof}

\begin{theorem}[Positive selection]\label{t:select}(i)--(iii) \derived; the count in (iii) \computed\ on every ray below.\
Let $\sum_{i=0}^hc_i(k)P(k+i)=0$ be a law of a ray through the origin, $c_0c_h\ne0$, let $D$ be the largest degree of the $c_i$, and $\chi(\lambda)=\sum_i[k^D]c_i\,\lambda^i$ its symbol, with roots $\lambda_c$; it vanishes at the native value of every nondegenerate native gauge (Cor.~35.6, p.~308), and a root need not be a native value.
\begin{enumerate}[label=(\roman*)]
\item For every solution $(x_k)_{k\ge1}$ of the law, $X(q)=\sum_kx_kq^{k-1}$ satisfies a linear differential equation in $q$ with polynomial coefficients, inhomogeneous by a polynomial that depends linearly on $x_1,\dots,x_h$; in $\mathbb C^\times$ its singular points are the $q_c=1/\lambda_c$. So $X$ continues along every path in $\mathbb C^\times$ that avoids them.
\item The values are the solution whose $X$ is holomorphic on $\mathbb C\setminus[f(\alpha^\star)^{-1},\infty)$ (Proposition~\ref{p:slit}): holomorphic at every $q_c$ off the slit, that is at every root $\lambda_c\notin(0,f(\alpha^\star)]$.
\item At a singular point $q_c$, holomorphy of $X$ is $\nu_c$ linear conditions on $(x_1,\dots,x_h)$, $\nu_c$ the number of independent local solutions of the equation that are not holomorphic there, at most the multiplicity of $\lambda_c$ as a root of $\chi$. At a simple root that is the native value of a nondegenerate native gauge it is one: the branch of Thm.~35.11, p.~310, read at that gauge in place of $\alpha^\star$, of exponent $\frac{n-3}2$, with a logarithm for odd $n$. With the corner of Corollary~\ref{c:dirlaw}, if $1+\sum_{q_c\text{ off the slit}}\nu_c=h$ and the conditions are independent, the law, the corner and positivity determine the values.
\item Then, for a basis $x^{(0)},\dots,x^{(h-1)}$ of rational solutions fixed by the corner, $P=x^{(0)}+\sum_{i\ge1}M_ix^{(i)}$, and the $M_i$ are the limits as $K\to\infty$ of the solutions of $x^{(0)}_{K+l}+\sum_iM_ix^{(i)}_{K+l}=0$, $0\le l<h-1$, with error $O\big((f(\alpha^\star)/\min_{c\ne\star}|\lambda_c|)^K\big)$ up to powers of $K$.
\end{enumerate}
\end{theorem}

\begin{proof}
(i) Multiply the law by $q^{k-1}$ and sum over $k\ge1$: $k$ acts as $q\partial_q+1$, a shift by $i$ as $q^{-i}$ after removing the first $i$ terms, which leaves a polynomial. The coefficient of the top power of $q\partial_q$ is $\sum_i[k^D]c_i\,q^{-i}$, which vanishes in $\mathbb C^\times$ exactly where $\chi(1/q)=0$. (ii) Proposition~\ref{p:slit}. (iii) The inhomogeneity is holomorphic at $q_c$, so the singular part of $X$ there lies in the span of the local homogeneous solutions that are not holomorphic, with coefficients linear in $(x_1,\dots,x_h)$. At a nondegenerate native gauge the value over a steepest-descent disc (Thm.~35.5, p.~308) gives one branch of exponent $\frac{n-1}2-1$; Thm.~35.11 is this computation at $\alpha^\star$, and it is local. (iv) The coefficients of $X$ grow at the rate $|\lambda_c|$ of the singular point nearest $0$ at which $X$ is singular. The rational solutions are singular at the $q_c$ off the slit and $P$ is not, so the combination annihilating their $h-1$ largest rates is $P$, and its own size $f(\alpha^\star)^K$ against the smallest of the others is the error.
\end{proof}

\begin{procedure}[The reading of a ray]\label{p:reading}
(1) From the family and the direction write $1/f$ and the walls $\kappa_{1,S}$, and $\mathrm{Cor}$ (Proposition~\ref{p:dirclose}). (2) The closings $x_j=[\alpha^0](1/f)^j$, exactly. (3) Their least recurrence by exact linear algebra, reflected (Corollary~\ref{c:dirlaw}): the law; its certificate by Thm.~35.3 as in Proposition~\ref{p:certs}. (4) The native values, the roots of the symbol, and the native gauges, the critical points of $\log f$: the signs and what a lean moves them to, and families. (5) The seeds: the corner, and a second corner wherever $c_0$ and $c_1$ share a zero (\S15.4). (6) Positivity, Theorem~\ref{t:select}(iii)--(iv): the masters as limits of rational solutions, at the rate $f(\alpha^\star)/\min_{c\ne\star}|\lambda_c|$. (7) Integer relations among the masters and against the members already in use; what is left over $\mathbb Q$ are the masters of the ray.
\end{procedure}

\begin{remark}[What sits on the slit]\label{r:onslit}
Positivity says nothing at a native value in $(0,f(\alpha^\star))$. For a ray strictly inside its walls $f$ has one critical point on the orthant (Thm.~8.18(ii)--(iii), p.~58), so such a value would sit at a native gauge off the orthant; on every ray strictly inside its walls read in \S\S15--18 there is none. On a ray along walls the slit carries the faces: on the line $1$ of $\binom{[5]}2$ (\S17.1) the native values $\frac13$ and $\frac14$ of Proposition~\ref{p:u25roots} lie below $\frac49$, they come from strata at the boundary of the chart, and the constants there are the faces' own, the three-mark centre among them (Theorem~\ref{t:u25masters}). That is Thm.~35.15, p.~312, read on the slit.
\end{remark}

\subsection*{18.3\quad Off the centre: a lean of three marks}

The centre is the fixed point of Cremona (Thm.~8.13, p.~57), where both letters take one exponent; a lean gives them two. For $\binom{\cE}1$, $K=\sum_e\alpha_e$ and $\Psi=\alpha^{\cE}\sum_e\alpha_e^{-1}$, and at uniform masses $m\one$ the balance is $nm=s+(n-1)t$.

\begin{proposition}[The lean is the index of the Gramma]\label{p:leanindex}\derived; receipt \computed\
For masses $\mvec$ on the balance wall and $\re s,\re t>0$,
\[
\Gamma(s)\Gamma(t)\,P_{\binom{\cE}1}(\mvec;s,t)=\int_0^\infty y^{t-1}\prod_{e\in\cE}\sG_{m_e-t;(1,-1)}(1,y)\,dy ,\qquad \sum_e(m_e-t)=s-t .
\]
At the centre this is Proposition~\ref{p:bananapin}; the lean and the offsets move the index slots of the marks' Grammas. In the direction $(2\one;4,1)$ of three marks, $P_{\binom{[3]}1}(2\one;4,1)=\frac16\int_0^\infty\sG_{1;(1,-1)}(1,y)^3\,dy=0.030080945282262304399884274898349847\ldots$ in ball arithmetic.
\end{proposition}

\begin{proof}
On the share side (Proposition~\ref{p:share}) the mark $e$ carries the share $a_e$ of the basis $\{e\}$ in $K$ and the share $b_e$ of the cobasis $\cE\setminus e$ in $\Psi$, and the pin at $e$ reads $a_e+\sum_{e'\ne e}b_{e'}=m_e$, that is $a_e=b_e+m_e-t$. So each mark carries $\Gamma(b_e+m_e-t)\Gamma(b_e)$, the share kernel of $\sG_{m_e-t;(1,-1)}(1,y)$ (Prop.~4.2, p.~25), and the Mellin read of the product pins $\sum_eb_e=t$, as in Proposition~\ref{p:bananapin}.
\end{proof}

\begin{proposition}[The native gauges of a lean of $\binom{\cE}1$]\label{p:leangauges}\derived\
In the direction $(m\one;s,t)$, $s\ne t$, put $u_e=\alpha_e/K(\alpha)$, $v_e=\alpha_e^{-1}/K(\alpha^{-1})$ and $\lambda=\big(K(\alpha)K(\alpha^{-1})\big)^{-1}$. At a native gauge $su_e-tv_e=\frac{s-t}n$ and $u_ev_e=\lambda$ for every $e$, so every $u_e$ is a root of $su^2-\frac{s-t}nu-t\lambda=0$; with $j$ marks at one root $u_-$ and $n-j$ at the other $u_+$,
\[
(n-2j)\,u_+=1-j\,\frac{s-t}{ns},\qquad f=\Big(\prod_eu_e\Big)^{\frac{s-t}n}\lambda^t .
\]
At $s=t$ the two roots are $\pm u_+$ and this is Proposition~\ref{p:onlysigns}. For three marks in the direction $(2\one;4,1)$: $j=0$ is the table's gauge, $f=\frac1{243}$; $j=1$ is $\alpha\propto(3,3,-2)$ and its images, $f=-\frac{27}{64}$, the sign $(1,1,-1)$ moved along the real line.
\end{proposition}

\begin{proof}
Nativity is $m=sR_e(\alpha)+tR^\ast_e(\alpha)$ with $R_e=u_e$ and $R^\ast_e=1-v_e$ (Thm.~7.14, p.~50), and $u_ev_e=1/\big(K(\alpha)K(\alpha^{-1})\big)$ for every $e$. The $u_e$ sum to $1$, which fixes $u_+$. For the value, $\alpha^{\one}=K(\alpha)^n\prod_eu_e$ and $K(\alpha^{-1})=1/(K(\alpha)\lambda)$, and the powers of $K(\alpha)$ cancel by the balance.
\end{proof}

\noindent Receipts \computed. The walls of $(2\one;4,1)$ are $2$ on marks and $3$ on pairs, so $\mathrm{Cor}=6/(2\cdot3)=1$. The closings are $x_j=(4j)!\,j!\sum_{b_1+b_2+b_3=j}\prod_f\big(b_f!\,(b_f+j)!\big)^{-1}=1,36,4620,762300,140216076,\dots$; sixty of them have a least recurrence of order $2$ and degree $6$, and reflected it is the law
\begin{gather*}
-3k^2(3k+1)(3k+2)(265k^2+636k+383)\,P(k)\\
+(1721705k^6+7575502k^5+13047276k^4+11068046k^3+4766659k^2+965100k+75600)\,P(k+1)\\
+972(k+1)^2(4k+5)(4k+7)(265k^2+106k+12)\,P(k+2)=0,
\end{gather*}
$P(k)=P_{\binom{[3]}1}(2k\one;4k,k)$, with the symbol $265(243\lambda-1)(64\lambda+27)$: the two native values of Proposition~\ref{p:leangauges}, which were read off the list before the closings were computed. The corner gives $12600\,P(1)+68040\,P(2)=383$; positivity gives
\[
P_{\binom{[3]}1}(2\one;4,1)=0.0300809452822623043998842748983498467336738511\ldots,
\]
about two digits per step (the rate $\frac{64}{6561}$), against $0.95$ for the centre, and Proposition~\ref{p:leanindex} confirms it. Over $\mathbb Q$,
\[
54\,P_{\binom{[3]}1}(2\one;4,1)=11-4\,\Pth\qquad(380\text{ digits}).
\]
The Cremona partner $(3\one;1,4)$ has the same closings, by $\alpha\mapsto\alpha^{-1}$, and the same law. So the lean is a new coordinate system on the content of the centre: a different law, different native gauges, a different corner, the same master.

The rational coordinates of this ray do not give a gap (Thm.~34.3, p.~300): on $(1,\Pth)$ their denominators beyond $d_k^2$ grow like $e^{8.7k}$, with primes of the apparent factor $265k^2+\cdots$ among them, against the decay $243^{-k}=e^{-5.49k}$ \computed; at the centre the excess is $3^{1.88k}$, against $9^{-k}$.

\subsection*{18.4\quad The fourth list on four marks}

On four marks the connected lists are $\binom{[4]}1$, $\binom{[4]}2$, $\binom{[4]}3$ and the triangle $H$ with one mark doubled; the first two were read in \S15.1, and the third is the dual of the first (Prop.~8.14, p.~57). $H$ is the list of \S17.4, where it was a step toward the triangle with two marks doubled; its own centre was not read there. With the doubled class $\{a,b\}$ and the marks $2,3$ its bases are the pairs other than $\{a,b\}$, and with $A=\alpha_a+\alpha_b$, $B=\alpha_2+\alpha_3$
\[
K=AB+\alpha_2\alpha_3,\qquad \Psi=AB+\alpha_a\alpha_b :
\]
the letters exchange under $\{a,b\}\leftrightarrow\{2,3\}$, so $H$ is self-dual up to that relabeling. The widths are $1$ on marks, on $\{a,b\}$, on $\{2,3\}$ and on triples, and $2$ on the four other pairs (Def.~7.10, p.~49), so at the centre $\mathrm{Cor}=8\cdot1+16\cdot\frac12=16$.

\begin{proposition}[The native gauges of the centre of $H$]\label{p:Hgauges}\derived; that there are no others \computed\
The table's gauge, of native value $\frac1{25}$; the sign $(-1,-1,1,1)$, of native value $\frac19$; and, of native value $1$, the three curves
\[
\{A=0,\ \alpha_a^2=\alpha_2\alpha_3\},\qquad \{B=0,\ \alpha_2^2=\alpha_a\alpha_b\},\qquad \{A=B=0\},
\]
the second the image of the first under the self-duality; on them the class letters vanish and $K$, $\Psi$ collapse to monomials, the case of Proposition~\ref{p:families}(ii) with one class. Every sign of value $1$ lies on them.
\end{proposition}

\begin{proof}
At the centre nativity is $R_e(\alpha)=R_e(\alpha^{-1})$ (\S15.2). On $A=0$ also $A(\alpha^{-1})=A/(\alpha_a\alpha_b)=0$, so $K(\alpha)=\alpha_2\alpha_3$, $K(\alpha^{-1})=1/(\alpha_2\alpha_3)$ and $R_2=R_3=1$ at $\alpha$ and at $\alpha^{-1}$, while $R_a(\alpha)=\alpha_aB/(\alpha_2\alpha_3)$ and $R_a(\alpha^{-1})=B/\alpha_a$ agree exactly when $B=0$ or $\alpha_a^2=\alpha_2\alpha_3$, and likewise at $b$; there $K(\alpha)K(\alpha^{-1})=1$. The second curve follows by the self-duality. At a sign $K(\eta)=(\eta_a+\eta_b)(\eta_2+\eta_3)+\eta_2\eta_3\in\{\pm5,\pm3,\pm1\}$ (Proposition~\ref{p:signsparallel}), and $K(\eta)=\pm1$ forces $\eta_a=-\eta_b$ or $\eta_2=-\eta_3$, which puts the sign on one of the curves. Newton's method from several thousand starts in $(\mathbb C^\times)^3$ finds no critical point of $\log f$ off these.
\end{proof}

\noindent Receipts \computed. By Proposition~\ref{p:parallel} the closings are
\[
x_k=\sum_{i+j+l=k}\Big(\frac{k!}{i!\,j!\,l!}\Big)^2\binom{2i+2j}{i+j}=1,5,53,761,12661,229705,4410665,\dots,
\]
and the least recurrence of $160$ of them has order $4$ and degree $7$ (156 equations, 40 unknowns, one solution). Reflected, the law of the centre of $H$, $P(k)=P_H(k\one;k,k)$, is
\begin{gather*}
k^3(k+1)(144k^3+864k^2+1693k+1075)\,P(k)\\
-2(k+1)(2592k^6+19440k^5+56322k^4+80296k^3+60004k^2+23017k+3580)\,P(k+1)\\
+2(21168k^7+211680k^6+884919k^5+2003456k^4\\
\qquad+2646984k^3+2031207k^2+830703k+137260)\,P(k+2)\\
-2(k+2)^2(34848k^5+261360k^4+741842k^3+984642k^2+598948k+127215)\,P(k+3)\\
+225(k+2)^2(k+3)^2(144k^3+432k^2+397k+102)\,P(k+4)=0,
\end{gather*}
with the symbol $144(25\lambda-1)(9\lambda-1)(\lambda-1)^2$: the curves double the root $1$, an integral direction in the exceptional set of Thm.~35.19(i), p.~313. (A search by total size first returns a recurrence of order $5$ and degree $5$: this law composed with a factor of order one, whose symbol has the extra root $\frac{131003}{4771109}$.) The corner relation is $17200-7160P(1)+274520P(2)-1017720P(3)+826200P(4)=0$, and $c_0$, $c_1$ share the zero $k=-1$; the second corner (\S15.4) is rational,
\[
57\,P(1)-807\,P(2)+1050\,P(3)=196\qquad(530\text{ digits}),
\]
the constant term at $k=-1$ being $1176$. Positivity (Theorem~\ref{t:select}: three conditions, one at $\frac19$ and two at the double root) gives
\begin{gather*}
P_H(\one;1,1)=4.096236998337882850055553228351268637634670896\ldots,\\
P_H(2\one;2,2)=0.04771649869151311029741208089066366661776262936\ldots .
\end{gather*}

\begin{proposition}[The centre of $H$ on its share side]\label{p:Hshare}\derived; receipt \computed\
$H$ is a series root over the parallel cell of $\{a,b\}$ and the two marks left in place (Def.~22.17(5), p.~191), so, as in Proposition~\ref{p:triangles},
\[
\Gamma(\sigma)^2P_H(\sigma\one;\sigma,\sigma)=\frac1{2\pi i}\int_{(c)}\Gamma(d)^2\,\Gamma(\sigma-d)^2\,P_{\binom{[2]}1}(d\one;d,d)^2\,dd,\qquad P_{\binom{[2]}1}(d\one;d,d)=\Gst21{d,\ d}{2d},
\]
$0<c<\sigma$. At $\sigma=1$, on $c=\frac12$, the integral is $4.0962369983378828500555532283512686376346709\ldots$ in ball arithmetic: the value of the positive selection, to $44$ digits.
\end{proposition}

\begin{proof}
A series node multiplies the transforms of its cells (Thm.~22.13, p.~188): the two marks give $\sG_{0;(1,-1)}(1,y)^2$, whose Mellin read is $\Gamma(z)^2P_{\binom{[2]}1}(z\one;z,z)$ (Proposition~\ref{p:bananapin}), and the parallel cell gives the transform whose Mellin read is $\Gamma(w)^2P_{\binom{[2]}1}((\sigma-w)\one;\sigma-w,\sigma-w)$ (Proposition~\ref{p:triangles}); the Mellin read of the product pins the sum of the shares, and $d=\sigma-w$.
\end{proof}

The integrand is $\Gamma(d)^6\Gamma(\sigma-d)^2/\Gamma(2d)^2$. Its poles at $-\mathbb N$ have order four and residues that grow like $16^m$, and to the right the double poles of $\Gamma(\sigma-d)^2$ meet a factor that grows faster than every power: neither closing of Thm.~4.5, p.~26, converges. The centre of $H$ is a fibre integral of stacked Gammas that does not land on a lattice, as the five-mark centre does not (Remark~\ref{r:live}).

\noindent The lean. The direction $(2\one;3,1)$ has the walls $2$ on marks, $3$ on $\{a,b\}$, $1$ on $\{2,3\}$, $4$ on the other pairs and $2$ on triples, so $\mathrm{Cor}=\frac13+1+1=\frac73$. On the slice $\alpha_a=\alpha_b=x$, $\alpha_2=\alpha_3=1$ the stride is $x^3/\big((4x+1)^3(x+4)\big)$, whose critical points are the roots of $2x^2-x-6=(x-2)(2x+3)$ \derived: the native gauge of the direction on the orthant is $(2,2,1,1)$, where $K=9$, $\Psi=12$ and $f=\frac4{2187}$, and the sign has moved to $(-\frac32,-\frac32,1,1)$, $f=\frac{27}{2500}$; the curves move to curves of value $-\frac{16}{27}$ \computed. The closings $x_j=[\alpha^{2j\one}]K^{3j}\Psi^j=1,48,10240,3131184,\dots$ have a least recurrence of order $4$ and degree $25$ ($196$ equations, $104$ unknowns, one solution), with the symbol a multiple of $(2187\lambda-4)(2500\lambda-27)(27\lambda+16)^2$; positivity converges at the rate $\frac{10000}{59049}$, about $0.77$ digits per step, and the lean's own second corner is rational: $413667900\,P'(1)+6131684550\,P'(2)+12031031250\,P'(3)=24789199$, $P'(k)=P_H(2k\one;3k,k)$.

\begin{theorem}[One content, two rays]\label{t:Hcontent}\computed\
\begin{gather*}
28\,P_H(2\one;3,1)=14-3\,P_H(\one;1,1)-P_H(2\one;2,2),\\
5670000\,P_H(4\one;6,2)=126002\,P_H(\one;1,1)+143398\,P_H(2\one;2,2)-522781,\\
95681250000\,P_H(6\one;9,3)=3048368669-731189098\,P_H(\one;1,1)-1115791502\,P_H(2\one;2,2),
\end{gather*}
to $530$ digits. The centre ray and the lean ray of $H$, with two laws, two corners and two positive selections, land on the one space $\mathbb Q+\mathbb Q\,P_H(\one;1,1)+\mathbb Q\,P_H(2\one;2,2)$.
\end{theorem}

\begin{remark}[What the masters of $H$ are not]\label{r:Hnot}\computed\
There is no rational relation among $1$, $P_H(\one;1,1)$, $P_H(2\one;2,2)$ and the members in use, $\Pth$, $\Gamma(\frac12)^4$, $\Tx{3}{\IZ}{0}$, $\Gamma(\frac12)^6/\sqrt3$, $\sqrt3\,\Gamma(\frac12)^2$, $\sqrt3\,\Gamma(\frac12)^2\log3$, $\Gamma(\frac12)^4\log2$, $\Gamma(\frac12)^4\log3$, $\log2$, $\log3$, $\Pth\log2$, $\Pth\log3$, $\Gamma(\frac13)^6/\Gamma(\frac12)^4$ and $\Gamma(\frac12)^8/\Gamma(\frac13)^6$, with coefficients below $10^{24}$ ($420$ digits); nor with $\Pth$, the two masters of $\binom{[5]}1$, the stacked Gamma $X$ at fifteenths of Proposition~\ref{p:stackedhalf}, $\Gamma(\frac12)^8/X$, $\Gamma(\frac12)^4$, $\sqrt3\,\Gamma(\frac12)^2$, $\sqrt{15}$, and $\sqrt{15}$ times $X$, $\Gamma(\frac12)^8/X$ and $\Gamma(\frac12)^4$, below $10^{28}$ ($380$ digits). In that framework, over $\mathbb Q$, the two masters of $H$ are new numbers. $H$ and $\binom{[5]}1$ have the same native values $\frac1{25},\frac19,1$ at their centres and different masters: the native values do not decide the numbers.
\end{remark}

\noindent What the tree expects \notyet. Opened at the mark $2$ (Prop.~19.19, p.~144), $H$ has the children $\binom{[3]}1$ and the pair $\{a,b\}$ with the coloop $3$, the strides $X_{H/1}$ and $X_{H\setminus1}$ of \S17.4, and its input is $\Psi_\sigma$ of \S17.4 without the factor $P_{\binom{[2]}1}$. Its lattice values are three-mark leans, in $\mathbb Q+\mathbb Q\,\Pth$, and its slopes should have one free direction (Proposition~\ref{p:r4zero}). As for $\binom{[5]}2$ (Theorem~\ref{t:u25masters}), the masters over the field of the weights of that closing should be $\Pth$ and one free slope, the weights carrying what Remark~\ref{r:Hnot} sees as new.

\subsection*{18.5\quad Two observations}

(i) \emph{The exact pairing belongs to the banana.} Propositions~\ref{p:pairing} and~\ref{p:edge} give
\[
\mathrm{Cor}\cdot\det H=-\frac{\lambda^\star\chi'(\lambda^\star)}{[k^{n-1}]c_0}=\prod_{c\ne\star}\Big(1-\frac{\lambda^\star}{\lambda_c}\Big)
\]
whenever the closings solve the adjoint of the law and $c_0$ has degree $n-1$; here $H$ in $\det H$ is the Hessian of Proposition~\ref{p:edge}. The laws of $\binom{[n]}1$ for $n=2,3,4,5$ (Theorem~\ref{t:cert}, \S15.7, Proposition~\ref{p:pairing}) are self-adjoint up to sign, $c_i(-k-i)=\pm c_i(k)$ \computed, and the identity is
\[
n!\cdot\frac{2^{n-1}}{n^n}=\prod_{j=1}^{\lceil n/2\rceil-1}\Big(1-\frac{(n-2j)^2}{n^2}\Big)=\prod_{j=1}^{\lceil n/2\rceil-1}\frac{4j(n-j)}{n^2}\quad\derived,
\]
which holds for every $n$: the low end of strength, the corner (Thm.~38.5(i), p.~330), and the high end, the Gaussian at the native gauge (Thm.~38.5(ii)), tied through the native values. On $H$, whose law is not self-adjoint, and on $\binom{[5]}2$ (\S15.5) the closings solve no adjoint and the identity fails.

(ii) \emph{The content belongs to the list.} On $\binom{[3]}1$ and on $H$, two rays each, the masters of the ray are those of the list. The letter and mark relations (Thm.~8.3, p.~54) keep integer data on the balance wall with rational coefficients, so every integer period of a list lies in the span they leave, and each ray samples that span. A ray with the symmetries of the list may sample only their invariant part, which offsets would break. Hypothesis: every ray through the origin strictly inside its walls, without the list's symmetries, spans the whole content. \notyet

\subsection*{18.6\quad Open}

\begin{enumerate}[label=(\roman*)]
\item Certificates of Thm.~35.3 for the three laws of \S\S18.3--18.4, and the converse of Corollary~\ref{c:dirlaw} in general. \notyet
\item The slit: for rays strictly inside their walls, no native value in $(0,f(\alpha^\star))$; for rays along walls, exactly the faces' native values there, carrying the faces' masters. \notyet
\item The content of a list: an offset ray, and a ray breaking the symmetries $a\leftrightarrow b$, $2\leftrightarrow3$ of $H$. \notyet
\item The masters of $H$ through its hinge at a mark, as in \S\S17.1--17.3, with the weights. \notyet
\item $H$ against $\binom{[5]}1$: the local exponents at $\frac19$ and $1$, the frames (Prop.~35.8, p.~309), and whether one law is a transform of the other. \notyet
\item The period matrix of the centre ray of $H$: the solutions selected at its other native gauges, and whether they are stacked Gammas, as the lowered settings of $\binom{[5]}1$ are (Proposition~\ref{p:stackedhalf}). \notyet
\item A carrier built for the purpose: a list on six marks that is not a graph, with native gauges off the signs and off the torus, read by certificates rather than closings. \notyet
\end{enumerate}

\noindent Aliases, last: the $x_j$ are the classical periods of the Laurent polynomial $1/f$, and the law is the Picard--Fuchs operator of its pencil read on coefficients; Theorem~\ref{t:select}(iv) is the mechanism of Ap\'ery limits, the selection read off the slit (Pincherle's minimal solution, Perron's theorem); $H$ is the cycle matroid of a triangle with one edge doubled; $\Gamma(\frac12)^6/\sqrt3$ stands for $L(\chi_{-3},3)=\frac{4\pi^3}{81\sqrt3}$; $\sG_{\nu;(1,-1)}(1,y)=2y^{\nu/2}K_\nu(2\sqrt y)$, so Proposition~\ref{p:leanindex} reads the lean as the order of the Bessel functions.


\renewcommand{\thesection}{19}
\setcounter{proposition}{0}
\renewcommand{\theproposition}{19.\arabic{proposition}}

\section{The three-mark centre, set as an examination}

Remark~\ref{r:paper} read the three-mark centre off its law as a list of ratios, and Procedure~\ref{p:reading} states the reading of every ray through the origin. This entry runs the reading once, on that one ray, written as an examination question with its answer key and its marking scheme. A key has to say what a complete solution shows at every step and what each omission costs. So every hypothesis that \S\S15--18 checked once in passing is checked here, at the step that uses it, and the places where a solution can look right and be wrong are named, with the credit they earn. The ratios of Remark~\ref{r:paper} become a certified bracket. The receipts are in the script exam\_s19.py (python-flint: exact arithmetic for the solution, ball arithmetic only for the independent check of \S19.5).

\subsection*{19.1\quad The question}

\begin{center}
\fbox{\parbox{0.94\linewidth}{\smallskip
\textbf{Question} (40 points). Let $F=\binom{[3]}1$ be the list on the marks $\{1,2,3\}$ whose bases are the singletons, with unit weights, so that its letters are
\[
K=\alpha_1+\alpha_2+\alpha_3,\qquad \Psi=\alpha_1\alpha_2+\alpha_1\alpha_3+\alpha_2\alpha_3,
\]
and let $\ell=P_F(\one;1,1)$ be the period of its centre. By the reading of its centre ray (Procedure~\ref{p:reading}), that is, the law of the ray certified by the relations along a ray (Thm.~35.3, p.~308), its corner, and positive selection (Theorem~\ref{t:select}), find a rational number $q$ with
\[
0<\ell-q<10^{-6},
\]
and prove both inequalities.

\smallskip
You may cite the handbook and \S\S15--18 of this notebook. When you cite a statement, check its hypotheses in the case at hand. Exact computation (integers, rationals, polynomials, linear algebra over $\mathbb Q$) may be done by machine; report it so that it can be repeated. Floating-point arithmetic, numerical integration, integer-relation searches and comparison with a known decimal value of $\ell$ may be used as checks. They earn no credit.
\smallskip}}
\end{center}

\subsection*{19.2\quad The answer and the scheme}

\emph{Answer.} $q=a_9=\dfrac{2447622221}{1044248800}=2.3439071426\ldots$, with
\[
\frac{2447622221}{1044248800}<\ell<\frac{6414407843}{2736630400},\qquad b_9-a_9=\frac{6561}{59511099776}<1.11\cdot10^{-7}.
\]
Any rational $q$ with a complete proof of $0<\ell-q<10^{-6}$ earns full credit, whatever bounds the proof uses (other consequences of positivity, such as the determinants of Thm.~35.10(iii), are admissible), and the points of the scheme then go to the corresponding steps of that proof; \S19.4, items 1, 5 and 6, lists the alternatives a grader will meet most often. What the key proves is this.

\begin{theorem}[The three-mark centre, bracketed]\label{t:bracket}\derived; the table \computed\ exactly.\
Let $c$ and $r$ be the solutions of
\[
k^2x(k)-(10k^2+10k+3)\,x(k+1)+9(k+1)^2\,x(k+2)=0\qquad(k\ge1)
\]
with $c(1)=1$, $c(2)=\frac13$ and $r(1)=0$, $r(2)=-\frac23$. For every integer $K\ge2$, $c(K)>0$ and $9c(K)>c(K-1)$, and
\begin{gather*}
a_K=-\frac{r(K)}{c(K)}\ <\ \Pth\ <\ b_K=\frac{r(K-1)-9\,r(K)}{9\,c(K)-c(K-1)},\\
b_K-a_K=\frac{2\cdot9^{2-K}}{3(K-1)^2\,c(K)\big(9c(K)-c(K-1)\big)} .
\end{gather*}
\end{theorem}

\begin{center}
\small
\begin{tabular}{@{}r p{0.72\linewidth} r@{}}
\toprule
Step & What is shown & Points\\ \midrule
1 & The ray: the chart and its measure, the stride, the values, the walls; $\ell=P(1)$ & 3\\
2 & The stride: $f\le\frac19$; $\frac19e^{-\varphi}\le f\le e^{-\varphi}$; $\int e^{-k\varphi}dy=6/k^2$; what follows at every real $k>0$ & 5\\
3 & The closings, their recurrence, the reflected candidate law & 3\\
4 & A certificate for Thm.~35.3 and its exact verification & 5\\
5 & The law at every real $k>0$: the integration by parts, with the bounds of Step 2 & 5\\
6 & The native values; what positivity can select & 2\\
7 & The corner: $k^2P(k)\to6$, so $9P(2)=3\ell-6$ & 5\\
8 & The rational solutions, their signs, the exact table & 4\\
9 & The bracket $a_K<\ell<b_K$, the choice of $K$, the answer & 8\\ \midrule
& & 40\\
\bottomrule
\end{tabular}
\end{center}

\subsection*{19.3\quad The key}

Throughout, $e_1,e_2,e_3$ are the elementary symmetric polynomials of $\alpha_1,\alpha_2,\alpha_3$, so that $K=e_1$ and $\Psi=e_2$, and $D_e=\alpha_e\,\partial/\partial\alpha_e$.

\paragraph{Step 1 (3 points). The ray.}
Take the chart $\alpha_3=1$ of the mark $e_0=3$, with $\alpha_1=e^{y_1}$ and $\alpha_2=e^{y_2}$, so that
\[
dy=dy_1\,dy_2=\frac{d\alpha_1\,d\alpha_2}{\alpha_1\alpha_2}.
\]
The centre ray at the strength $\sigma=1$ (Def.~35.9, p.~309) has origin $0$ and direction $(\one;1,1)$. Its stride and its values are
\[
f=\frac{\alpha^{\one}}{K\Psi}=\frac{e_3}{e_1e_2},\qquad P(k)=P_F(k\one;k,k)=\int_{\mathbb R^2}f^k\,dy\qquad(k=1,2,\dots),
\]
so that
\[
\ell=P(1)=\int_0^\infty\!\!\int_0^\infty\frac{d\alpha_1\,d\alpha_2}{(1+\alpha_1+\alpha_2)(\alpha_1+\alpha_2+\alpha_1\alpha_2)} .
\]
For real $k>0$ the same integral defines $P(k)$; Step 2 shows that it converges. The walls of the direction are $\kappa_{1,S}=|S|-d(S)-d^\ast(S)$ (Def.~8.1, p.~53), the width $c(S)$ of Def.~7.10, p.~49. Every nonempty $S$ has rank $r(S)=1$, so for $S$ nonempty and proper $d(S)=|S|-r^\ast(S)=r-r(\cE\setminus S)=0$ and $d^\ast(S)=|S|-r(S)=|S|-1$, and
\[
\kappa_{1,S}=1\qquad(\varnothing\neq S\subsetneq\cE).
\]
All proper walls are positive, the hypothesis of Thm.~35.3 and of Thm.~35.10, p.~309.

\smallskip\noindent\emph{Marking.} 1 point for the stride with the chart and its measure; 1 for $\ell=P(1)$ with the definition of the values cited; 1 for the walls. The measure $d\alpha_1d\alpha_2$ in place of $dy$ makes the integral diverge (the integrand is $\sim1/(\alpha_1+\alpha_2)$ at infinity): $-1$.

\paragraph{Step 2 (5 points). The stride.}
(a) \emph{The gauge.} The identity
\[
e_1e_2-9e_3=\alpha_1(\alpha_2-\alpha_3)^2+\alpha_2(\alpha_3-\alpha_1)^2+\alpha_3(\alpha_1-\alpha_2)^2
\]
holds because both sides equal $\sum_{i\neq j}\alpha_i^2\alpha_j-6e_3$, the left since $e_1e_2=\sum_{i\neq j}\alpha_i^2\alpha_j+3e_3$. So $0<f\le\frac19$ on the orthant, with equality exactly on the diagonal $\alpha_1=\alpha_2=\alpha_3$, which on the chart is the single point $y=0$. This is Thm.~35.10(i) here, with $\alpha^\star=\one$ and $f(\alpha^\star)=\frac19$.

(b) \emph{The decay.} Order the coordinates, $\alpha_{(1)}\ge\alpha_{(2)}\ge\alpha_{(3)}$. Then $\alpha_{(1)}\le e_1\le3\alpha_{(1)}$, $\alpha_{(1)}\alpha_{(2)}\le e_2\le3\alpha_{(1)}\alpha_{(2)}$ and $e_3=\alpha_{(1)}\alpha_{(2)}\alpha_{(3)}$, so
\[
\frac19\,\frac{\alpha_{(3)}}{\alpha_{(1)}}\le f\le\frac{\alpha_{(3)}}{\alpha_{(1)}} .
\]
On the chart $\alpha_{(3)}/\alpha_{(1)}=e^{-\varphi(y)}$ with
\[
\varphi(y)=\max(y_1,y_2,0)-\min(y_1,y_2,0),
\]
and since $0$ and $y_i$ are both among the three entries, $\varphi(y)\ge|y_i|$: $\varphi(y)\ge\|y\|_\infty$.

(c) \emph{The volume.} For real $k>0$,
\[
\int_{\mathbb R^2}e^{-k\varphi(y)}\,dy=\frac6{k^2}.
\]
The lines $y_1=y_2$, $y_1=0$, $y_2=0$ cut the plane into six cones, one for each order of the entries $(y_1,y_2,0)$, and on each cone $\varphi$ is the largest entry less the smallest:
\begin{gather*}
\int_{y_1\ge y_2\ge0}e^{-ky_1}dy=\int_0^\infty y_1e^{-ky_1}dy_1=\frac1{k^2},\qquad
\int_{y_1\ge0\ge y_2}e^{-k(y_1-y_2)}dy=\frac1k\cdot\frac1k,\\
\int_{0\ge y_1\ge y_2}e^{ky_2}dy=\int_0^\infty ue^{-ku}du=\frac1{k^2},
\end{gather*}
and the three cones with $y_1$ and $y_2$ exchanged give the same.

(d) \emph{Consequences}, for real $k>0$:
\[
\frac{6\cdot9^{-k}}{k^2}\le P(k)\le\frac6{k^2},\qquad P(k+1)<\frac19\,P(k),
\]
the second from (a), strict because $f<\frac19$ off the single point $y=0$. And $P$ is continuous on $(0,\infty)$: on $[k_1,\infty)$ the integrands $f^k$ are dominated by $f^{k_1}$, because $0<f<1$, and $f^{k_1}$ is integrable.

\smallskip\noindent\emph{Marking.} (a) 2 points; the identity must be shown or verified by expansion. The two inequalities of the arithmetic and geometric means, $e_1\ge3e_3^{1/3}$ and $e_2\ge3e_3^{2/3}$, are an accepted proof of $f\le\frac19$. The strict inequalities of Step 9 need only $f<\frac19$ on a set of positive measure; without any strictness the bracket is closed at the top, which suffices for the answer. (b) 1 point for the upper side, the decay; citing Thm.~35.10(i) for $f\le C'e^{-c'\|y\|}$ on the chart earns it, and gives convergence at every real $k>0$. 1 point for the lower side of (b) and the volume (c) together; they serve only the corner, and a solution that obtains the corner by a correct citation of Proposition~\ref{p:corner} or~\ref{p:dirclose} in Step 7 earns this point there. (d) 1 point: convergence, positivity, continuity and $P(k+1)<\frac19P(k)$. Step 2(d) credits the proof of the last inequality and Step 9(b) its use.

\paragraph{Step 3 (3 points). The candidate law.}
Here $1/f=e_1e_2/e_3=(\alpha_1+\alpha_2+\alpha_3)(\alpha_1^{-1}+\alpha_2^{-1}+\alpha_3^{-1})$, and the closings of the direction (Proposition~\ref{p:dirclose}) are
\[
x_j=[\alpha^0]\Big(\frac1f\Big)^j=\sum_{a_1+a_2+a_3=j}\Big(\frac{j!}{a_1!\,a_2!\,a_3!}\Big)^2=1,\ 3,\ 15,\ 93,\ 639,\ 4653,\ 35169,\ 272835,\ 2157759,\ \dots,
\]
since the constant term pairs the monomial $\alpha^a$ of $(\sum_e\alpha_e)^j$ with the monomial $\alpha^{-a}$ of $(\sum_e\alpha_e^{-1})^j$, and both have the coefficient $j!/(a_1!a_2!a_3!)$. Exact linear algebra with the ansatz $\sum_{i=0}^2p_i(j)\,x_{j+i}=0$, $\deg p_i\le2$, on $x_0,\dots,x_{15}$ (14 equations, 9 unknowns) has a one-dimensional space of solutions,
\[
(j+2)^2x_{j+2}-(10j^2+30j+23)\,x_{j+1}+9(j+1)^2x_j=0 .
\]
Reflected as in Corollary~\ref{c:dirlaw}, $c_i(k)=p_{2-i}(-k-2)$:
\begin{equation}\tag{L}
k^2P(k)-(10k^2+10k+3)\,P(k+1)+9(k+1)^2\,P(k+2)=0 .
\end{equation}
At this point (L) is a candidate. Corollary~\ref{c:dirlaw} runs the other way, from a law of $P$ to the recurrence of the closings, and sixteen closings do not prove their own recurrence.

\smallskip\noindent\emph{Marking.} 1 point each for the closings, the recurrence (with the size of the system, or any other exact derivation) and the reflection. Any route to (L) earns the 3 points: the closings are how the law is found, not part of its proof. A solution that takes (L) as proved here keeps these 3 points and earns nothing in Steps 4 and 5.

\paragraph{Step 4 (5 points). The certificate.}
The origin is $0$, so the defects of the origin vanish, and the defect of the direction is (Lemma~35.2(ii), p.~308)
\[
c_e=D_e\log f=1-R_e-R^\ast_e,\qquad R_e=\frac{\alpha_e}{e_1},\qquad R^\ast_e=\frac{\alpha_e(e_1-\alpha_e)}{e_2},
\]
the occupation of $e$ under the basis law and under the cobasis law. Since $\sum_eR_e=1$ and $\sum_eR^\ast_e=2$, $\sum_ec_e=0$. A \emph{certificate} of (L) is a choice of insertions $g_1,g_2,g_3$ (Def.~35.1: regular functions on the letter complement; here rational functions of degree $0$ whose denominators are products of $e_1$ and $e_2$, with coefficients polynomial in $k$) such that
\begin{equation}\tag{C}
k^2-(10k^2+10k+3)\,f+9(k+1)^2f^2=\sum_{e=1}^3\big(D_eg_e+k\,c_e\,g_e\big).
\end{equation}
\emph{Ansatz:} $g_e=N_e/(e_1^2e_2)$, with $N_e$ homogeneous of degree $4$ with coefficients in $\mathbb Q+\mathbb Qk$, $g_e$ the image of $g_1$ under the exchange of the marks $1$ and $e$, and $N_1$ symmetric in $\alpha_2,\alpha_3$. Multiplied by $e_1^3e_2^2$, (C) becomes a polynomial identity, of degree $7$ in $\alpha$ and $2$ in $k$, through
\[
e_1^3e_2^2\big(D_eg_e+k\,c_eg_e\big)=e_1e_2\,D_eN_e-N_e\big(2\alpha_ee_2+\alpha_e(e_1-\alpha_e)e_1\big)+k\,N_e\big(e_1e_2-\alpha_ee_2-\alpha_e(e_1-\alpha_e)e_1\big),
\]
and it is linear in the $18$ coefficients of $N_1$: $99$ equations of rank $10$. One solution is
\begin{equation}\tag{G}
g_1=-\frac{\alpha_1^2\big(k\,\alpha_1(\alpha_2+\alpha_3)+(2+3k)(\alpha_2^2+\alpha_3^2)-(1-k)\,\alpha_2\alpha_3\big)}{e_1^2e_2},
\end{equation}
and within the ansatz the others differ from it by one insertion $h$ added at every mark, $h$ with a symmetric numerator (an $8$-dimensional space). Such an $h$ changes nothing, because $\sum_eD_eh=0$ (Euler's identity in degree $0$) and $\sum_ec_eh=0$.

\emph{Verification.} Write $g_e=g^{(0)}_e+k\,g^{(1)}_e$. The three powers of $k$ in (C) read
\begin{align*}
k^2:&\quad \textstyle\sum_ec_eg^{(1)}_e=1-10f+9f^2=(1-f)(1-9f),\\
k^1:&\quad \textstyle\sum_e\big(D_eg^{(1)}_e+c_eg^{(0)}_e\big)=-10f+18f^2,\\
k^0:&\quad \textstyle\sum_eD_eg^{(0)}_e=-3f+9f^2 .
\end{align*}
Times $e_1^3e_2^2$ each sum over the marks is a symmetric polynomial of degree $7$, and in the basis of products of $e_1,e_2,e_3$ exact expansion gives
\begin{gather*}
\textstyle\sum_ee_1^3e_2^2\,c_eg^{(1)}_e=e_1^3e_2^2-10\,e_1^2e_2e_3+9\,e_1e_3^2,\\
\textstyle\sum_ee_1^3e_2^2\,D_eg^{(1)}_e=-2\,e_1e_2^3-3\,e_1^2e_2e_3+9\,e_1e_3^2,\qquad
\sum_ee_1^3e_2^2\,c_eg^{(0)}_e=2\,e_1e_2^3-7\,e_1^2e_2e_3+9\,e_1e_3^2,\\
\textstyle\sum_ee_1^3e_2^2\,D_eg^{(0)}_e=-3\,e_1^2e_2e_3+9\,e_1e_3^2,
\end{gather*}
which, with $f=e_3/(e_1e_2)$, are the three right sides times $e_1^3e_2^2$. So (C) holds identically in $\alpha$ and $k$. Proposition~\ref{p:certs} reported a certificate of this law with the same denominator.

\smallskip\noindent\emph{Marking.} 1 point for the defects $c_e$, in any equivalent form. 2 points for a certificate. Any certificate is accepted; it need not be symmetric, and any denominator is accepted for which Step 5 can be carried out. 2 points for the verification: by exact expansion, by the symmetric-function form above, or by exact evaluation of the cleared identity on a grid large enough to decide it. After clearing, both sides are polynomials of degree at most $7$ in each $\alpha_i$ and $2$ in $k$, homogeneous in $\alpha$, so agreement at $\alpha_3=1$ on $8$ values of $\alpha_1$, $8$ of $\alpha_2$ and $3$ of $k$ (all $192$ points) proves it. Agreement at a few random points is a test, not a proof. To test a certificate other than (G), a grader evaluates both sides of (C) at $\alpha=(1,2,4)$, $k=\frac13$, where the left side is $-\frac{2809}{7203}$ (there $f=\frac4{49}$, and (G) gives $g_1,g_2,g_3=-\frac{85}{1029},-\frac{310}{1029},-\frac{424}{1029}$); a certificate written on the chart, with $g_3=0$ and $D_1,D_2$ read as $\partial_{y_1},\partial_{y_2}$, is tested at the same point $(\alpha_1,\alpha_2)=(\frac14,\frac12)$. A certificate with floating-point coefficients, or verified only in floating point, earns 0 of the 4: the identity has to hold exactly. How the certificate was found does not matter: rational coefficients found by a numerical search and then verified exactly earn full credit.

\paragraph{Step 5 (5 points). The law at every real $k>0$.}
Thm.~35.3 states the relations for all \emph{large} $k$. The law is needed at $k=1,\dots,K-2$ (Step 8) and as $k\to0^+$ (Step 7), so the integration by parts is carried out here at every real $k>0$, with the bounds of Step 2. Fix a real $k>0$.

(a) $D_e(f^k)=k\,c_ef^k$, so the right side of (C), times $f^k$, is $\sum_eD_e(g_ef^k)$.

(b) On the chart $D_1=\partial_{y_1}$ and $D_2=\partial_{y_2}$, and for every function $h$ of degree $0$, $D_3h=-D_1h-D_2h$ (Euler's identity). So
\[
\sum_{e=1}^3D_e(g_ef^k)=\partial_{y_1}H_1+\partial_{y_2}H_2,\qquad H_i=(g_i-g_3)f^k .
\]

(c) \emph{Bounds.} Every monomial of every $N_e$ occurs in the expansion of $e_1^2e_2$: for $N_1$, $\alpha_1^3\alpha_2=\alpha_1^2\cdot\alpha_1\alpha_2$, $\alpha_1^2\alpha_2^2=\alpha_1\alpha_2\cdot\alpha_1\alpha_2$ and $\alpha_1^2\alpha_2\alpha_3=\alpha_1^2\cdot\alpha_2\alpha_3$, with the images under $2\leftrightarrow3$, and the other $N_e$ are images of $N_1$. The coefficients of $e_1^2e_2$ are positive, so $\alpha^v\le e_1^2e_2$ on the orthant for each such monomial, and $|g_e|\le\|N_e\|_1$, the sum of the absolute values of the coefficients. For every mark $i$,
\[
D_ig_e=\frac{D_iN_e}{e_1^2e_2}-g_e\,\big(2R_i+R^\ast_i\big),
\]
where $D_iN_e$ has the monomials of $N_e$ or fewer and $0\le R_i,R^\ast_i\le1$; and $|c_i|\le1$. With Step 2(b) there is, for fixed $k$, a constant $C$ with
\[
|H_i|\le C\,e^{-k\varphi},\qquad |\partial_{y_j}H_i|\le C\,e^{-k\varphi} .
\]

(d) If $H$ is $C^1$ on $\mathbb R^2$, tends to $0$ at infinity and $\partial_{y_1}H$ is integrable, then $\int\partial_{y_1}H\,dy=0$: for every $y_2$, $\int_{-R}^R\partial_{y_1}H\,dy_1=H(R,y_2)-H(-R,y_2)\to0$, and Fubini applies. By (c) and Step 2, $H_1$ and $H_2$ qualify ($e^{-k\varphi}\le e^{-k\|y\|_\infty}\to0$, and $e^{-k\varphi}$ is integrable), so the right side of (C), times $f^k$, integrates to $0$.

(e) The left side of (C), times $f^k$, integrates to the left side of (L), all three integrals convergent by Step 2. Hence (L) holds at every real $k>0$.

\smallskip\noindent\emph{Marking.} 1 point for (a)--(b), including the mark $3$ on its own chart. 2 points for the bounds (c). 1 point for (d). 1 point for the conclusion at every real $k>0$, with the reason it is needed. Citing Thm.~35.3 at $k=1,2,\dots$ and as $k\to0^+$ without (c): 2 of 5. Any complete extension earns full credit; for instance (C) integrated at large integers by Thm.~35.3 as stated (its hypothesis checked in Step 1), then carried to every $k$ with $\re k>0$ by Carlson's theorem: the left side of (L) is holomorphic there, it is bounded by a polynomial in $|k|$ on $\re k\ge1$ because $|P(k)|\le P(\re k)\le6/(\re k)^2$ (Step 2), and Carlson's theorem applied to $z\mapsto F(N+z)$, with $F$ the left side of (L) and $N\ge1$ an integer from which on (L) holds at every integer, makes $F$ vanish identically. Integrating $D_3(g_3f^k)$ by parts ``in $y_3$'' on the chart $\alpha_3=1$, where $y_3$ is not a coordinate: $-1$, unless repaired by Euler's identity or by a certificate with $g_3=0$.

\paragraph{Step 6 (2 points). The native values; what positivity can select.}
The symbol of (L) (Cor.~35.6, p.~308) is
\[
\chi(\lambda)=1-10\lambda+9\lambda^2=(1-\lambda)(1-9\lambda).
\]
The native gauges of the centre ray are the signs $\eta$ with $\eta_3=1$ (Proposition~\ref{p:onlysigns}, Thm.~35.12, p.~310), of native value $K(\eta)^{-2}$: the table's gauge $\eta=(1,1,1)=\alpha^\star$, with $K(\eta)=3$ and the value $\frac19=f(\alpha^\star)$; and $\eta=(-1,1,1),(1,-1,1),(-1,-1,1)$, with $K(\eta)=\pm1$ and the value $1$. They are nondegenerate. By Thm.~35.12(ii) the Hessian of $-\log f$ at $\eta$ is $2\big(\kappa_\eta(X_e,X_{e'})\big)_{e,e'\neq3}$, the covariances of the indicators under the signed measure $\mathbb P_\eta(\{e\})=\eta_e/K(\eta)$. Its determinant is $\frac4{27}$ at $(1,1,1)$ and $-4$ at each of the other three signs. So Cor.~35.6, with $g=1$, makes both native values roots of $\chi$, and since $\chi$ has degree $2$ they are all its roots. (Cor.~35.6 runs from native values to roots; a root need not be a native value, as the extra root $\frac{131003}{4771109}$ in \S18.4 shows.)

In Theorem~\ref{t:select} the singular points are $q=9$, where the slit $[9,\infty)$ begins, and $q=1$, off the slit. At $q=1$, a simple root and the value of nondegenerate gauges, positivity imposes $\nu=1$ condition, and with the corner $1+\nu=2=h$. The condition is independent, and here that is seen directly: by Step 8, $c(K)\ge3^{1-K}$ and $9c(K)-c(K-1)\ge2c(K-1)\ge2\cdot3^{2-K}$, so the width of Step 9(d) is at most $1/(K-1)^2$ and the bracket closes on one number. So the law, the corner and positivity determine $P$ (Theorem~\ref{t:select}(iii)). The rational solutions of (L) have the exponential rate $1$ of the signs and $P$ the rate $\frac19$ of the gauge, so the bracket shrinks by about $\frac19$ per step, $\log_{10}9=0.954$ digits: about seven steps for $10^{-6}$, before constants. The constants decide in Step 9.

\smallskip\noindent\emph{Marking.} 1 point for the native values with their nondegeneracy, matched to the roots of the symbol in the direction Cor.~35.6 gives. ``The roots are native values by Cor.~35.6'' reverses the corollary and earns 0 for this point unless the roots are matched to gauges directly, as above. 1 point for the count $1+\nu=h$ with the independence, or any equivalent statement that the corner supplies the condition positivity does not. The step is not used in the proofs of the bounds; a complete solution without it loses these 2 points only.

\paragraph{Step 7 (5 points). The corner.}
(a) By Step 2(d), $6\cdot9^{-k}\le k^2P(k)\le6$ for real $k>0$, so
\[
\lim_{k\to0^+}k^2P(k)=6 .
\]
(b) Take (L) at real $k\in(0,1)$ (Step 5) and let $k\to0^+$. By the continuity of Step 2(d), $P(k+1)\to P(1)=\ell$ and $P(k+2)\to P(2)$, while $10k^2+10k+3\to3$ and $9(k+1)^2\to9$. So $6-3\ell+9P(2)=0$:
\begin{equation}\tag{S}
P(2)=\frac{\ell-2}3 .
\end{equation}
(c) This is Corollary~\ref{c:dirlaw} at $j=0$,
\[
[k^2]c_0\cdot\mathrm{Cor}+c_1(0)P(1)+c_2(0)P(2)=0,\qquad \mathrm{Cor}=\sum_\pi\prod_{i=2}^3\kappa_{1,U_i}^{-1}=3!\cdot1=6
\]
(Proposition~\ref{p:dirclose}; for $\binom{\cE}1$ the corner is $n!$, Proposition~\ref{p:corner}). The six cones of Step 2(c) are the six orders $\pi$, each contributing $\prod_i\kappa_{1,U_i}^{-1}=1$.

\smallskip\noindent\emph{Marking.} 2 points for (a), 2 for (b), 1 for (c). A correct citation of Proposition~\ref{p:corner} or~\ref{p:dirclose} in place of (a), with the walls of Step 1 and the sum over the $3!$ orders, earns the 2 points of (a). Errors:
\begin{itemize}
\item Reading (L) at $k=0$ as $0\cdot P(0)-3P(1)+9P(2)=0$. Here $P(0)=\int dy=\infty$, and the product $0\cdot\infty$ is the corner, $6$, not $0$. The reading gives $P(2)=\ell/3$, which is false ($P(2)=0.1146\ldots$, $\ell/3=0.78\ldots$) and contradicts Step 2(d) at once: $P(2)<\frac19P(1)$ would say $\ell/3<\ell/9$. After it every $P(K)$ is $\ell$ times a positive rational, positivity fixes only ratios, and the bracket is empty. 0 for Step 7, and in Steps 8--9 only the points that do not depend on (S): the signs and their use, 2 points.
\item Passing the limit under the integral, $k^2f^k\to0$ pointwise, so ``$k^2P(k)\to0$'': 0 for (a). There is no domination, since $\int dy=\infty$.
\item The wrong power ($kP(k)$, the corner of a ray on two marks) or $\mathrm{Cor}=1$ (one order): 0 for (a); (b) is marked as a consequence.
\item Citing Corollary~\ref{c:dirlaw} for (S). Its hypothesis is the law \emph{identically}, as meromorphic functions of $k$; Step 5 gives it at real $k>0$, and the continuation of Proposition~\ref{p:dirclose} extends it. Citing the corollary with the law proved only at the integers leaves a gap: $-1$.
\end{itemize}

\paragraph{Step 8 (4 points). The rational solutions and their signs.}
For integers $k\ge1$, $9(k+1)^2\neq0$, and (L) gives
\[
P(k+2)=\frac{(10k^2+10k+3)\,P(k+1)-k^2P(k)}{9(k+1)^2}.
\]
Let $c$ and $r$ be the solutions of (L) with $c(1)=1$, $c(2)=\frac13$ and $r(1)=0$, $r(2)=-\frac23$. From $P(1)=\ell$, (S) and induction,
\[
P(K)=c(K)\,\ell+r(K)\qquad(K\ge1).
\]
\emph{Signs, at every $K$.} Put $x'_j=9^jc(j+1)$. Then (L) at $k=j+1$ reads $(j+2)^2x'_{j+2}=(10j^2+30j+23)\,x'_{j+1}-9(j+1)^2x'_j$, with $x'_0=1$, $x'_1=3$. If $x'_{j+1}\ge3x'_j>0$, then
\[
(j+2)^2x'_{j+2}\ge\big(10j^2+30j+23-3(j+1)^2\big)x'_{j+1}=(7j^2+24j+20)\,x'_{j+1}\ge3(j+2)^2x'_{j+1},
\]
since $7j^2+24j+20-3(j+2)^2=4(j+1)(j+2)\ge0$. So $x'_{j+1}\ge3x'_j>0$ for every $j$, hence $x'_j\ge3^j$, that is,
\[
c(K)\ge3^{1-K}>0,\qquad 9c(K)\ge3c(K-1)>c(K-1)\qquad(K\ge2).
\]
The $x'_j$ obey the recurrence of the closings with the same first two terms, so $x'_j=x_j$ and $c(K)=9^{1-K}x_{K-1}$ wherever that recurrence holds: in the table below by inspection, and for every $j$ by Corollary~\ref{c:dirlaw} once (L) holds identically. The proof of the signs does not use this.

\begin{center}
\small
\renewcommand{\arraystretch}{1.12}
\begin{tabular}{@{}r r r r l@{}}
\toprule
$K$ & $c(K)$ & $r(K)$ & $a_K=-r(K)/c(K)$ & \\ \midrule
1 & $1$ & $0$ & $0$ & \\
2 & $1/3$ & $-2/3$ & $2$ & \\
3 & $5/27$ & $-23/54$ & $23/10$ & $2.3$\\
4 & $31/243$ & $-145/486$ & $145/62$ & $2.33870967\ldots$\\
5 & $71/729$ & $-1331/5832$ & $1331/568$ & $2.34330985\ldots$\\
6 & $517/6561$ & $-242353/1312200$ & $242353/103400$ & $2.34383945\ldots$\\
7 & $11723/177147$ & $-5495507/35429400$ & $5495507/2344600$ & $2.34389959\ldots$\\
8 & $10105/177147$ & $-46422941/347208120$ & $46422941/19805800$ & $2.34390638\ldots$\\
9 & $26639/531441$ & $-2447622221/20832487200$ & $2447622221/1044248800$ & $2.34390714\ldots$\\
10 & $5773279/129140163$ & $-530455590781/5062294389600$ & $530455590781/226312536800$ & $2.34390722\ldots$\\
\bottomrule
\end{tabular}
\end{center}
Two checks a grader can run on a student's table: $9^{K-1}c(K)$ is the closing $x_{K-1}$ ($1,3,15,93,\dots$), and the frame of the two solutions (Prop.~35.8, p.~309: $C_{k+1}=\frac{c_0(k)}{c_2(k)}C_k=\frac{k^2}{9(k+1)^2}C_k$, $C_1=-\frac23$) is
\[
c(k)\,r(k+1)-c(k+1)\,r(k)=-\frac23\cdot\frac{9^{1-k}}{k^2} .
\]

\smallskip\noindent\emph{Marking.} 2 points for $P(K)=c(K)\ell+r(K)$ with the initial values from (S). 1 point for the signs $c(K)>0$ and $9c(K)-c(K-1)>0$, at every $K$ as above or at the $K$ used from the exact table. 1 point for the table, exact, through the $K$ used. A decimal table: 0 for the last point.

\paragraph{Step 9 (8 points). The bracket and the answer.}
(a) \emph{Below, by positivity.} $P(K)=c(K)\ell+r(K)>0$, divided by $c(K)$, gives
\[
\ell>a_K=-\frac{r(K)}{c(K)} .
\]
(b) \emph{Above, by the gauge.} By Step 2(d), $P(K)<\frac19P(K-1)$, that is, $c(K)\ell+r(K)<\frac19\big(c(K-1)\ell+r(K-1)\big)$, or $\big(9c(K)-c(K-1)\big)\ell<r(K-1)-9r(K)$, and divided by $9c(K)-c(K-1)$,
\[
\ell<b_K=\frac{r(K-1)-9\,r(K)}{9\,c(K)-c(K-1)} .
\]
(c) \emph{The signs.} Both divisions keep the direction of the inequality: $c(K)>0$ and $9c(K)-c(K-1)>0$ by Step 8.

At $K=2$ this is the first bracket: $P(2)=\frac{\ell-2}3>0$ gives $\ell>2$, and $P(2)<\frac19P(1)$ gives $\ell<3$.

(d) \emph{The width and the answer.} With the frame of Step 8,
\[
b_K-a_K=\frac{r(K-1)\,c(K)-r(K)\,c(K-1)}{c(K)\big(9c(K)-c(K-1)\big)}=\frac{2\cdot9^{2-K}}{3(K-1)^2\,c(K)\big(9c(K)-c(K-1)\big)},
\]
and in exact arithmetic:
\begin{center}
\small
\renewcommand{\arraystretch}{1.12}
\begin{tabular}{@{}r r r r@{}}
\toprule
$K$ & $b_K$ & $b_K-a_K$ & width with $P(K)<9^{1-K}\ell$ instead\\ \midrule
2 & $3$ & $1$ & $1$\\
3 & $19/8$ & $7.50\cdot10^{-2}$ & $1.64\cdot10^{-1}$\\
4 & $61/26$ & $7.44\cdot10^{-3}$ & $2.54\cdot10^{-2}$\\
5 & $3413/1456$ & $7.84\cdot10^{-4}$ & $3.67\cdot10^{-3}$\\
6 & $104539/44600$ & $8.43\cdot10^{-5}$ & $5.04\cdot10^{-4}$\\
7 & $149014/63575$ & $9.17\cdot10^{-6}$ & $6.66\cdot10^{-5}$\\
8 & $909876251/388187800$ & $1.0036\cdot10^{-6}$ & $8.59\cdot10^{-6}$\\
9 & $6414407843/2736630400$ & $1.1025\cdot10^{-7}$ & $1.0863\cdot10^{-6}$\\
10 & $232184895407/99058909600$ & $1.21\cdot10^{-8}$ & $1.35\cdot10^{-7}$\\
\bottomrule
\end{tabular}
\end{center}
At $K=8$ the certified width is $1.0036\cdot10^{-6}>10^{-6}$; at $K=9$ it is $\frac{6561}{59511099776}<1.11\cdot10^{-7}$. So
\[
\frac{2447622221}{1044248800}=a_9<\ell<b_9=\frac{6414407843}{2736630400},
\]
and $q=a_9$ satisfies $0<\ell-q<b_9-a_9<1.11\cdot10^{-7}<10^{-6}$.

\smallskip\noindent\emph{Marking.} (a) 2 points. (b) 2 points: 1 for using $P(K)<\frac19P(K-1)$ (its proof is credited in Step 2(d)), 1 for $b_K$. (c) 1 point for invoking the two signs at the divisions (their proof is credited in Step 8; a sign proved there and used here silently loses nothing). (d) 3 points: 1 for a $K$ at which the width certified by (a)--(c) is below $10^{-6}$, 1 for the exact comparison, 1 for a stated $q$ whose inequalities $0<\ell-q<10^{-6}$ follow from the bracket.

\subsection*{19.4\quad Where a solution that looks right is wrong}

\begin{enumerate}[label=\arabic*.]
\item \emph{The answer one step early.} $a_8=\frac{46422941}{19805800}$ is a correct answer: $\ell-a_8=8.58\cdot10^{-7}$. But the bracket at $K=8$ has width $1.0036\cdot10^{-6}$, so a solution that certifies $a_8$ from $K=8$ alone has an error in it. The upper bound at $K=9$ certifies it: $\ell-a_8<b_9-a_8=8.72\cdot10^{-7}$. Full credit with that; with the bound claimed at $K=8$, 5 of the 8 points of Step 9. The true error falls below the tolerance one step before the certified one does.
\item \emph{The observed error.} Comparing $a_K$ with a decimal value of $\ell$ (from the share side, from a quadrature, or from a closed form recalled elsewhere) proves nothing about $\ell-a_K$: 0 for Step 9(d).
\item \emph{The rate as a bound.} Theorem~\ref{t:select}(iv) gives the error $O(9^{-K})$ up to powers of $K$, with no constant: 0 for Step 9(d). Nor is the last step a bound: $a_9-a_8=7.62\cdot10^{-7}$ is smaller than $\ell-a_8=8.58\cdot10^{-7}$, because the $a_K$ increase to $\ell$ with ratio about $\frac19$, so the remaining error is about $\frac98$ of the next step.
\item \emph{An unproved upper bound.} ``$P(K)\le9^{-K}\ell$, since $f\le\frac19$''. The gauge gives $P(K)<\frac19P(K-1)$, and iterated $P(K)<9^{1-K}P(1)$; the factor $9^{-K}$ does not follow from $f\le\frac19$, it would need $\int dy<\infty$. (The inequality is in fact true for $K\ge8$ and false for $K\le7$: $9^KP(K)$ decreases by Step 2(d), and $9^7P(7)/\ell=1.03$, $9^8P(8)/\ell=0.90$. Nothing in such a solution proves it, and a grader who checks it numerically should not be misled.) With it the width at $K=9$ is $1.21\cdot10^{-7}$, while the simple bound actually proved gives $1.09\cdot10^{-6}$ there, so the solution certifies nothing below $10^{-6}$. The stated $q=a_9$ is right and the proof is not: 1 of 2 in Step 9(b), 0 in 9(d), 4 of the 8 points of Step 9.
\item \emph{The simple bound, done right.} $P(K)<9^{1-K}P(1)=9^{1-K}\ell$ gives $\big(c(K)-9^{1-K}\big)\ell<-r(K)$, and $c(K)-9^{1-K}=9^{1-K}(x'_{K-1}-1)>0$ for $K\ge2$ by Step 8, so $\ell<-r(K)/\big(c(K)-9^{1-K}\big)$: valid and cruder (the last column of the table). It needs $K=10$, with $q=a_{10}=\frac{530455590781}{226312536800}$. Full credit, with the sign stated.
\item \emph{Decimals.} $q=2.343907$ is accepted, with the exact comparisons $\frac{2343907}{10^6}<a_9$ and $b_9-\frac{2343907}{10^6}=2.53\cdot10^{-7}$. A decimal above $a_9$, such as $2.3439072$, is a correct answer ($\ell-2.3439072=3.9\cdot10^{-8}$) but needs a new lower bound: $a_{10}=2.34390722\ldots$ supplies it. A decimal rounded above the lower bound in hand: $-1$ in Step 9(d), its last point.
\item \emph{The signs.} Signs of $c(K)$ and of $9c(K)-c(K-1)$ neither proved nor checked: $-2$ in all, the sign point of Step 8 and Step 9(c). They hold here at every $K$ (Step 8), but where a rational solution changes sign the inequality reverses; on the line $1$ of $\binom{[5]}2$, Remark~\ref{r:paper} checked the sign only for $m\le80$.
\item \emph{The corner.} Reading the law at $k=0$ with $0\cdot P(0)=0$, or passing the limit $k\to0^+$ under the integral: the approximation is lost (Step 7). The first gives $P(2)=\ell/3$, every $P(K)$ is then $\ell$ times a positive rational, positivity gives only $\ell>0$, and the gauge bound at $K=2$, $\ell/3<\ell/9$, gives $\ell<0$: the student's own bracket is empty, which the grader can point to.
\item \emph{The theorem at small $k$.} Thm.~35.3 cited at $k=1,2,\dots$ or as $k\to0^+$ without the decay of Step 2: Step 5.
\item \emph{A certificate checked at sample points only}: Step 4.
\item \emph{The measure} $d\alpha_1d\alpha_2$ on the chart: Step 1.
\item \emph{A quadrature for $\ell$ or for $P(2)$ used in the solution}: no credit at the step where it is used. It is a check (\S19.5).
\end{enumerate}

\subsection*{19.5\quad Receipts and the independent check}

The script exam\_s19.py prints every number of this entry. Exact: the sum of squares; the closings $x_0,\dots,x_{15}$ and their recurrence (14 equations, 9 unknowns, nullity 1); the reflection; the certificate (G) and the identity (C) after clearing $e_1^3e_2^2$; its three powers of $k$ in the basis of products of $e_1,e_2,e_3$; the rank $10$ of the symmetric ansatz in $18$ unknowns; the monomial check of Step 5(c); the spot value $-\frac{2809}{7203}$; the symbol and the signs; the walls and $\mathrm{Cor}=6$; the table and the brackets to $K=10$; and to $K=16$ the signs with $c(K)\ge3^{1-K}$, the identity $9^{K-1}c(K)=x_{K-1}$, the frame, and the two forms of the width with its bound $1/(K-1)^2$.

\noindent Check \computed, not part of the solution, in ball arithmetic, with the integrals cut where their tails are below $10^{-70}$. On the chart, the integral over $\alpha_2$ is elementary, $\int_0^\infty\frac{d\alpha_2}{(\alpha_2+A)(A\alpha_2+B)}=\frac{\log(A^2/B)}{A^2-B}$ with $A=1+\alpha_1$, $B=\alpha_1$, so
\[
\ell=\int_0^\infty\frac{\log\big((1+x)^2/x\big)}{1+x+x^2}\,dx=2.343907238689458890601562288872277069095\ldots .
\]
On the share side (Proposition~\ref{p:bananapin}), $\Gamma(k)^2P(k)=\int_0^\infty y^{k-1}\sG_{0;(1,-1)}(1,y)^3\,dy$ gives the same $\ell$ to $40$ digits, $P(2)-\frac{\ell-2}3=0$ to $47$ digits, and $P(9)=4.8152090557463543143\cdot10^{-9}=c(9)\ell+r(9)$. With this value,
\[
\ell-a_8=8.5773\cdot10^{-7},\qquad \ell-a_9=9.6062\cdot10^{-8},\qquad b_9-\ell=1.4186\cdot10^{-8}.
\]
The upper side is the sharper. Its slack is $\frac19P(K-1)-P(K)$, which is about $P(K)/(K-1)$ because $P(k)$ decays like $9^{-k}/k$ (for three marks the branch of Thm.~35.11, p.~310, is logarithmic), so $b_K-\ell$ is about $\frac9{8(K-1)}(\ell-a_K)$.

\subsection*{19.6\quad What the example shows}

The reading needed three inputs, and each was elementary once it was named. The analytic input is one inequality, $e_1e_2\ge9e_3$, sharpened to $f\asymp e^{-\varphi}$ by the tropical stride $\varphi$: it gives the convergence at every real $k>0$, the corner, and the upper side of the bracket. The algebraic input is one certificate: it gives the law at every real $k>0$, not only at large $k$. The arithmetic is exact. Positivity gives the lower side and the gauge the upper side. On every ray through the origin the first two inputs exist \derived: $f\le f(\alpha^\star)$ (Thm.~35.10(i)), and the tropical stride, whose volume is $\mathrm{Cor}$ (the six cones of Step 2(c) are the six orders of Proposition~\ref{p:dirclose}). A certificate exists for every law that follows from the relations (Thm.~35.19, p.~313), but Step 5 needs more than existence: bounds on the $g_e$ and their derivatives against the decay of $f^k$ down to $k\to0^+$. Here they came from the numerators lying among the monomials of $e_1^2e_2$, as the exponent condition did in Proposition~\ref{p:certs}; that a certificate with such bounds exists on every ray is \notyet. With $h-1$ masters, positivity of $P(K),\dots,P(K+h-2)$ and the gauge bounds $P(K+i)<f(\alpha^\star)P(K+i-1)$ cut out a polytope in the coordinates of the masters, and Theorem~\ref{t:select}(iv) is the rate at which it shrinks \notyet. The width of the bracket is a ratio of integers from the other end of the pencil (Remark~\ref{r:twochains}):
\[
b_K-a_K=\frac{2\cdot9^{K-1}}{3(K-1)^2\,x_{K-1}\,(x_{K-1}-x_{K-2})}\qquad\derived,
\]
from $c(K)=9^{1-K}x_{K-1}$, which holds at every $K$ by Corollary~\ref{c:dirlaw}: by Step 5 and the continuation of Proposition~\ref{p:dirclose}, (L) holds identically.

\noindent Open.
\begin{enumerate}[label=(\roman*)]
\item The same examination for the lean $(2\one;4,1)$ of \S18.3: its law has no certificate yet (\S18.6(i)), and the signs of its rational solutions are unproved. \notyet
\item The polytope for the two masters of $H$ (\S18.4), with a certificate for its law of order $4$. \notyet
\item The corner as the boundary term of the certificate at $k=0$, where $g^{(0)}_e$ does not vanish at infinity, instead of a limit. \notyet
\end{enumerate}

\subsection*{19.7\quad Corrections}

\begin{enumerate}[label=(\roman*)]
\item The receipts of Proposition~\ref{p:certs} say that Thm.~35.3 ``applies as stated'' to the banana rays. As stated it gives the relations for large $k$, and Remark~\ref{r:paper} used the three-mark law from $k=1$ and its corner from $k\to0^+$. Step 5 closes that gap for the three-mark law; the four-mark banana ray needs the same bounds, with its own tropical stride.
\item In Remark~\ref{r:paper} the ratios $-r/c$ are lower bounds at every $q$, since $c(q)>0$ at every $q$ (Step 8), not only at the $q$ computed.
\item Theorem~\ref{t:select} said that the roots of the symbol are native values, citing Cor.~35.6. The corollary gives the converse, that the native value of every nondegenerate native gauge is a root, and a root need not be a native value (the extra root $\frac{131003}{4771109}$ of \S18.4). The statement is corrected in place: (ii) speaks of roots, and the count of one condition in (iii) is for a simple root that is the value of a nondegenerate native gauge. Step 6 checks that hypothesis here.
\end{enumerate}

\noindent Aliases, last: $\ell=3L(\chi_{-3},2)$ \recalled, which agrees with the check to $40$ digits; the closings are the sums of squares of trinomial coefficients, the even moments of the planar walk of three unit steps; (L) is an Ap\'ery-like recurrence, and Step 9 is Ap\'ery's limit made two-sided by positivity (Pincherle's minimal solution, read with signs); the certificate is a creative-telescoping certificate for the integral; with $\sG_{0;(1,-1)}(1,y)=2K_0(2\sqrt y)$ the share side is the Bessel moment, $P(k)=4^{2-k}\Gamma(k)^{-2}\int_0^\infty x^{2k-1}K_0(x)^3dx$; $\varphi$ is the tropicalization of $-\log f$, and $\int e^{-k\varphi}dy=6/k^2$ its volume.

\end{document}
